% Matrices 2×2 et applications linéaires du plan — Mathématiques Tle TI — Cours signé OnBuch+
% Programme officiel MINESEC. Compiler avec : tectonic M20-matrices-2x2-et-applications-lineaires.tex
\newcommand{\DOCMATIERE}{Mathématiques}
\newcommand{\DOCNIVEAU}{Tle TI}
\newcommand{\DOCMODULE}{Configurations et transformations élémentaires du plan}
\newcommand{\DOCLECON}{20}
\newcommand{\DOCTITRE}{Matrices 2×2 et applications linéaires du plan}
% OnBuch+ — préambule « cours signés » ENRICHI (compilé avec tectonic / XeLaTeX)
% Système visuel « Pop » d'OnBuch+ : fond crème (jamais blanc), bordures encre
% épaisses, ombres « dures » décalées, titres Archivo Black en capitales.
% Version MATHÉMATIQUES : graphes (pgfplots/tikz), boîtes pédagogiques
% (définition, propriété, méthode, attention, le savais-tu, exercice résolu,
% exercice, TP, bilan), page de garde et signature OnBuch+ ancrée en bas.
%
% Macros attendues dans main.tex AVANT \input{preamble} :
%   \DOCMATIERE  \DOCNIVEAU  \DOCMODULE  \DOCLECON (numéro)  \DOCTITRE
\documentclass[11pt]{article}

\usepackage{fontspec}
\usepackage[french]{babel}
\usepackage[a4paper,margin=2cm,top=3.4cm,bottom=2.8cm,headheight=1.4cm,footskip=1.3cm]{geometry}
\usepackage{amsmath,amssymb,amsfonts}
\usepackage{mathtools} % \xmapsto, \coloneqq... souvent utilisés par les modèles
\usepackage{cancel} % simplifications barrées
% \abs{z} (module, valeur absolue) et \norm{u} (norme), utilisés par les modèles
\providecommand{\abs}{}\let\abs\relax\DeclarePairedDelimiter{\abs}{\lvert}{\rvert}
\providecommand{\norm}{}\let\norm\relax\DeclarePairedDelimiter{\norm}{\lVert}{\rVert}
\usepackage[table]{xcolor} % [table] : fournit aussi \rowcolors (alternance de lignes)
\usepackage{pagecolor}
\usepackage{tikz}
\usetikzlibrary{arrows.meta,positioning,calc,shapes.geometric,shapes.misc,shapes.arrows,decorations.pathmorphing,decorations.markings,patterns,shadows,mindmap,trees,fit,backgrounds,matrix,intersections}
\usepackage{pgfplots}
\usepackage{pgf-pie} % diagrammes circulaires \pie (statistiques)
\pgfplotsset{compat=1.18}
\usepgfplotslibrary{fillbetween,groupplots} % aires entre courbes (calcul intégral)
\usepackage{enumitem}
\usepackage{fancyhdr}
% [most] charge toutes les bibliothèques tcolorbox (listings, minted, raster,
% documentation...) et gonfle inutilement la mémoire TeX sur un document long
% (~300 boîtes) : on ne charge que ce qui est réellement utilisé.
\usepackage[skins,breakable]{tcolorbox}
\usepackage{graphicx}
\usepackage{colortbl}
\usepackage{booktabs}
\usepackage{tabularx}
\usepackage{diagbox} % case d'en-tête diagonale des tableaux à double entrée
% Colonnes extensibles centrée / gauche / droite (C, L, R), souvent utilisées par les modèles
\newcolumntype{C}{>{\centering\arraybackslash}X}
\newcolumntype{L}{>{\raggedright\arraybackslash}X}
\newcolumntype{R}{>{\raggedleft\arraybackslash}X}
\usepackage{array}
\usepackage{multicol}
\usepackage{siunitx}
% parse-numbers=false : accepte aussi \SI{2,5 \times 10^{-3}}{...}
\sisetup{locale=FR,output-decimal-marker={,},group-minimum-digits=4,parse-numbers=false}
\DeclareSIUnit\ohm{\ensuremath{\symup{\Omega}}} % U+2126 absent de la police
\usepackage{qrcode}
% \degree : commande fréquemment utilisée par les modèles pour un angle ou une
% température en dehors de \SI{}{} (non fournie par les paquets chargés).
\providecommand{\degree}{^{\circ}}
% \qed (fin de démonstration), utilisé par les modèles sans amsthm
\providecommand{\qed}{\hfill$\square$}
% \barre : double barre des valeurs interdites dans un tableau de variations (tkz-tab)
\providecommand{\barre}{\ensuremath{\|}}
% environnement proof (amsthm non chargé) utilisé par les modèles
\ifdefined\proof\else\newenvironment{proof}[1][Démonstration]{\par\noindent\textit{#1.}\ }{\qed\par}\fi
\usepackage{lastpage}
\usepackage{hyperref}
\hypersetup{hidelinks}

% ============================================================
% Palette « Pop » OnBuch+ — mêmes hex que la classe P (plus_ui.dart)
% ============================================================
\definecolor{poporange}{HTML}{F59321}
\definecolor{popdark}{HTML}{E07A0C}
\definecolor{popcream}{HTML}{FFF6E8}
\definecolor{popink}{HTML}{1C1714}
\definecolor{poppurple}{HTML}{7A5AE0}
\definecolor{popgreen}{HTML}{2E9E6A}
\definecolor{poppink}{HTML}{E0457B}
\definecolor{popblue}{HTML}{2F7DE1}
\definecolor{popgold}{HTML}{C98A12}
\definecolor{popmuted}{HTML}{8A7F76}
\definecolor{popline}{HTML}{EADFD2}
\definecolor{popgray}{HTML}{8A7F76}
\colorlet{popgrey}{popgray}
% Alias tolérants (le modèle nomme parfois les couleurs différemment)
\colorlet{popwhite}{white}
\colorlet{popblack}{popink}
\colorlet{popred}{poppink}
\colorlet{popyellow}{popgold}
% Teintes claires pour fonds de tableaux / zones
\colorlet{popblueL}{popblue!10!white}
\colorlet{poporangeL}{poporange!14!white}
\colorlet{popgreenL}{popgreen!12!white}
\colorlet{poppurpleL}{poppurple!12!white}
\colorlet{poppinkL}{poppink!12!white}
\colorlet{popgoldL}{popgold!14!white}

\pagecolor{popcream}

% ============================================================
% Typographie
% ============================================================
\defaultfontfeatures{Ligatures=TeX}
\setmainfont{PlusJakartaSans-Medium.ttf}[Path=../../../fonts/,BoldFont=PlusJakartaSans-Bold.ttf]
% Maths en sans-serif (Fira Math) avec chiffres et lettres droites de Plus
% Jakarta, pour que les formules se fondent dans le texte.
\usepackage[math-style=ISO,bold-style=ISO]{unicode-math}
\setmathfont{FiraMath-Regular.otf}[Path=../../../fonts/]
\setmathfont{PlusJakartaSans-Medium.ttf}[Path=../../../fonts/,range={up/num,up/{Latin,latin},\mathup}]
\usepackage{icomma} % virgule décimale sans espace en mode math
% Caractères mathématiques Unicode absents des polices de texte (Plus Jakarta,
% Archivo Black) : rendus par leur équivalent mathématique, en texte comme en
% formule — sinon ils s'affichent en carré vide (ex. « Arithmétique dans ℤ »).
\usepackage{newunicodechar}
\newunicodechar{ℝ}{\ensuremath{\mathbb{R}}}
\newunicodechar{ℤ}{\ensuremath{\mathbb{Z}}}
\newunicodechar{ℂ}{\ensuremath{\mathbb{C}}}
\newunicodechar{ℕ}{\ensuremath{\mathbb{N}}}
\newunicodechar{ℚ}{\ensuremath{\mathbb{Q}}}
\newunicodechar{𝔻}{\ensuremath{\mathbb{D}}}
\newunicodechar{α}{\ensuremath{\alpha}}
\newunicodechar{β}{\ensuremath{\beta}}
\newunicodechar{γ}{\ensuremath{\gamma}}
\newunicodechar{δ}{\ensuremath{\delta}}
\newunicodechar{ε}{\ensuremath{\varepsilon}}
\newunicodechar{θ}{\ensuremath{\theta}}
\newunicodechar{λ}{\ensuremath{\lambda}}
\newunicodechar{μ}{\ensuremath{\mu}}
\newunicodechar{σ}{\ensuremath{\sigma}}
\newunicodechar{τ}{\ensuremath{\tau}}
\newunicodechar{φ}{\ensuremath{\varphi}}
\newunicodechar{ψ}{\ensuremath{\psi}}
\newunicodechar{ω}{\ensuremath{\omega}}
\newunicodechar{Σ}{\ensuremath{\Sigma}}
% majuscules produites par \MakeUppercase dans les titres (α-aminés -> Α)
\newunicodechar{Α}{\ensuremath{\alpha}}
\newunicodechar{Β}{\ensuremath{\beta}}
\newunicodechar{Γ}{\ensuremath{\Gamma}}
\newunicodechar{Δ}{\ensuremath{\Delta}}
\newunicodechar{Ε}{\ensuremath{\varepsilon}}
\newunicodechar{Θ}{\ensuremath{\Theta}}
\newunicodechar{Λ}{\ensuremath{\Lambda}}
\newunicodechar{Μ}{\ensuremath{\mu}}
\newunicodechar{Τ}{\ensuremath{\tau}}
\newunicodechar{Φ}{\ensuremath{\Phi}}
\newunicodechar{Ψ}{\ensuremath{\Psi}}
\newunicodechar{Ω}{\ensuremath{\Omega}}
\newunicodechar{↦}{\ensuremath{\mapsto}}
\newunicodechar{∈}{\ensuremath{\in}}
\newunicodechar{∉}{\ensuremath{\notin}}
\newunicodechar{∧}{\ensuremath{\wedge}}
\newunicodechar{⇔}{\ensuremath{\Leftrightarrow}}
\newunicodechar{⟺}{\ensuremath{\Longleftrightarrow}}
\newunicodechar{⇒}{\ensuremath{\Rightarrow}}
\newunicodechar{⟹}{\ensuremath{\Longrightarrow}}
\newunicodechar{∘}{\ensuremath{\circ}}
\newunicodechar{∪}{\ensuremath{\cup}}
\newunicodechar{∩}{\ensuremath{\cap}}
\newunicodechar{≡}{\ensuremath{\equiv}}
\newunicodechar{⊂}{\ensuremath{\subset}}
\newunicodechar{⊥}{\ensuremath{\perp}}
\newunicodechar{∃}{\ensuremath{\exists}}
\newunicodechar{∀}{\ensuremath{\forall}}
\newunicodechar{∛}{\ensuremath{\sqrt[3]{\,}}}
\newunicodechar{∜}{\ensuremath{\sqrt[4]{\,}}}
\newunicodechar{⃗}{\ensuremath{{}^{\to}}}
\newunicodechar{ⁿ}{\ifmmode^{n}\else\textsuperscript{\ensuremath{n}}\fi}
\newunicodechar{ˣ}{\ifmmode^{x}\else\textsuperscript{\ensuremath{x}}\fi}
\newunicodechar{ᵇ}{\ifmmode^{b}\else\textsuperscript{\ensuremath{b}}\fi}
\newunicodechar{ʳ}{\ifmmode^{r}\else\textsuperscript{\ensuremath{r}}\fi}
\newunicodechar{ᵘ}{\ifmmode^{u}\else\textsuperscript{\ensuremath{u}}\fi}
\newunicodechar{ʸ}{\ifmmode^{y}\else\textsuperscript{\ensuremath{y}}\fi}
\newunicodechar{ᵃ}{\ifmmode^{a}\else\textsuperscript{\ensuremath{a}}\fi}
\newunicodechar{ᵉ}{\ifmmode^{e}\else\textsuperscript{\ensuremath{e}}\fi}
\newunicodechar{ᵏ}{\ifmmode^{k}\else\textsuperscript{\ensuremath{k}}\fi}
\newunicodechar{ᵖ}{\ifmmode^{p}\else\textsuperscript{\ensuremath{p}}\fi}
\newunicodechar{ᴺ}{\ifmmode^{N}\else\textsuperscript{\ensuremath{N}}\fi}
\newunicodechar{ᶜ}{\ifmmode^{c}\else\textsuperscript{\ensuremath{c}}\fi}
\newunicodechar{ʰ}{\ifmmode^{h}\else\textsuperscript{\ensuremath{h}}\fi}
\newunicodechar{ᵠ}{\ifmmode^{q}\else\textsuperscript{\ensuremath{q}}\fi}
\newunicodechar{ₙ}{\ifmmode_{n}\else\textsubscript{\ensuremath{n}}\fi}
\newunicodechar{ₐ}{\ifmmode_{a}\else\textsubscript{\ensuremath{a}}\fi}
\newunicodechar{ᵢ}{\ifmmode_{i}\else\textsubscript{\ensuremath{i}}\fi}
\newunicodechar{ᵣ}{\ifmmode_{r}\else\textsubscript{\ensuremath{r}}\fi}
\newunicodechar{ₖ}{\ifmmode_{k}\else\textsubscript{\ensuremath{k}}\fi}
\newunicodechar{ᵦ}{\ifmmode_{\beta}\else\textsubscript{\ensuremath{\beta}}\fi}
\newunicodechar{ₓ}{\ifmmode_{x}\else\textsubscript{\ensuremath{x}}\fi}
\newfontfamily\popdisplay{ArchivoBlack-Regular.ttf}[Path=../../../fonts/]
\newfontfamily\poplogo{SpaceGrotesk-Bold.ttf}[Path=../../../fonts/]
\newfontfamily\popsemi{PlusJakartaSans-SemiBold.ttf}[Path=../../../fonts/]
\newfontfamily\popxbold{PlusJakartaSans-ExtraBold.ttf}[Path=../../../fonts/]
\newfontfamily\popxboldls{PlusJakartaSans-ExtraBold.ttf}[Path=../../../fonts/,LetterSpace=3.0]

\setlist[enumerate]{leftmargin=1.7em,itemsep=5pt,topsep=4pt}
\setlist[itemize]{leftmargin=1.7em,itemsep=4pt,topsep=4pt,label={\color{poporange}\textbullet}}
\setlength{\parindent}{0pt}
\setlength{\parskip}{5pt}
\renewcommand{\arraystretch}{1.25}

% pgfplots : style Pop par défaut
\pgfplotsset{
  every axis/.append style={axis line style={popink,line width=1pt},tick style={popink},
    grid style={popline,line width=0.5pt},label style={font=\small},tick label style={font=\footnotesize}},
  popcurve/.style={poporange,line width=1.8pt,no markers,smooth},
}
% Même style disponible dans un \draw TikZ simple (hors pgfplots)
\tikzset{popcurve/.style={poporange,line width=1.8pt}}

% ============================================================
% Pill / squiggle
% ============================================================
\newcommand{\poppill}[3]{%
  \tikz[baseline=-0.55ex]\node[draw=popink,line width=1.2pt,rounded corners=2.6pt,%
    fill=#2,inner xsep=6.5pt,inner ysep=3pt]{%
    \popxboldls\fontsize{7}{7}\selectfont\color{#3}\MakeUppercase{#1}};%
}
\newcommand{\popsquiggle}[1]{%
  \tikz[baseline=-0.4ex]{
    \draw[#1,line width=2.6pt,line cap=round]
      plot [smooth] coordinates {(0,0.09) (0.34,0.01) (0.68,0.21) (1.02,0.01) (1.36,0.10)};
  }%
}
% Mot-clé surligné (marqueur orange sous le texte)
\newcommand{\cle}[1]{{\popxbold\color{popdark}#1}}

% ============================================================
% En-tête / pied
% ============================================================
\pagestyle{fancy}
\fancyhf{}
\renewcommand{\headrulewidth}{0pt}
\renewcommand{\footrulewidth}{0pt}
\lhead{\raisebox{-0.15ex}{\poplogo\fontsize{13}{13}\selectfont\color{popink}OnBuch\color{poporange}+}}
\rhead{\poppill{\DOCMATIERE\ \textbullet\ \DOCNIVEAU\ \textbullet\ Leçon \DOCLECON}{white}{popink}}
\lfoot{\popxbold\fontsize{7.6}{7.6}\selectfont\color{popmuted}\MakeUppercase{Cours signé OnBuch+}\ \textbullet\ {\popsemi\DOCTITRE}}
\rfoot{\tikz[baseline=-0.6ex]\node[rounded corners=8.5pt,fill=popink,minimum height=17pt,minimum width=17pt,inner xsep=5pt,inner sep=0]{\popxbold\fontsize{8}{8}\selectfont\color{popcream}\thepage\,/\,\pageref*{LastPage}};}

% ============================================================
% Titres
% ============================================================
\newcommand{\chaptitre}[2]{%
  \begin{tcolorbox}[enhanced,colback=poporange,colframe=popink,boxrule=2.6pt,arc=11pt,
      drop shadow=popink,left=15pt,right=15pt,top=12pt,bottom=11pt,before skip=3pt,after skip=18pt]
    \poppill{#2}{popink}{popcream}\par\vspace{9pt}
    {\raggedright\hyphenpenalty=10000\exhyphenpenalty=10000\popdisplay\fontsize{22}{24}\selectfont\color{popcream}\MakeUppercase{#1}\par}
  \end{tcolorbox}
}
% Section numérotée : pastille orange avec le numéro + titre Archivo
\newcounter{popsec}
\newcommand{\coursec}[1]{%
  \stepcounter{popsec}\setcounter{popsub}{0}%
  \par\vspace{16pt}\noindent
  \begin{minipage}{\linewidth}
  \tikz[baseline=-0.7ex]\node[draw=popink,line width=1.6pt,fill=poporange,rounded corners=4pt,minimum size=22pt,inner sep=0]{\popdisplay\fontsize{12}{12}\selectfont\color{popcream}\thepopsec};%
  \hspace{8pt}{\popdisplay\fontsize{15}{17}\selectfont\color{popink}\MakeUppercase{#1}}\par
  \vspace{4pt}\noindent{\color{poporange}\rule{36pt}{3.6pt}}
  \end{minipage}\par\nopagebreak\vspace{8pt}
}
\newcounter{popsub}
\newcommand{\courssub}[1]{%
  \stepcounter{popsub}%
  \par\vspace{9pt}\noindent{\popxbold\fontsize{11.5}{13}\selectfont\color{popink}{\color{poporange}\thepopsec.\thepopsub}\hspace{6pt}#1}\par\nopagebreak\vspace{4pt}
}

% ============================================================
% Boîtes pédagogiques — même recette : fond blanc, bordure épaisse colorée,
% ombre dure encre, badge plein. Titre optionnel [#1].
% ============================================================
\tcbset{popbase/.style={breakable,enhanced,colback=white,boxrule=2.3pt,arc=9pt,
  shadow={3pt}{-3pt}{0pt}{popink},left=14pt,right=14pt,top=11pt,bottom=11pt,before skip=11pt,after skip=15pt,
  fontupper=\popsemi\fontsize{10.6}{14.5}\selectfont\color{popink},
  fontlower=\popsemi\fontsize{10.6}{14.5}\selectfont\color{popink}}}
% Coche dessinée (le glyphe ✓ n'existe pas dans les polices du document)
\newcommand{\popcheck}[1]{\tikz[baseline=-0.5ex]\draw[#1,line width=1.6pt,line cap=round,line join=round](-0.13,0.0)--(-0.04,-0.09)--(0.13,0.1);}
\newcommand{\popboxhead}[3]{\poppill{#1}{#2}{white}\ifx\relax#3\relax\else\hspace{6pt}{\popxbold\fontsize{10.6}{12}\selectfont\color{popink}#3}\fi\par\vspace{7pt}}

\newtcolorbox{aretenir}[1][]{popbase,colframe=popgreen,before upper={\popboxhead{À retenir}{popgreen}{#1}}}
\newtcolorbox{exemplebox}[1][]{popbase,colframe=poporange,before upper={\popboxhead{Exemple}{poporange}{#1}}}
\newtcolorbox{definition}[1][]{popbase,colframe=popblue,before upper={\popboxhead{Définition}{popblue}{#1}}}
\newtcolorbox{propriete}[1][]{popbase,colframe=poppurple,before upper={\popboxhead{Propriété}{poppurple}{#1}}}
\newtcolorbox{methode}[1][]{popbase,colframe=popgold,before upper={\popboxhead{Méthode}{popgold}{#1}}}
\newtcolorbox{attention}[1][]{popbase,colframe=poppink,before upper={\popboxhead{Attention}{poppink}{#1}}}
\newtcolorbox{experience}[1][]{popbase,colframe=popblue,colback=popblueL,before upper={\popboxhead{Expérience}{popblue}{#1}}}
% « Le savais-tu ? » : carte violette pleine (contexte, culture, Cameroun)
\newtcolorbox{savaistu}[1][]{popbase,colback=poppurple,colframe=popink,
  fontupper=\popsemi\fontsize{10.4}{14.2}\selectfont\color{white},
  before upper={\poppill{Le savais-tu\,?}{white}{poppurple}\ifx\relax#1\relax\else\hspace{6pt}{\popxbold\color{white}#1}\fi\par\vspace{7pt}}}
% Exercice résolu : énoncé en haut, \tcblower puis la solution sur fond crème
\newcounter{popexo}\newcounter{popexr}
\newtcolorbox{exoresolu}[1][]{popbase,colframe=poporange,colbacklower=popcream,
  segmentation style={popink,line width=1.2pt,dashed},
  before upper={\stepcounter{popexr}\popboxhead{Exercice résolu \thepopexr}{poporange}{#1}},
  before lower={\poppill{Solution}{popink}{popcream}\par\vspace{6pt}}}
% Exercice (non résolu en place) — la correction est dans la partie Corrigés
\newtcolorbox{exercice}[1][]{popbase,colframe=popink,
  before upper={\stepcounter{popexo}\popboxhead{Exercice \thepopexo}{popink}{#1}}}
% Corrigé
\newtcolorbox{corrige}[1][]{popbase,colframe=popgreen,colback=popgreenL,
  before upper={\popboxhead{Corrigé}{popgreen}{#1}}}
% Objectifs / prérequis / bilan
\newtcolorbox{objectifs}{popbase,colframe=popink,colback=poporangeL,
  before upper={\popboxhead{Objectifs de la leçon}{popink}{}}}
\newtcolorbox{prerequis}{popbase,colframe=popink,colback=popblueL,
  before upper={\popboxhead{Prérequis}{popblue}{}}}
\newtcolorbox{bilan}{popbase,colframe=popink,colback=white,boxrule=2.8pt,
  before upper={\popboxhead{Fiche bilan}{poporange}{L'essentiel à connaître pour l'examen}}}
% Figure encadrée avec légende
\newtcolorbox{popfigure}[1][]{enhanced,breakable=false,colback=white,colframe=popline,boxrule=1.4pt,arc=8pt,
  left=10pt,right=10pt,top=10pt,bottom=8pt,before skip=10pt,after skip=12pt,halign=center,
  lower separated=false,halign lower=center,
  fontlower=\popsemi\fontsize{9}{11.5}\selectfont\color{popmuted},
  #1}
\newcommand{\legende}[1]{\par\vspace{6pt}{\popsemi\fontsize{9}{11.5}\selectfont\color{popmuted}#1}}

% ============================================================
% Signature OnBuch+
% ============================================================
\newcommand{\poplogoinline}[1]{{\poplogo\fontsize{#1}{#1}\selectfont\color{popink}OnBuch\color{poporange}+}}
\newcommand{\signatureonbuch}{%
  \begin{tcolorbox}[enhanced,colback=popink,colframe=popink,boxrule=2.2pt,arc=10pt,
      drop shadow=poporange,left=14pt,right=14pt,top=10pt,bottom=10pt,before skip=0pt,after skip=0pt]
    \tikz[baseline=-0.7ex]\node[circle,fill=popgreen,draw=popcream,line width=1.3pt,minimum size=22pt,inner sep=0]{\popcheck{white}};%
    \hspace{9pt}\begin{minipage}[c]{0.62\linewidth}
      {\popxbold\fontsize{12}{13}\selectfont\color{popcream}Signé\ \textbullet\ {\poplogo OnBuch\color{poporange}+}}\par\vspace{2pt}
      {\raggedright\popsemi\fontsize{8}{10.5}\selectfont\color{popline}\MakeUppercase{Équipe pédagogique OnBuch+}\\\MakeUppercase{Conforme au programme officiel MINESEC}\par}
    \end{minipage}\hfill
    {\poplogo\fontsize{22}{22}\selectfont\color{popcream}OnBuch\color{poporange}+}
  \end{tcolorbox}%
}

% ============================================================
% Page de garde
% #1 titre · #2 matière · #3 niveau · #4 module · #5 n° leçon · #6 durée ·
% #7 description / objectifs (texte ou itemize)
% ============================================================
\newcommand{\pagegarde}[7]{%
  \thispagestyle{empty}
  \newgeometry{a4paper,margin=1.7cm,top=1.6cm,bottom=1.4cm}%
  \begin{tikzpicture}[remember picture,overlay]
    \fill[poporange!18] ([xshift=-3.2cm,yshift=-3.6cm]current page.north east) circle (4.6cm);
    \fill[poppurple!14] ([xshift=2.4cm,yshift=7.6cm]current page.south west) circle (3.8cm);
  \end{tikzpicture}%
  {\poplogo\fontsize{30}{30}\selectfont\color{popink}OnBuch\color{poporange}+}\hfill
  \raisebox{0.6ex}{\poppill{Cours signé \textbullet\ Premium}{popink}{popcream}}\par
  \vspace{1pt}\popsquiggle{poppurple}\par
  \vspace{8pt}
  \begin{tcolorbox}[enhanced,colback=poporange,colframe=popink,boxrule=2.8pt,arc=13pt,
      drop shadow=popink,left=18pt,right=18pt,top=12pt,bottom=13pt,after skip=10pt]
    \poppill{#2 \textbullet\ #3}{popink}{popcream}\hspace{5pt}\poppill{#4}{white}{popink}\par\vspace{8pt}
    {\popdisplay\fontsize{14}{14}\selectfont\color{popink}LEÇON #5}\par\vspace{4pt}
    {\raggedright\hyphenpenalty=10000\exhyphenpenalty=10000\popdisplay\fontsize{25}{27}\selectfont\color{popcream}\MakeUppercase{#1}\par}
    \vspace{8pt}
    {\popxbold\fontsize{9.5}{9.5}\selectfont\color{popink}\MakeUppercase{Durée conseillée : #6}}
  \end{tcolorbox}
  \begin{tcolorbox}[enhanced,colback=white,colframe=popink,boxrule=2.2pt,arc=10pt,
      drop shadow=popink,left=16pt,right=16pt,top=10pt,bottom=10pt,after skip=8pt]
    \poppill{Dans cette leçon}{poppurple}{white}\par\vspace{5pt}
    \popsemi\fontsize{9.3}{12.6}\selectfont\color{popink}#7
  \end{tcolorbox}
  \noindent\begin{minipage}[t]{0.487\linewidth}
    \begin{tcolorbox}[enhanced,colback=popgreen,colframe=popink,boxrule=2.2pt,arc=10pt,
        drop shadow=popink,left=11pt,right=11pt,top=10pt,bottom=10pt,height=3.1cm,valign=center]
      \tcbox[colback=white,colframe=popink,boxrule=1.3pt,arc=5pt,boxsep=0pt,nobeforeafter,
        left=3pt,right=3pt,top=3pt,bottom=3pt,valign=center]{\qrcode[height=1.9cm]{https://play.google.com/store/apps/details?id=com.luvvix.onbuch&hl=fr}}%
      \hspace{8pt}\begin{minipage}[c]{0.5\linewidth}
        {\popdisplay\fontsize{11}{12.5}\selectfont\color{white}\MakeUppercase{Télécharge OnBuch}}\par\vspace{4pt}
        {\raggedright\popsemi\fontsize{8.4}{11}\selectfont\color{white}Gratuit \textbullet\ Google Play\\Tuteur IA, annales, concours, résultats\par}
      \end{minipage}
    \end{tcolorbox}
  \end{minipage}\hfill
  \begin{minipage}[t]{0.487\linewidth}
    \begin{tcolorbox}[enhanced,colback=poppurple,colframe=popink,boxrule=2.2pt,arc=10pt,
        drop shadow=popink,left=11pt,right=11pt,top=10pt,bottom=10pt,height=3.1cm,valign=center]
      \tikz[baseline=-0.7ex]\node[circle,fill=white,draw=popink,line width=1.3pt,minimum size=1.6cm,inner sep=0]{\popdisplay\fontsize{24}{24}\selectfont\color{poppurple}+};%
      \hspace{8pt}\begin{minipage}[c]{0.6\linewidth}
        {\popdisplay\fontsize{11}{12.5}\selectfont\color{white}\MakeUppercase{Passe à OnBuch+}}\par\vspace{4pt}
        {\raggedright\popsemi\fontsize{8.4}{11}\selectfont\color{white}Tous les cours signés, Tuteur IA illimité, fiches PDF — onglet OnBuch+ de l'app\par}
      \end{minipage}
    \end{tcolorbox}
  \end{minipage}
  \vspace{8pt}
  \popcontenu
  \vfill
  \signatureonbuch
  \restoregeometry
  \newpage
}

% Rangée « Ce cours contient » (page de garde)
\newcommand{\poptile}[3]{% #1 couleur #2 gros texte #3 légende
  \begin{tcolorbox}[enhanced,colback=#1,colframe=popink,boxrule=2pt,arc=9pt,shadow={2.5pt}{-2.5pt}{0pt}{popink},
      left=6pt,right=6pt,top=9pt,bottom=9pt,halign=center,valign=center,height=1.75cm,width=\linewidth,nobeforeafter]
    {\popdisplay\fontsize{12.5}{13}\selectfont\color{white}\MakeUppercase{#2}}\par\vspace{3pt}
    {\centering\popsemi\fontsize{7.8}{9.5}\selectfont\color{white}#3\par}
  \end{tcolorbox}}
\newcommand{\popcontenu}{%
  \par\noindent{\popxboldls\fontsize{7.5}{7.5}\selectfont\color{popmuted}CE COURS SIGNÉ CONTIENT}\par\vspace{7pt}
  \noindent\begin{minipage}[t]{0.235\linewidth}\poptile{popblue}{Cours}{complet, illustré, pas à pas}\end{minipage}\hfill
  \begin{minipage}[t]{0.235\linewidth}\poptile{poporange}{Méthodes}{et exercices résolus}\end{minipage}\hfill
  \begin{minipage}[t]{0.235\linewidth}\poptile{poppink}{Exercices}{type examen + corrigés}\end{minipage}\hfill
  \begin{minipage}[t]{0.235\linewidth}\poptile{popgreen}{Fiche bilan}{l'essentiel à retenir}\end{minipage}\par}

% Sommaire
\newtcolorbox{sommaire}{enhanced,colback=white,colframe=popink,boxrule=2.2pt,arc=10pt,
  drop shadow=popink,left=18pt,right=18pt,top=13pt,bottom=13pt,before skip=6pt,after skip=20pt,
  before upper={\poppill{Sommaire}{poppurple}{white}\par\vspace{9pt}\popsemi\fontsize{10.6}{14.5}\selectfont\color{popink}}}

% Signature de fin de cours — poussée tout en bas de la dernière page
\newcommand{\findecours}{%
  \par\vfill\nopagebreak
  \begin{center}\popsquiggle{poporange}\end{center}\vspace{4pt}
  \signatureonbuch
}
\AtBeginDocument{\def\llbracket{\mathopen{[\mkern-3.5mu[}}\def\rrbracket{\mathclose{]\mkern-3.5mu]}}} % intervalles d'entiers (glyphe ⟦ absent de la police)
% Glyphes absents de Fira Math : redéfinis à partir de glyphes disponibles
\AtBeginDocument{%
  \def\setminus{\mathbin{\backslash}}\let\smallsetminus\setminus
  \def\vdots{\mathord{\vbox{\baselineskip3.2pt\lineskiplimit0pt\kern4pt\hbox{.}\hbox{.}\hbox{.}}}}%
  \def\ddots{\mathinner{\mkern1mu\raise6.5pt\hbox{.}\mkern2mu\raise3.6pt\hbox{.}\mkern2mu\raise0.7pt\hbox{.}\mkern1mu}}%
  \def\triangle{\mathord{\scalebox{0.9}{$\symup{\Delta}$}}}%
  \def\nabla{\mathord{\raisebox{\depth}{\scalebox{1}[-1]{$\symup{\Delta}$}}}}%
  \def\checkmark{\popcheck{popdark}}%
  \def\star{\mathbin{\ast}}\def\bigstar{\mathord{\ast}}%
  \def\diamond{\mathbin{\scalebox{0.6}{$\Diamond$}}}%
}

%% FIGFIX-BEGIN v5 — correctifs globaux des figures (docs/onbuch-generation/tools/figfix)
\usepackage{adjustbox}\usepackage{etoolbox}\tcbuselibrary{raster}
% toute figure TikZ/pgfplots plus large que la ligne est réduite à la largeur disponible
\BeforeBeginEnvironment{tikzpicture}{\begin{adjustbox}{max width=\linewidth}}
\AfterEndEnvironment{tikzpicture}{\end{adjustbox}}
% légendes pgfplots lisibles par-dessus les courbes
\pgfplotsset{every axis/.append style={legend style={fill=white,fill opacity=0.92,text opacity=1,draw=black!15,inner sep=3pt}}}
% étiquettes d'arêtes (syntaxe edge["..."]) avec fond blanc
\tikzset{every edge quotes/.append style={fill=white,inner sep=1.2pt}}
%% FIGFIX-END
\begin{document}

\pagegarde{Matrices 2×2 et applications linéaires du plan}{Mathématiques}{Tle TI}{Configurations et transformations élémentaires du plan}{20}{6 h}{Cette leçon introduit les matrices carrées d'ordre 2, leurs opérations, le déterminant et l'inversion, puis les relie aux applications linéaires du plan. Elle fournit un outil unifié pour résoudre les systèmes linéaires 2×2 et étudier les transformations géométriques du plan, compétence centrale au Bac TI.
\vspace{4pt}
\begin{multicols}{2}\small
\begin{itemize}
\item À la fin de cette leçon, je sais effectuer la somme de deux matrices 2×2, le produit d'une matrice par un réel et le produit de deux matrices 2×2
\item À la fin de cette leçon, je sais calculer le déterminant d'une matrice 2×2 et décider si elle est inversible
\item À la fin de cette leçon, je sais calculer l'inverse d'une matrice 2×2 inversible et vérifier le résultat
\item À la fin de cette leçon, je sais écrire un système linéaire 2×2 sous forme matricielle AX = B et le résoudre à l'aide de la matrice inverse
\item À la fin de cette leçon, je sais déterminer la matrice d'une application linéaire du plan dans une base donnée
\item À la fin de cette leçon, je sais calculer les coordonnées de l'image d'un vecteur par produit matriciel
\end{itemize}\end{multicols}}

\begin{sommaire}
\begin{multicols}{2}
\begin{enumerate}
\item Matrices 2×2 : définition et opérations
\item Déterminant et inverse d'une matrice 2×2
\item Matrices et systèmes linéaires 2×2
\item Matrice d'une application linéaire du plan
\item Opérations sur les applications linéaires et automorphismes
\item Applications géométriques et synthèse
\item Activité d'intégration
\item Méthodes et exercices résolus
\item Exercices
\item Corrigés des exercices
\item Fiche bilan
\end{enumerate}
\end{multicols}
\end{sommaire}

\begin{objectifs}
\begin{itemize}
\item À la fin de cette leçon, je sais effectuer la somme de deux matrices 2×2, le produit d'une matrice par un réel et le produit de deux matrices 2×2
\item À la fin de cette leçon, je sais calculer le déterminant d'une matrice 2×2 et décider si elle est inversible
\item À la fin de cette leçon, je sais calculer l'inverse d'une matrice 2×2 inversible et vérifier le résultat
\item À la fin de cette leçon, je sais écrire un système linéaire 2×2 sous forme matricielle AX = B et le résoudre à l'aide de la matrice inverse
\item À la fin de cette leçon, je sais déterminer la matrice d'une application linéaire du plan dans une base donnée
\item À la fin de cette leçon, je sais calculer les coordonnées de l'image d'un vecteur par produit matriciel
\item À la fin de cette leçon, je sais déterminer la matrice d'une somme et d'une composée d'applications linéaires
\item À la fin de cette leçon, je sais déterminer la matrice de l'application réciproque d'un automorphisme et l'interpréter géométriquement
\end{itemize}
\end{objectifs}

\begin{prerequis}
\begin{itemize}
\item Vecteurs du plan : coordonnées, opérations, colinéarité (Première)
\item Résolution de systèmes linéaires 2×2 par substitution ou combinaison (Seconde/Première)
\item Notion d'application et de composée d'applications (Première)
\item Transformations du plan : symétries, homothéties, projections (Première)
\item Calcul littéral et manipulation de fractions
\end{itemize}
\end{prerequis}

\newpage

\coursec{Matrices $2\times 2$ : définition et opérations}

\textbf{Situation d'accroche.} Maman Ngo Bell tient un dépôt de boissons à Douala. Chaque vendredi soir, son stock se transforme : les ventes de la semaine font baisser le nombre de casiers de jus et d'eau minérale, puis la livraison du samedi matin le fait remonter. Elle rêve d'un outil qui, à partir du stock de ce vendredi, lui donnerait \emph{automatiquement} celui du vendredi suivant, sans tout recompter casier par casier. À Yaoundé, un opérateur de transfert d'argent mobile connaît les totaux encaissés et les commissions prélevées sur deux transactions, mais il a égaré les montants initiaux : il doit \emph{remonter} le calcul.

Ces deux problèmes, en apparence très différents, reposent sur la même idée mathématique : transformer des \emph{couples de nombres} par des règles simples (additions et multiplications), puis éventuellement \emph{inverser} la transformation. L'outil qui code ces transformations s'appelle une \cle{matrice} : une simple grille de quatre nombres.

\textbf{Problématique.} Comment une grille de quatre nombres peut-elle résoudre un système d'équations, décrire une rotation du plan, ou retrouver une information perdue ? Tout commence par apprendre à \emph{calculer} avec ces grilles : c'est l'objet de cette première section.

\courssub{Qu'est-ce qu'une matrice carrée d'ordre 2 ?}

\textbf{Intuition.} Revenons au dépôt de Maman Ngo Bell. Ses ventes du vendredi peuvent se résumer ainsi :

\begin{center}
\begin{tabularx}{\linewidth}{|l|X|X|}
\hline
\rowcolor{popblueL} & Jus (casiers) & Eau minérale (casiers) \\
\hline
Semaine 1 & $12$ & $20$ \\
\hline
Semaine 2 & $15$ & $18$ \\
\hline
\end{tabularx}
\end{center}

Si l'on enlève les étiquettes, il reste l'essentiel : quatre nombres rangés en carré. C'est exactement cela, une matrice.

\begin{definition}[Matrice carrée d'ordre 2]
On appelle \cle{matrice carrée d'ordre 2} tout tableau de quatre nombres réels rangés en $2$ lignes et $2$ colonnes :
\[ A = \begin{pmatrix} a & b \\ c & d \end{pmatrix} \]
Les nombres $a$, $b$, $c$, $d$ sont les \cle{coefficients} de la matrice. Le coefficient situé sur la $i$-ème ligne et la $j$-ème colonne se note $a_{ij}$. Ainsi $a_{11}=a$, $a_{12}=b$, $a_{21}=c$, $a_{22}=d$.
\end{definition}

Le vocabulaire est important et tombe souvent au Bac :

\begin{itemize}
\item la \cle{première ligne} est $(a\ \ b)$, la \cle{deuxième ligne} est $(c\ \ d)$ ;
\item la \cle{première colonne} est $\begin{pmatrix} a \\ c \end{pmatrix}$, la \cle{deuxième colonne} est $\begin{pmatrix} b \\ d \end{pmatrix}$ ;
\item la \cle{diagonale principale} est formée des coefficients $a$ et $d$ (en haut à gauche, en bas à droite) ; l'autre diagonale, formée de $b$ et $c$, est la diagonale secondaire.
\end{itemize}

\begin{popfigure}
\begin{center}
\begin{tikzpicture}[scale=1.15, every node/.style={font=\large}]
% matrice
\node (a) at (0,1) {$a$};
\node (b) at (1.6,1) {$b$};
\node (c) at (0,0) {$c$};
\node (d) at (1.6,0) {$d$};
% parenthèses
\draw[popdark, thick] (-0.6,1.5) -- (-0.8,1.5) -- (-0.8,-0.5) -- (-0.6,-0.5);
\draw[popdark, thick] (2.2,1.5) -- (2.4,1.5) -- (2.4,-0.5) -- (2.2,-0.5);
% diagonale principale
\draw[poporange, very thick, dashed] (a.north west) -- (d.south east);
% annotations lignes
\draw[popblue, thick, -{Stealth}] (-1.9,1) -- (-0.95,1);
\node[popblue, left] at (-1.9,1) {ligne 1};
\draw[popblue, thick, -{Stealth}] (-1.9,0) -- (-0.95,0);
\node[popblue, left] at (-1.9,0) {ligne 2};
% annotations colonnes
\draw[popgreen, thick, -{Stealth}] (0,2.1) -- (0,1.4);
\node[popgreen, above] at (0,2.1) {colonne 1};
\draw[popgreen, thick, -{Stealth}] (1.6,2.1) -- (1.6,1.4);
\node[popgreen, above] at (1.6,2.1) {colonne 2};
% diagonale label
\node[poporange, right, align=left] at (3.0,0.5) {diagonale\\ principale};
\draw[poporange, thick, -{Stealth}] (3.0,0.45) -- (1.75,0.35);
\end{tikzpicture}
\end{center}
\legende{Anatomie d'une matrice carrée d'ordre 2 : lignes, colonnes et diagonale principale.}
\end{popfigure}

\begin{definition}[Égalité, matrice nulle, matrice identité]
Soient $A = \begin{pmatrix} a & b \\ c & d \end{pmatrix}$ et $A' = \begin{pmatrix} a' & b' \\ c' & d' \end{pmatrix}$.
\begin{itemize}
\item On dit que $A$ et $A'$ sont \cle{égales}, et on écrit $A = A'$, lorsque leurs coefficients sont égaux \emph{place par place} : $a=a'$, $b=b'$, $c=c'$, $d=d'$.
\item La \cle{matrice nulle}, notée $O$ (ou $0_2$), est la matrice dont tous les coefficients sont nuls : $O = \begin{pmatrix} 0 & 0 \\ 0 & 0 \end{pmatrix}$.
\item La \cle{matrice identité} (ou matrice unité), notée $I_2$, est la matrice qui porte des $1$ sur la diagonale principale et des $0$ ailleurs : $I_2 = \begin{pmatrix} 1 & 0 \\ 0 & 1 \end{pmatrix}$.
\end{itemize}
\end{definition}

\begin{attention}[L'égalité se lit coefficient par coefficient]
Dire que $\begin{pmatrix} x & 2 \\ -1 & y \end{pmatrix} = \begin{pmatrix} 3 & 2 \\ -1 & 5 \end{pmatrix}$, c'est affirmer \emph{en une seule phrase} que $x = 3$ \textbf{et} $y = 5$. Une égalité matricielle $2\times 2$ contient donc \emph{quatre} égalités numériques. C'est un réflexe précieux pour résoudre des équations où l'inconnue est une matrice.
\end{attention}

\courssub{Somme de deux matrices et produit par un réel}

\textbf{Intuition.} Maman Ngo Bell veut additionner ses ventes de deux semaines : casiers de jus avec casiers de jus, casiers d'eau avec casiers d'eau. On additionne donc les tableaux \emph{case par case}. De même, si les ventes doublent (période de fêtes de fin d'année !), chaque case est multipliée par $2$.

\begin{definition}[Somme de deux matrices]
La \cle{somme} de deux matrices $A = \begin{pmatrix} a & b \\ c & d \end{pmatrix}$ et $B = \begin{pmatrix} a' & b' \\ c' & d' \end{pmatrix}$ est la matrice obtenue en additionnant les coefficients de même position :
\[ A + B = \begin{pmatrix} a+a' & b+b' \\ c+c' & d+d' \end{pmatrix}. \]
\end{definition}

\begin{definition}[Produit d'une matrice par un réel]
Pour tout réel $k$, le \cle{produit} de $A$ par $k$ est la matrice obtenue en multipliant \emph{chaque} coefficient par $k$ :
\[ kA = \begin{pmatrix} ka & kb \\ kc & kd \end{pmatrix}. \]
On note $-A = (-1)A$ la matrice opposée, et $A - B = A + (-B)$.
\end{definition}

\begin{exemplebox}[Calculs de sommes et de produits par un réel]
Soient $A = \begin{pmatrix} 12 & 20 \\ 15 & 18 \end{pmatrix}$ et $B = \begin{pmatrix} 3 & 5 \\ 2 & 4 \end{pmatrix}$.
\begin{align*}
A + B &= \begin{pmatrix} 12+3 & 20+5 \\ 15+2 & 18+4 \end{pmatrix} = \begin{pmatrix} 15 & 25 \\ 17 & 22 \end{pmatrix}, \\
3B &= \begin{pmatrix} 9 & 15 \\ 6 & 12 \end{pmatrix}, \\
A - 2B &= \begin{pmatrix} 12-6 & 20-10 \\ 15-4 & 18-8 \end{pmatrix} = \begin{pmatrix} 6 & 10 \\ 11 & 10 \end{pmatrix}.
\end{align*}
\end{exemplebox}

\begin{propriete}[Propriétés de la somme et du produit par un réel — admises]
Pour toutes matrices $A$, $B$, $C$ d'ordre 2 et tous réels $k$, $k'$ :
\begin{itemize}
\item \textbf{Commutativité :} $A + B = B + A$ ;
\item \textbf{Associativité :} $(A + B) + C = A + (B + C)$ ;
\item $A + O = A$ et $A + (-A) = O$ ;
\item $k(A+B) = kA + kB$ ; \quad $(k+k')A = kA + k'A$ ; \quad $k(k'A) = (kk')A$ ; \quad $1A = A$.
\end{itemize}
Ces propriétés se vérifient sans peine sur des exemples (fais-le avec les matrices $A$ et $B$ ci-dessus !) ; on les admet dans le cas général. Elles découlent directement des propriétés de l'addition et de la multiplication des nombres réels, appliquées coefficient par coefficient.
\end{propriete}

\courssub{Produit de deux matrices : la règle « ligne $\times$ colonne »}

\textbf{Intuition.} Voici le point délicat — et le plus riche. Supposons que Maman Ngo Bell vende chaque casier de jus à $\num{12000}$ FCFA et chaque casier d'eau à $\num{5000}$ FCFA. Pour calculer la recette d'une semaine, elle multiplie les \emph{quantités} (une ligne du tableau des ventes) par les \emph{prix} (une colonne) et elle additionne : $12 \times \num{12000} + 20 \times \num{5000}$. C'est exactement le mécanisme du produit matriciel : \textbf{on fait se rencontrer une ligne et une colonne}.

\begin{definition}[Produit de deux matrices $2\times 2$]
Soient $A = \begin{pmatrix} a & b \\ c & d \end{pmatrix}$ et $B = \begin{pmatrix} a' & b' \\ c' & d' \end{pmatrix}$. Le \cle{produit} $AB$ est la matrice définie par :
\[ AB = \begin{pmatrix} aa' + bc' & ab' + bd' \\ ca' + dc' & cb' + dd' \end{pmatrix}. \]
Autrement dit, le coefficient de la $i$-ème ligne et de la $j$-ème colonne de $AB$ s'obtient en multipliant la $i$-ème \textbf{ligne} de $A$ par la $j$-ème \textbf{colonne} de $B$, terme à terme, puis en additionnant les résultats.
\end{definition}

\begin{popfigure}
\begin{center}
\begin{tikzpicture}[scale=1.0, every node/.style={font=\normalsize}]
% Matrice A
\node at (-3.4,2.3) {$A$};
\node (a11) at (-3.9,1.6) {$1$};
\node (a12) at (-3.1,1.6) {$2$};
\node (a21) at (-3.9,0.8) {$0$};
\node (a22) at (-3.1,0.8) {$3$};
\draw[popdark] (-4.35,1.95) -- (-4.5,1.95) -- (-4.5,0.45) -- (-4.35,0.45);
\draw[popdark] (-2.65,1.95) -- (-2.5,1.95) -- (-2.5,0.45) -- (-2.65,0.45);
% symbole fois
\node at (-1.9,1.2) {$\times$};
% Matrice B
\node at (-0.4,2.3) {$B$};
\node (b11) at (-0.9,1.6) {$2$};
\node (b12) at (-0.1,1.6) {$-1$};
\node (b21) at (-0.9,0.8) {$1$};
\node (b22) at (-0.1,0.8) {$1$};
\draw[popdark] (-1.35,1.95) -- (-1.5,1.95) -- (-1.5,0.45) -- (-1.35,0.45);
\draw[popdark] (0.35,1.95) -- (0.5,1.95) -- (0.5,0.45) -- (0.35,0.45);
% égal
\node at (1.1,1.2) {$=$};
% Matrice AB
\node at (2.6,2.3) {$AB$};
\node (p11) at (2.1,1.6) {$4$};
\node (p12) at (3.1,1.6) {$1$};
\node (p21) at (2.1,0.8) {$3$};
\node (p22) at (3.1,0.8) {$3$};
\draw[popdark] (1.65,1.95) -- (1.5,1.95) -- (1.5,0.45) -- (1.65,0.45);
\draw[popdark] (3.55,1.95) -- (3.7,1.95) -- (3.7,0.45) -- (3.55,0.45);
% flèches ligne 1 de A vers colonne 1 de B vers coefficient (1,1)
\draw[poporange, very thick, -{Stealth}] (a11.north) .. controls (-3.9,3.2) and (-0.9,3.4) .. (b11.north);
\draw[poporange, very thick, -{Stealth}] (a12.north) .. controls (-3.1,3.0) and (-0.9,3.0) .. (b21.south west);
\draw[poporange, very thick, -{Stealth}] (b11.east) .. controls (0.2,2.6) and (1.6,2.4) .. (p11.north);
% flèche ligne 2 de A vers colonne 2 de B vers coefficient (2,2)
\draw[popblue, very thick, -{Stealth}] (a21.south) .. controls (-3.9,-0.6) and (-0.1,-0.6) .. (b12.south);
\draw[popblue, very thick, -{Stealth}] (a22.south) .. controls (-3.1,-0.3) and (-0.1,-0.3) .. (b22.south);
\draw[popblue, very thick, -{Stealth}] (b22.east) .. controls (0.9,0.2) and (2.6,0.2) .. (p22.south);
% légendes
\node[poporange, align=center] at (-1.6,3.7) {ligne 1 de $A$ $\times$ colonne 1 de $B$ :\\ $1\times 2 + 2\times 1 = 4$};
\node[popblue, align=center] at (-1.6,-1.3) {ligne 2 de $A$ $\times$ colonne 2 de $B$ :\\ $0\times(-1) + 3\times 1 = 3$};
\end{tikzpicture}
\end{center}
\legende{La règle « ligne $\times$ colonne » : chaque coefficient du produit est la rencontre d'une ligne de $A$ et d'une colonne de $B$.}
\end{popfigure}

\begin{methode}[Calculer un produit $AB$ pas à pas]
Pour calculer $AB$ avec $A = \begin{pmatrix} a & b \\ c & d \end{pmatrix}$ et $B = \begin{pmatrix} a' & b' \\ c' & d' \end{pmatrix}$ :
\begin{enumerate}
\item \textbf{Coefficient $(1,1)$} : ligne 1 de $A$ contre colonne 1 de $B$ : $aa' + bc'$.
\item \textbf{Coefficient $(1,2)$} : ligne 1 de $A$ contre colonne 2 de $B$ : $ab' + bd'$.
\item \textbf{Coefficient $(2,1)$} : ligne 2 de $A$ contre colonne 1 de $B$ : $ca' + dc'$.
\item \textbf{Coefficient $(2,2)$} : ligne 2 de $A$ contre colonne 2 de $B$ : $cb' + dd'$.
\end{enumerate}
Astuce de vérification : pose ton doigt gauche sur la ligne de $A$, ton doigt droit sur la colonne de $B$, et fais-les glisser ensemble.
\end{methode}

\begin{exemplebox}[Un produit calculé en détail]
Soient $A = \begin{pmatrix} 1 & 2 \\ 0 & 3 \end{pmatrix}$ et $B = \begin{pmatrix} 2 & -1 \\ 1 & 1 \end{pmatrix}$. Calculons $AB$ :
\begin{align*}
AB &= \begin{pmatrix} 1\times 2 + 2\times 1 & 1\times(-1) + 2\times 1 \\ 0\times 2 + 3\times 1 & 0\times(-1) + 3\times 1 \end{pmatrix} \\
&= \begin{pmatrix} 2+2 & -1+2 \\ 0+3 & 0+3 \end{pmatrix} = \begin{pmatrix} 4 & 1 \\ 3 & 3 \end{pmatrix}.
\end{align*}
\end{exemplebox}

\courssub{Le produit matriciel n'est pas commutatif}

Voici la grande surprise de cette leçon, et l'une des erreurs les plus fréquentes au Baccalauréat.

\begin{exoresolu}[Montrer que $AB \neq BA$ en général]
Énoncé. Soient $A = \begin{pmatrix} 1 & 2 \\ 0 & 3 \end{pmatrix}$ et $B = \begin{pmatrix} 2 & -1 \\ 1 & 1 \end{pmatrix}$. Calculer $AB$ et $BA$, puis comparer.
\tcblower
Solution détaillée. Nous avons déjà calculé $AB = \begin{pmatrix} 4 & 1 \\ 3 & 3 \end{pmatrix}$. Calculons maintenant $BA$ avec la même règle ligne $\times$ colonne :
\begin{align*}
BA &= \begin{pmatrix} 2\times 1 + (-1)\times 0 & 2\times 2 + (-1)\times 3 \\ 1\times 1 + 1\times 0 & 1\times 2 + 1\times 3 \end{pmatrix} \\
&= \begin{pmatrix} 2-0 & 4-3 \\ 1+0 & 2+3 \end{pmatrix} = \begin{pmatrix} 2 & 1 \\ 1 & 5 \end{pmatrix}.
\end{align*}
Comparons : $AB = \begin{pmatrix} 4 & 1 \\ 3 & 3 \end{pmatrix}$ et $BA = \begin{pmatrix} 2 & 1 \\ 1 & 5 \end{pmatrix}$. Les coefficients $(1,1)$ valent $4$ et $2$ : les matrices sont différentes. Donc $AB \neq BA$.
\end{exoresolu}

\begin{attention}[Deux pièges mortels sur le produit matriciel]
\begin{enumerate}
\item \textbf{Le produit n'est pas commutatif.} En général $AB \neq BA$. Il est donc interdit d'écrire $AB$ à la place de $BA$, ou de « développer » $(A+B)^2$ en $A^2 + 2AB + B^2$ : en réalité $(A+B)^2 = A^2 + AB + BA + B^2$, et on ne peut pas regrouper $AB$ et $BA$ !
\item \textbf{Le produit ne se fait PAS coefficient par coefficient.} $\begin{pmatrix} 1 & 2 \\ 0 & 3 \end{pmatrix}\begin{pmatrix} 2 & -1 \\ 1 & 1 \end{pmatrix}$ ne vaut \emph{pas} $\begin{pmatrix} 2 & -2 \\ 0 & 3 \end{pmatrix}$. C'est la règle ligne $\times$ colonne qui s'applique, toujours.
\end{enumerate}
\end{attention}

\begin{propriete}[Associativité, distributivité, rôle de $I_2$ — admises, vérifiées sur exemples]
Pour toutes matrices $A$, $B$, $C$ d'ordre 2 et tout réel $k$ :
\begin{itemize}
\item \textbf{Associativité :} $(AB)C = A(BC)$ ;
\item \textbf{Distributivité à gauche :} $A(B+C) = AB + AC$ ;
\item \textbf{Distributivité à droite :} $(A+B)C = AC + BC$ ;
\item \textbf{Compatibilité avec les scalaires :} $k(AB) = (kA)B = A(kB)$ ;
\item \textbf{Élément neutre :} $AI_2 = I_2A = A$.
\end{itemize}
Ces propriétés sont admises dans le cas général ; la démonstration n'est pas exigible au Bac, mais tu dois savoir les \emph{vérifier} sur des exemples numériques.
\end{propriete}

\begin{exemplebox}[Vérification de la distributivité sur un exemple]
Prenons $A = \begin{pmatrix} 1 & 2 \\ 0 & 3 \end{pmatrix}$, $B = \begin{pmatrix} 2 & -1 \\ 1 & 1 \end{pmatrix}$ et $C = \begin{pmatrix} 1 & 0 \\ 2 & 1 \end{pmatrix}$.

D'une part : $B + C = \begin{pmatrix} 3 & -1 \\ 3 & 2 \end{pmatrix}$, donc
\[ A(B+C) = \begin{pmatrix} 1\times 3 + 2\times 3 & 1\times(-1) + 2\times 2 \\ 0\times 3 + 3\times 3 & 0\times(-1) + 3\times 2 \end{pmatrix} = \begin{pmatrix} 9 & 3 \\ 9 & 6 \end{pmatrix}. \]

D'autre part :
\[ AB = \begin{pmatrix} 4 & 1 \\ 3 & 3 \end{pmatrix}, \qquad AC = \begin{pmatrix} 1+4 & 0+2 \\ 0+6 & 0+3 \end{pmatrix} = \begin{pmatrix} 5 & 2 \\ 6 & 3 \end{pmatrix}, \]
\[ AB + AC = \begin{pmatrix} 4+5 & 1+2 \\ 3+6 & 3+3 \end{pmatrix} = \begin{pmatrix} 9 & 3 \\ 9 & 6 \end{pmatrix}. \]
On retrouve bien $A(B+C) = AB + AC$. Vérifie de même, à titre d'entraînement, que $AI_2 = A$ : chaque ligne de $A$ « rencontre » les colonnes de $I_2$ qui ne laissent passer que le bon coefficient.
\end{exemplebox}

\begin{aretenir}[Ce qu'il faut retenir des opérations]
\begin{itemize}
\item Somme et produit par un réel : \textbf{coefficient par coefficient}.
\item Produit de deux matrices : \textbf{ligne $\times$ colonne}.
\item Le produit est \textbf{associatif} et \textbf{distributif} sur la somme, mais \textbf{pas commutatif}.
\item $I_2$ est l'élément neutre du produit : $AI_2 = I_2A = A$ ; $O$ est l'élément neutre de la somme : $A + O = A$.
\end{itemize}
\end{aretenir}

\courssub{Puissances d'une matrice}

\begin{definition}[Puissances d'une matrice carrée]
Soit $A$ une matrice carrée d'ordre 2. On définit ses \cle{puissances} par :
\[ A^2 = A \times A, \qquad A^3 = A^2 \times A = A \times A \times A, \]
et plus généralement, pour tout entier $n \geq 1$, $A^n = \underbrace{A \times A \times \cdots \times A}_{n \text{ facteurs}}$. Par convention, $A^0 = I_2$ et $A^1 = A$.
\end{definition}

Grâce à l'associativité, cette définition est cohérente : $A^3 = (A^2)A = A(A^2)$, peu importe où l'on place les parenthèses.

\begin{exemplebox}[Calcul de $A^2$ et $A^3$]
Soit $A = \begin{pmatrix} 1 & 1 \\ 0 & 2 \end{pmatrix}$.
\[ A^2 = \begin{pmatrix} 1 & 1 \\ 0 & 2 \end{pmatrix}\begin{pmatrix} 1 & 1 \\ 0 & 2 \end{pmatrix} = \begin{pmatrix} 1+0 & 1+2 \\ 0+0 & 0+4 \end{pmatrix} = \begin{pmatrix} 1 & 3 \\ 0 & 4 \end{pmatrix}. \]
\[ A^3 = A^2 \times A = \begin{pmatrix} 1 & 3 \\ 0 & 4 \end{pmatrix}\begin{pmatrix} 1 & 1 \\ 0 & 2 \end{pmatrix} = \begin{pmatrix} 1 & 7 \\ 0 & 8 \end{pmatrix}. \]
Observe : les coefficients diagonaux sont $1$, $2$, puis $1$, $4$, puis $1$, $8$... et le coefficient en haut à droite vaut $1$, $3$, $7$... Une conjecture se dessine !
\end{exemplebox}

\begin{experience}[Démonstration guidée : une formule pour $A^n$ par récurrence]
Soit $A = \begin{pmatrix} 1 & 1 \\ 0 & 2 \end{pmatrix}$. Nous allons démontrer que pour tout entier $n \geq 1$ :
\[ A^n = \begin{pmatrix} 1 & 2^n - 1 \\ 0 & 2^n \end{pmatrix}. \]
\textbf{Étape 1 — Initialisation.} Pour $n = 1$ : $\begin{pmatrix} 1 & 2^1 - 1 \\ 0 & 2^1 \end{pmatrix} = \begin{pmatrix} 1 & 1 \\ 0 & 2 \end{pmatrix} = A$. La propriété est vraie au rang $1$.

\textbf{Étape 2 — Hérédité.} Supposons la propriété vraie à un rang $n \geq 1$, c'est-à-dire $A^n = \begin{pmatrix} 1 & 2^n - 1 \\ 0 & 2^n \end{pmatrix}$. Alors :
\begin{align*}
A^{n+1} &= A^n \times A = \begin{pmatrix} 1 & 2^n - 1 \\ 0 & 2^n \end{pmatrix}\begin{pmatrix} 1 & 1 \\ 0 & 2 \end{pmatrix} \\
&= \begin{pmatrix} 1\times 1 + (2^n-1)\times 0 & 1\times 1 + (2^n - 1)\times 2 \\ 0\times 1 + 2^n \times 0 & 0\times 1 + 2^n \times 2 \end{pmatrix} \\
&= \begin{pmatrix} 1 & 1 + 2^{n+1} - 2 \\ 0 & 2^{n+1} \end{pmatrix} = \begin{pmatrix} 1 & 2^{n+1} - 1 \\ 0 & 2^{n+1} \end{pmatrix}.
\end{align*}
La propriété est donc vraie au rang $n+1$.

\textbf{Étape 3 — Conclusion.} Par le principe de récurrence, la formule est vraie pour tout entier $n \geq 1$. Ce type de raisonnement « conjecture puis récurrence » est un grand classique du Bac TI : sache le reproduire.
\end{experience}

\begin{exercice}[À toi de jouer]
Soit $M = \begin{pmatrix} 2 & 0 \\ 1 & 1 \end{pmatrix}$.
\begin{enumerate}
\item Calculer $M^2$ et $M^3$.
\item Conjecturer une expression de $M^n$ pour tout entier $n \geq 1$, puis la démontrer par récurrence.
\item Soient $P = \begin{pmatrix} 1 & 0 \\ 0 & 3 \end{pmatrix}$ et $Q = \begin{pmatrix} 0 & 1 \\ 1 & 0 \end{pmatrix}$. Calculer $PQ$ et $QP$. Que constates-tu ?
\end{enumerate}
\end{exercice}

\begin{corrige}[Exercice : puissances et non-commutativité]
1. $M^2 = \begin{pmatrix} 2 & 0 \\ 1 & 1 \end{pmatrix}\begin{pmatrix} 2 & 0 \\ 1 & 1 \end{pmatrix} = \begin{pmatrix} 4 & 0 \\ 3 & 1 \end{pmatrix}$, puis $M^3 = M^2 M = \begin{pmatrix} 4 & 0 \\ 3 & 1 \end{pmatrix}\begin{pmatrix} 2 & 0 \\ 1 & 1 \end{pmatrix} = \begin{pmatrix} 8 & 0 \\ 7 & 1 \end{pmatrix}$.

2. On conjecture $M^n = \begin{pmatrix} 2^n & 0 \\ 2^n - 1 & 1 \end{pmatrix}$. \textbf{Initialisation} : pour $n=1$, $\begin{pmatrix} 2 & 0 \\ 1 & 1 \end{pmatrix} = M$ : vrai. \textbf{Hérédité} : si $M^n = \begin{pmatrix} 2^n & 0 \\ 2^n - 1 & 1 \end{pmatrix}$, alors
\[ M^{n+1} = \begin{pmatrix} 2^n & 0 \\ 2^n - 1 & 1 \end{pmatrix}\begin{pmatrix} 2 & 0 \\ 1 & 1 \end{pmatrix} = \begin{pmatrix} 2^{n+1} & 0 \\ 2^{n+1} - 2 + 1 & 1 \end{pmatrix} = \begin{pmatrix} 2^{n+1} & 0 \\ 2^{n+1} - 1 & 1 \end{pmatrix}. \]
La propriété est héréditaire, donc vraie pour tout $n \geq 1$ par récurrence.

3. $PQ = \begin{pmatrix} 0 & 1 \\ 3 & 0 \end{pmatrix}$ et $QP = \begin{pmatrix} 0 & 3 \\ 1 & 0 \end{pmatrix}$. On constate que $PQ \neq QP$ : nouvelle confirmation que le produit matriciel n'est pas commutatif.
\end{corrige}

\begin{savaistu}[Les matrices sont partout autour de toi]
Les matrices ne sont pas qu'un jeu de l'esprit ! Une \textbf{image numérique} (photo prise avec ton téléphone) est une gigantesque matrice de nombres : chaque coefficient code la couleur d'un pixel, et les filtres (flou, netteté) sont des opérations matricielles. Les \textbf{graphismes 3D} des jeux vidéo reposent sur des produits de matrices qui décrivent rotations, zooms et déplacements des personnages à chaque image. Plus près de nous, les systèmes de \textbf{transfert d'argent mobile} (MTN Mobile Money, Orange Money) utilisent des calculs matriciels dans leurs algorithmes de cryptage pour sécuriser tes transactions : c'est grâce à des matrices difficiles à « inverser » sans la bonne clé que ton argent voyage en sécurité entre Douala et Maroua. La matrice que tu manipules aujourd'hui sur le papier est la même qui protège le portefeuille de Maman Ngo Bell !
\end{savaistu}

Dans la prochaine section, nous apprendrons à mesurer une matrice avec son \cle{déterminant} et à la « remonter » grâce à son \cle{inverse} : c'est la clé qui permettra à l'opérateur de Yaoundé de retrouver les montants initiaux de ses transactions.

\coursec{Déterminant et inverse d'une matrice 2×2}

Dans la section précédente, nous avons appris à additionner des matrices, à les multiplier par un réel et surtout à calculer le \cle{produit} de deux matrices $2\times 2$. Nous avons constaté que ce produit n'est pas commutatif et qu'il existe une matrice \emph{identité} $I_2 = \begin{pmatrix} 1 & 0 \\ 0 & 1 \end{pmatrix}$ qui joue le rôle du nombre $1$. Une question surgit alors naturellement : pour les nombres réels, tout nombre non nul $a$ possède un \emph{inverse} $\frac{1}{a}$ tel que $a \times \frac{1}{a} = 1$. Est-ce pareil pour les matrices ? Une matrice $A$ possède-t-elle une matrice $A^{-1}$ telle que $A \times A^{-1} = I_2$ ? Et comment la calculer ?

La réponse passe par un nombre remarquable attaché à chaque matrice : son \cle{déterminant}. C'est l'objet de cette section, sans doute la plus importante de la leçon, car elle est la clé de la résolution des systèmes d'équations (section suivante) et de l'étude des transformations du plan.

\courssub{Le déterminant d'une matrice 2×2}

\textbf{Une intuition pour commencer.}\par
Considère la matrice $A = \begin{pmatrix} a & b \\ c & d \end{pmatrix}$. Ses deux colonnes peuvent être vues comme les coordonnées de deux vecteurs du plan : $\vec{u} = \begin{pmatrix} a \\ c \end{pmatrix}$ et $\vec{v} = \begin{pmatrix} b \\ d \end{pmatrix}$. Ces deux vecteurs déterminent un parallélogramme. On peut montrer (et nous le visualiserons plus bas) que l'\emph{aire algébrique} de ce parallélogramme vaut exactement $ad - bc$. Ce nombre, fabriqué à partir des quatre coefficients de $A$, porte un nom : le déterminant de $A$.

\begin{definition}[Déterminant d'une matrice $2\times 2$]
Soit $A = \begin{pmatrix} a & b \\ c & d \end{pmatrix}$ une matrice carrée d'ordre $2$. On appelle \cle{déterminant} de $A$ le réel noté $\det(A)$ ou $|A|$, défini par :
\[ \det(A) = |A| = \begin{vmatrix} a & b \\ c & d \end{vmatrix} = ad - bc. \]
On retient : \emph{produit des termes de la diagonale principale} ($a$ et $d$) \emph{moins produit des termes de l'autre diagonale} ($b$ et $c$).
\end{definition}

\begin{attention}[Déterminant ou valeur absolue ?]
La notation $|A|$ avec des barres verticales est traditionnelle, mais le déterminant n'est \textbf{pas} une valeur absolue : il peut être \emph{négatif} ! Par exemple, pour $A = \begin{pmatrix} 1 & 2 \\ 3 & 4 \end{pmatrix}$, on a $\det(A) = 1\times 4 - 2\times 3 = 4 - 6 = -2$. Le signe du déterminant a un sens géométrique (il indique si la transformation conserve ou renverse l'orientation du plan), et sa valeur absolue donne le facteur d'aire.
\end{attention}

\begin{exemplebox}[Premiers calculs de déterminants]
\begin{enumerate}
\item $A = \begin{pmatrix} 3 & 1 \\ 5 & 2 \end{pmatrix}$ : $\det(A) = 3\times 2 - 1\times 5 = 6 - 5 = 1$.
\item $B = \begin{pmatrix} 2 & -3 \\ 4 & 1 \end{pmatrix}$ : $\det(B) = 2\times 1 - (-3)\times 4 = 2 + 12 = 14$.
\item $C = \begin{pmatrix} 2 & 4 \\ 1 & 2 \end{pmatrix}$ : $\det(C) = 2\times 2 - 4\times 1 = 4 - 4 = 0$. Ce cas, déterminant nul, sera crucial dans la suite.
\end{enumerate}
\end{exemplebox}

\textbf{Interprétation géométrique : le facteur d'aire.}\par
Nous verrons en section 4 qu'à toute matrice $A$ est associée une transformation du plan (une \emph{application linéaire}). Le carré unité, construit sur les vecteurs $\vec{i}$ et $\vec{j}$ de la base, est transformé en un parallélogramme dont l'aire vaut exactement $|\det(A)|$. Le déterminant mesure donc de combien la transformation \emph{gonfle} ou \emph{rétrécit} les aires.

\begin{popfigure}
\begin{center}
\begin{tikzpicture}[scale=1.6,>=Stealth]
% Carré unité
\fill[popblueL,opacity=0.6] (0,0) -- (1,0) -- (1,1) -- (0,1) -- cycle;
\draw[popblue,thick] (0,0) -- (1,0) -- (1,1) -- (0,1) -- cycle;
\draw[->,popdark,thick] (0,0) -- (1,0) node[midway,below] {$\vec{i}$};
\draw[->,popdark,thick] (0,0) -- (0,1) node[midway,left] {$\vec{j}$};
\node[popblue] at (0.5,0.5) {aire $=1$};
\node[below left] at (0,0) {$O$};
% Flèche de transformation
\draw[->,poporange,very thick] (1.6,0.5) -- (2.6,0.5) node[midway,above, align=center] {$A=\begin{pmatrix} 2 & 1 \\ 0.5 & 1 \end{pmatrix}$};
% Parallélogramme image : u=(2,0.5), v=(1,1)
\begin{scope}[xshift=3.2cm]
\fill[poporangeL,opacity=0.6] (0,0) -- (2,0.5) -- (3,1.5) -- (1,1) -- cycle;
\draw[poporange,thick] (0,0) -- (2,0.5) -- (3,1.5) -- (1,1) -- cycle;
\draw[->,popdark,thick] (0,0) -- (2,0.5) node[midway,below right] {$A\vec{i}$};
\draw[->,popdark,thick] (0,0) -- (1,1) node[midway,above left] {$A\vec{j}$};
\node[poporange] at (1.5,0.75) {aire $=|\det(A)|$};
\node[below left] at (0,0) {$O$};
\end{scope}
\end{tikzpicture}
\end{center}
\legende{Le carré unité (aire $1$) est transformé en un parallélogramme d'aire $|\det(A)| = |2\times 1 - 1\times 0{,}5| = 1{,}5$.}
\end{popfigure}

\courssub{Le théorème fondamental : inversibilité et déterminant}

Voici le résultat central de cette section. Il dit que le déterminant est le \emph{test} qui décide si une matrice possède un inverse, exactement comme « être non nul » est le test qui décide si un réel possède un inverse.

\begin{propriete}[Théorème : critère d'inversibilité]
Une matrice carrée $A$ d'ordre $2$ est \cle{inversible} si et seulement si $\det(A) \neq 0$. Dans ce cas, il existe une unique matrice notée $A^{-1}$ telle que :
\[ A \times A^{-1} = A^{-1} \times A = I_2. \]
\emph{Cette démonstration est formatrice et exigible : elle construit explicitement l'inverse.}
\end{propriete}

\begin{experience}[Démonstration guidée du théorème]
Posons $A = \begin{pmatrix} a & b \\ c & d \end{pmatrix}$ et notons $\Delta = ad - bc$ son déterminant.

\textbf{Étape 1 : une matrice candidate.} Construisons, à partir des coefficients de $A$, la matrice $B = \begin{pmatrix} d & -b \\ -c & a \end{pmatrix}$ (on a échangé $a$ et $d$, et changé les signes de $b$ et $c$). Calculons le produit $A \times B$ :
\[ A \times B = \begin{pmatrix} a & b \\ c & d \end{pmatrix} \begin{pmatrix} d & -b \\ -c & a \end{pmatrix} = \begin{pmatrix} ad - bc & -ab + ba \\ cd - dc & -cb + da \end{pmatrix} = \begin{pmatrix} \Delta & 0 \\ 0 & \Delta \end{pmatrix} = \Delta \, I_2. \]
Un calcul identique montre que $B \times A = \Delta \, I_2$.

\textbf{Étape 2 : si $\Delta \neq 0$, l'inverse existe.} On peut diviser par le réel non nul $\Delta$. Posons $A^{-1} = \dfrac{1}{\Delta} B$. Alors :
\[ A \times A^{-1} = A \times \frac{1}{\Delta} B = \frac{1}{\Delta}(A \times B) = \frac{1}{\Delta} \times \Delta\, I_2 = I_2, \]
et de même $A^{-1} \times A = I_2$. La matrice $A$ est donc inversible, et nous avons \emph{construit} son inverse.

\textbf{Étape 3 : si $\Delta = 0$, l'inverse n'existe pas.} Supposons $\Delta = ad - bc = 0$ et raisonnons par l'absurde : supposons qu'il existe une matrice $M$ telle que $A \times M = I_2$. Or le calcul de l'étape 1 donne $B \times A = \Delta\, I_2 = 0_2$ (matrice nulle). Multiplions cette égalité à droite par $M$ :
\[ (B \times A) \times M = 0_2 \times M = 0_2. \]
Mais par associativité, $(B \times A) \times M = B \times (A \times M) = B \times I_2 = B$. Donc $B = 0_2$, c'est-à-dire $a = b = c = d = 0$ : $A$ serait la matrice nulle, qui ne peut évidemment pas vérifier $A \times M = I_2$ (car $0_2 \times M = 0_2 \neq I_2$). Contradiction : $A$ n'est pas inversible. $\blacksquare$
\end{experience}

De cette démonstration, nous tirons immédiatement la formule explicite de l'inverse, à connaître par cœur :

\begin{propriete}[Formule de l'inverse d'une matrice $2\times 2$]
Si $A = \begin{pmatrix} a & b \\ c & d \end{pmatrix}$ avec $\det(A) = ad - bc \neq 0$, alors :
\[ A^{-1} = \frac{1}{\det(A)} \begin{pmatrix} d & -b \\ -c & a \end{pmatrix} = \frac{1}{ad-bc}\begin{pmatrix} d & -b \\ -c & a \end{pmatrix}. \]
\end{propriete}

\begin{popfigure}
\begin{center}
\begin{tikzpicture}[scale=1.1]
\node[draw=popdark,thick,rounded corners,fill=popcream,inner sep=8pt, align=center] (A) at (0,0) {$A=\begin{pmatrix} \textcolor{popblue}{a} & \textcolor{poporange}{b} \\ \textcolor{poporange}{c} & \textcolor{popblue}{d} \end{pmatrix}$};
\draw[->,popdark,very thick] (2.2,0) -- (3.4,0);
\node[draw=popdark,thick,rounded corners,fill=popcream,inner sep=8pt, align=center] (B) at (6.2,0) {$A^{-1}=\dfrac{1}{\textcolor{popgreen}{ad-bc}}\begin{pmatrix} \textcolor{popblue}{d} & \textcolor{poporange}{-b} \\ \textcolor{poporange}{-c} & \textcolor{popblue}{a} \end{pmatrix}$};
% Annotations
\draw[<->,popblue,thick] (-0.55,0.55) to[bend left=40] node[above,popblue,font=\small]{on échange} (0.55,0.55);
\draw[<->,popblue,thick] (5.6,0.55) to[bend left=40] (6.75,0.55);
\node[poporange,font=\small] at (6.2,-1.1) {on change les signes de $b$ et $c$};
\node[popgreen,font=\small] at (6.2,1.15) {on divise tout par le déterminant};
\end{tikzpicture}
\end{center}
\legende{Schéma-mémo de la formule de l'inverse : échanger $a$ et $d$, changer les signes de $b$ et $c$, diviser par $\det(A)$.}
\end{popfigure}

\begin{methode}[Inverser une matrice $2\times 2$ en trois étapes]
Pour inverser $A = \begin{pmatrix} a & b \\ c & d \end{pmatrix}$ :
\begin{enumerate}
\item \textbf{Calculer} le déterminant $\det(A) = ad - bc$. S'il est nul, \emph{stop} : la matrice n'est pas inversible, on le justifie et on s'arrête là.
\item \textbf{Écrire} la matrice $\begin{pmatrix} d & -b \\ -c & a \end{pmatrix}$ : échange de $a$ et $d$, changement de signe de $b$ et $c$.
\item \textbf{Diviser} chaque coefficient par $\det(A)$, puis \textbf{vérifier} en calculant le produit $A \times A^{-1}$ : on doit obtenir $I_2$.
\end{enumerate}
L'étape de vérification n'est pas optionnelle au Baccalauréat : elle sécurise ta copie et te fait gagner des points de rigueur.
\end{methode}

\begin{exemplebox}[Inverse de $A = \begin{pmatrix} 3 & 1 \\ 5 & 2 \end{pmatrix}$, avec vérification]
\textbf{Étape 1.} $\det(A) = 3\times 2 - 1\times 5 = 6 - 5 = 1 \neq 0$ : $A$ est inversible.

\textbf{Étape 2.} On échange $3$ et $2$, on change les signes de $1$ et $5$ : $\begin{pmatrix} 2 & -1 \\ -5 & 3 \end{pmatrix}$.

\textbf{Étape 3.} On divise par $\det(A) = 1$ : $A^{-1} = \begin{pmatrix} 2 & -1 \\ -5 & 3 \end{pmatrix}$.

\textbf{Vérification.} Calculons $A \times A^{-1}$ :
\[ \begin{pmatrix} 3 & 1 \\ 5 & 2 \end{pmatrix} \begin{pmatrix} 2 & -1 \\ -5 & 3 \end{pmatrix} = \begin{pmatrix} 3\times 2 + 1\times(-5) & 3\times(-1) + 1\times 3 \\ 5\times 2 + 2\times(-5) & 5\times(-1) + 2\times 3 \end{pmatrix} = \begin{pmatrix} 1 & 0 \\ 0 & 1 \end{pmatrix} = I_2. \]
Le produit dans l'autre sens donne aussi $I_2$ (vérifie-le !). La formule fonctionne.
\end{exemplebox}

\begin{attention}[Les deux erreurs classiques]
\begin{itemize}
\item \textbf{Oublier de diviser par le déterminant.} Dans l'exemple précédent, $\det(A)=1$, donc la division est invisible... et beaucoup d'élèves l'oublient ensuite quand $\det(A) \neq 1$ ! Pour $B = \begin{pmatrix} 2 & -3 \\ 4 & 1 \end{pmatrix}$, $\det(B) = 14$, donc $B^{-1} = \dfrac{1}{14}\begin{pmatrix} 1 & 3 \\ -4 & 2 \end{pmatrix}$, et \emph{non} $\begin{pmatrix} 1 & 3 \\ -4 & 2 \end{pmatrix}$.
\item \textbf{Se tromper dans les signes.} Ce sont $b$ et $c$ (hors diagonale) qui changent de signe ; $a$ et $d$ gardent leur signe mais \emph{échangent leurs places}. Ne change pas les signes de $a$ et $d$ !
\end{itemize}
Dans les deux cas, la vérification $A \times A^{-1} = I_2$ détecte immédiatement l'erreur : fais-la systématiquement.
\end{attention}

\courssub{Déterminant d'un produit et déterminant de l'inverse}

\begin{propriete}[Déterminant d'un produit (admise)]
Pour toutes matrices carrées $A$ et $B$ d'ordre $2$ :
\[ \det(A \times B) = \det(A) \times \det(B). \]
\emph{Cette propriété est admise ; le calcul complet, fastidieux mais sans difficulté, consiste à développer les deux membres.}
\end{propriete}

\begin{exemplebox}[Vérification sur un exemple]
Prenons $A = \begin{pmatrix} 3 & 1 \\ 5 & 2 \end{pmatrix}$ et $B = \begin{pmatrix} 2 & -3 \\ 4 & 1 \end{pmatrix}$. On a $\det(A) = 1$ et $\det(B) = 14$.
\[ A \times B = \begin{pmatrix} 3\times 2 + 1\times 4 & 3\times(-3) + 1\times 1 \\ 5\times 2 + 2\times 4 & 5\times(-3) + 2\times 1 \end{pmatrix} = \begin{pmatrix} 10 & -8 \\ 18 & -13 \end{pmatrix}. \]
Alors $\det(A\times B) = 10\times(-13) - (-8)\times 18 = -130 + 144 = 14$, et effectivement $\det(A)\times\det(B) = 1 \times 14 = 14$. La propriété est confirmée.
\end{exemplebox}

\begin{propriete}[Déterminant de l'inverse]
Si $A$ est inversible, alors $\det(A^{-1}) = \dfrac{1}{\det(A)}$.
\end{propriete}

\begin{experience}[Preuve en trois lignes]
On sait que $A \times A^{-1} = I_2$. Prenons le déterminant des deux membres et appliquons la propriété du produit :
\[ \det(A) \times \det(A^{-1}) = \det(I_2) = 1\times 1 - 0\times 0 = 1. \]
Comme $\det(A) \neq 0$, on peut diviser : $\det(A^{-1}) = \dfrac{1}{\det(A)}$. $\blacksquare$
\end{experience}

\courssub{Matrices non inversibles : comprendre le cas det(A) = 0}

Que se passe-t-il géométriquement quand $\det(A) = 0$ ? Reprenons l'interprétation en termes d'aire : le parallélogramme construit sur les colonnes $\vec{u} = \begin{pmatrix} a \\ c \end{pmatrix}$ et $\vec{v} = \begin{pmatrix} b \\ d \end{pmatrix}$ a une aire nulle, ce qui signifie que ce parallélogramme est \emph{aplati} : les deux vecteurs colonnes sont \cle{colinéaires}. La transformation associée écrase tout le plan sur une droite (ou sur un point) : il est alors impossible de « revenir en arrière », car des points distincts ont la même image. D'où l'absence d'inverse.

\begin{exemplebox}[Des matrices non inversibles]
\begin{itemize}
\item $C = \begin{pmatrix} 2 & 4 \\ 1 & 2 \end{pmatrix}$ : $\det(C) = 4 - 4 = 0$. Les colonnes $\begin{pmatrix} 2 \\ 1 \end{pmatrix}$ et $\begin{pmatrix} 4 \\ 2 \end{pmatrix}$ sont colinéaires (la seconde est le double de la première).
\item $D = \begin{pmatrix} 3 & 6 \\ 0 & 0 \end{pmatrix}$ : $\det(D) = 0$. La seconde colonne est le double de la première.
\item La matrice nulle $0_2$ : évidemment non inversible.
\end{itemize}
À retenir : $\det(A) = 0 \iff$ les colonnes de $A$ sont colinéaires $\iff$ $A$ n'est pas inversible.
\end{exemplebox}

\begin{savaistu}[Pourquoi « déterminant » ?]
C'est le mathématicien allemand Gottfried Wilhelm Leibniz qui, vers 1693, introduit ces quantités pour \emph{déterminer} si un système d'équations possède une solution unique — d'où le nom ! Nous retrouverons exactement cette idée dans la section suivante : un système linéaire $2\times 2$ a une solution unique précisément quand le déterminant de sa matrice est non nul. Les déterminants sont aussi au cœur des calculs d'images satellites et de cartographie, utilisés par exemple pour estimer les surfaces des parcelles agricoles dans la plaine des Mbo à partir de coordonnées GPS.
\end{savaistu}

\begin{exoresolu}[Inversibilité et calcul d'inverse]
On considère les matrices $M = \begin{pmatrix} 4 & -2 \\ 3 & 1 \end{pmatrix}$ et $N = \begin{pmatrix} 6 & -3 \\ -4 & 2 \end{pmatrix}$.
\begin{enumerate}
\item Chacune de ces matrices est-elle inversible ? Justifier.
\item Calculer l'inverse de celle qui l'est et vérifier le résultat.
\item Calculer $\det(M^{-1})$ sans calculer à nouveau un déterminant $2\times 2$.
\end{enumerate}
\tcblower
\textbf{1.} $\det(M) = 4\times 1 - (-2)\times 3 = 4 + 6 = 10 \neq 0$ : $M$ est inversible. $\det(N) = 6\times 2 - (-3)\times(-4) = 12 - 12 = 0$ : $N$ n'est \textbf{pas} inversible (ses colonnes sont colinéaires : $\begin{pmatrix} -3 \\ 2 \end{pmatrix} = -\frac{1}{2}\begin{pmatrix} 6 \\ -4 \end{pmatrix}$).

\textbf{2.} $M^{-1} = \dfrac{1}{10}\begin{pmatrix} 1 & 2 \\ -3 & 4 \end{pmatrix}$. Vérification :
\[ M \times M^{-1} = \frac{1}{10}\begin{pmatrix} 4 & -2 \\ 3 & 1 \end{pmatrix}\begin{pmatrix} 1 & 2 \\ -3 & 4 \end{pmatrix} = \frac{1}{10}\begin{pmatrix} 4+6 & 8-8 \\ 3-3 & 6+4 \end{pmatrix} = \frac{1}{10}\begin{pmatrix} 10 & 0 \\ 0 & 10 \end{pmatrix} = I_2. \]

\textbf{3.} $\det(M^{-1}) = \dfrac{1}{\det(M)} = \dfrac{1}{10} = 0{,}1$.
\end{exoresolu}

\begin{exercice}[À toi de jouer]
Soit $A = \begin{pmatrix} 5 & 3 \\ 3 & 2 \end{pmatrix}$ et $B = \begin{pmatrix} 1 & 2 \\ 2 & 4 \end{pmatrix}$.
\begin{enumerate}
\item Montrer que $A$ est inversible et calculer $A^{-1}$. Vérifier que $A^{-1} \times A = I_2$.
\item La matrice $B$ est-elle inversible ? Interpréter géométriquement.
\item Calculer $\det(A^2)$ de deux façons différentes.
\end{enumerate}
\end{exercice}

\begin{corrige}[Exercice]
\textbf{1.} $\det(A) = 5\times 2 - 3\times 3 = 10 - 9 = 1 \neq 0$, donc $A$ est inversible et $A^{-1} = \begin{pmatrix} 2 & -3 \\ -3 & 5 \end{pmatrix}$. Vérification : $A^{-1}A = \begin{pmatrix} 10-9 & 15-15 \\ -15+15 & -9+10 \end{pmatrix} = I_2$.

\textbf{2.} $\det(B) = 4 - 4 = 0$ : $B$ n'est pas inversible. Ses colonnes sont colinéaires (la seconde est le double de la première) : la transformation associée écrase le plan sur une droite.

\textbf{3.} Première façon : $\det(A^2) = \det(A)\times\det(A) = 1$. Deuxième façon : $A^2 = \begin{pmatrix} 34 & 21 \\ 21 & 13 \end{pmatrix}$, et $\det(A^2) = 34\times 13 - 21\times 21 = 442 - 441 = 1$. Les deux résultats concordent.
\end{corrige}

\begin{aretenir}[L'essentiel de la section]
\begin{itemize}
\item $\det\begin{pmatrix} a & b \\ c & d \end{pmatrix} = ad - bc$ ; géométriquement, $|\det(A)|$ est le facteur d'aire de la transformation associée.
\item $A$ est inversible $\iff \det(A) \neq 0 \iff$ les colonnes de $A$ ne sont pas colinéaires.
\item Si $\det(A) \neq 0$ : $A^{-1} = \dfrac{1}{ad-bc}\begin{pmatrix} d & -b \\ -c & a \end{pmatrix}$ — échanger $a,d$ ; changer les signes de $b,c$ ; diviser par le déterminant.
\item Toujours vérifier : $A \times A^{-1} = A^{-1} \times A = I_2$.
\item $\det(A\times B) = \det(A)\det(B)$ et $\det(A^{-1}) = \dfrac{1}{\det(A)}$.
\end{itemize}
\end{aretenir}

Nous sommes maintenant armés pour l'une des applications les plus concrètes des matrices : la résolution des systèmes linéaires de deux équations à deux inconnues, que nous aborderons dans la section suivante grâce à l'écriture matricielle $AX = B$.

\coursec{Matrices et systèmes linéaires 2×2}

Dans la section précédente, tu as appris à calculer le déterminant d'une matrice $2\times 2$ et, lorsqu'il est non nul, son inverse. Tu t'es peut-être demandé : à quoi cela sert-il concrètement ? La réponse est magnifique : l'inversion de matrices est un outil de résolution de systèmes d'équations. Un système de deux équations à deux inconnues — que tu résous depuis la classe de Seconde par substitution ou par combinaison — peut s'écrire comme une seule équation matricielle $AX = B$, et se résoudre d'un seul coup par la formule $X = A^{-1}B$. C'est exactement ce que fait un ordinateur (ou ta calculatrice) lorsqu'on lui confie un système, sauf qu'il le fait avec des matrices de grande taille. Dans cette section, tu vas découvrir ce pont fondamental entre l'algèbre des matrices et la géométrie des droites du plan.

\courssub{Écriture matricielle d'un système linéaire 2×2}

Partons d'une situation familière. Au marché Mokolo de Yaoundé, une vendeuse annonce : « 3 kg de tomates et 1 kg d'oignons coûtent 11 000 FCFA ; 5 kg de tomates et 2 kg d'oignons coûtent 19 000 FCFA. » Si $x$ désigne le prix du kilogramme de tomates (en milliers de francs) et $y$ celui des oignons, on obtient le système :
\[
\begin{cases} 3x + y = 11 \\ 5x + 2y = 19 \end{cases}
\]
Observons attentivement le membre de gauche. Les expressions $3x + y$ et $5x + 2y$ ressemblent étrangement aux coefficients d'un produit de matrices. En effet, rappelons la règle du produit matriciel :
\[
\begin{pmatrix} 3 & 1 \\ 5 & 2 \end{pmatrix}\begin{pmatrix} x \\ y \end{pmatrix} = \begin{pmatrix} 3x + 1y \\ 5x + 2y \end{pmatrix}.
\]
Le système tout entier s'écrit donc comme une égalité entre deux matrices colonnes, c'est-à-dire comme une unique équation matricielle.

\begin{definition}[Écriture matricielle d'un système]
Tout système linéaire de deux équations à deux inconnues
\[
\begin{cases} ax + by = e \\ cx + dy = f \end{cases}
\]
s'écrit de manière équivalente sous la forme matricielle
\[
AX = B, \quad \text{avec} \quad A = \begin{pmatrix} a & b \\ c & d \end{pmatrix}, \quad X = \begin{pmatrix} x \\ y \end{pmatrix}, \quad B = \begin{pmatrix} e \\ f \end{pmatrix}.
\]
La matrice $A$ est appelée la \cle{matrice des coefficients} du système, $X$ la \cle{colonne des inconnues} et $B$ le \cle{second membre}.
\end{definition}

\begin{attention}[Bien respecter l'ordre des coefficients]
La première ligne de $A$ contient les coefficients de la première équation, la deuxième ligne ceux de la deuxième équation ; la première colonne correspond à $x$, la deuxième à $y$. Avant d'écrire $A$, vérifie que chaque équation est bien rangée dans l'ordre $ax + by = e$. Par exemple, l'équation $2y - 5x = 3$ doit d'abord être réécrite $-5x + 2y = 3$.
\end{attention}

\courssub{Théorème d'existence et d'unicité : le rôle du déterminant}

Voici l'idée directrice. Pour résoudre l'équation numérique $ax = b$ (avec $a \neq 0$), tu multiplies des deux côtés par l'inverse de $a$ : $x = \dfrac{b}{a}$. Pour l'équation matricielle $AX = B$, on ne peut pas « diviser » par une matrice, mais on peut multiplier à gauche par $A^{-1}$, \emph{à condition que $A$ soit inversible}, c'est-à-dire que $\det(A) \neq 0$. C'est précisément le contenu du théorème fondamental de cette section.

\begin{propriete}[Théorème : existence et unicité de la solution]
Soit le système $AX = B$, où $A = \begin{pmatrix} a & b \\ c & d \end{pmatrix}$ et $B = \begin{pmatrix} e \\ f \end{pmatrix}$.
Si $\det(A) = ad - bc \neq 0$, alors le système admet une \cle{unique solution}, donnée par
\[
X = A^{-1}B.
\]
\end{propriete}

\begin{experience}[Démonstration du théorème (exigible au Bac TI)]
Cette démonstration est courte, élégante, et tu dois savoir la refaire. Elle procède en deux temps : on montre que si une solution existe, elle est forcément $A^{-1}B$ (unicité), puis que $A^{-1}B$ est bien solution (existence).

\textbf{Étape 1 : unicité.} Supposons que $X$ soit une solution, c'est-à-dire $AX = B$. Comme $\det(A) \neq 0$, la matrice $A$ est inversible. Multiplions les deux membres de l'égalité \emph{à gauche} par $A^{-1}$ :
\[
A^{-1}(AX) = A^{-1}B.
\]
Par associativité du produit matriciel, $A^{-1}(AX) = (A^{-1}A)X = I_2 X = X$, où $I_2 = \begin{pmatrix} 1 & 0 \\ 0 & 1 \end{pmatrix}$ est la matrice identité. Donc
\[
X = A^{-1}B.
\]
Toute solution est donc nécessairement égale à $A^{-1}B$ : il y a au plus une solution.

\textbf{Étape 2 : existence.} Vérifions que $X_0 = A^{-1}B$ convient effectivement :
\[
A X_0 = A(A^{-1}B) = (AA^{-1})B = I_2\, B = B.
\]
Donc $X_0$ est bien solution du système.

\textbf{Conclusion.} Le système $AX = B$ admet une solution et une seule : $X = A^{-1}B$. $\blacksquare$
\end{experience}

\begin{aretenir}[La formule à retenir]
\[
\boxed{\;AX = B \ \text{et}\ \det(A) \neq 0 \quad \Longleftrightarrow \quad X = A^{-1}B\;}
\]
avec, pour $A = \begin{pmatrix} a & b \\ c & d \end{pmatrix}$, la formule de l'inverse vue à la section précédente :
\[
A^{-1} = \dfrac{1}{ad - bc}\begin{pmatrix} d & -b \\ -c & a \end{pmatrix}.
\]
\end{aretenir}

\courssub{Méthode complète de résolution par la matrice inverse}

\begin{methode}[Résoudre un système 2×2 par la matrice inverse : quatre étapes]
\begin{enumerate}
\item \textbf{Écrire le système sous forme matricielle} $AX = B$, en rangeant chaque équation dans l'ordre $ax + by = e$.
\item \textbf{Calculer le déterminant} $\det(A) = ad - bc$.
\begin{itemize}
\item Si $\det(A) = 0$ : la méthode ne s'applique pas (voir la discussion plus bas).
\item Si $\det(A) \neq 0$ : continuer.
\end{itemize}
\item \textbf{Calculer l'inverse} $A^{-1} = \dfrac{1}{\det(A)}\begin{pmatrix} d & -b \\ -c & a \end{pmatrix}$ (on échange $a$ et $d$, on change les signes de $b$ et $c$, et on divise tout par le déterminant).
\item \textbf{Calculer le produit} $X = A^{-1}B$, puis \textbf{vérifier} la solution en la reportant dans les deux équations de départ.
\end{enumerate}
\end{methode}

\begin{exemplebox}[Résolution complète du système du marché]
Résolvons le système de notre vendeuse de tomates :
\[
\begin{cases} 3x + y = 11 \\ 5x + 2y = 19 \end{cases}
\]
\textbf{Étape 1.} Forme matricielle : $A = \begin{pmatrix} 3 & 1 \\ 5 & 2 \end{pmatrix}$, $X = \begin{pmatrix} x \\ y \end{pmatrix}$, $B = \begin{pmatrix} 11 \\ 19 \end{pmatrix}$.

\textbf{Étape 2.} Déterminant : $\det(A) = 3 \times 2 - 1 \times 5 = 6 - 5 = 1 \neq 0$. Le système admet une unique solution.

\textbf{Étape 3.} Inverse :
\[
A^{-1} = \dfrac{1}{1}\begin{pmatrix} 2 & -1 \\ -5 & 3 \end{pmatrix} = \begin{pmatrix} 2 & -1 \\ -5 & 3 \end{pmatrix}.
\]

\textbf{Étape 4.} Produit $X = A^{-1}B$ :
\begin{align*}
X &= \begin{pmatrix} 2 & -1 \\ -5 & 3 \end{pmatrix}\begin{pmatrix} 11 \\ 19 \end{pmatrix} = \begin{pmatrix} 2 \times 11 + (-1) \times 19 \\ (-5) \times 11 + 3 \times 19 \end{pmatrix} = \begin{pmatrix} 22 - 19 \\ -55 + 57 \end{pmatrix} = \begin{pmatrix} 3 \\ 2 \end{pmatrix}.
\end{align*}
La solution est $x = 3$ et $y = 2$ : le kilogramme de tomates coûte 3 000 FCFA et celui d'oignons 2 000 FCFA.

\textbf{Vérification :} $3 \times 3 + 2 = 11$ \checkmark\ et $5 \times 3 + 2 \times 2 = 15 + 4 = 19$ \checkmark.
\end{exemplebox}

\begin{attention}[Erreur fréquente : multiplier à droite au lieu de multiplier à gauche]
Le produit matriciel \textbf{n'est pas commutatif} : en général, $A^{-1}B \neq BA^{-1}$. La solution du système $AX = B$ est $X = A^{-1}B$ : on multiplie $B$ \emph{à gauche} par $A^{-1}$. D'ailleurs, ici $BA^{-1}$ n'a même pas de sens en général, car $B$ est une matrice $2 \times 1$ et $A^{-1}$ une matrice $2 \times 2$ : le produit $BA^{-1}$ n'est pas défini (le nombre de colonnes de $B$, à savoir 1, ne coïncide pas avec le nombre de lignes de $A^{-1}$, à savoir 2). Au Bac, écrire $X = BA^{-1}$ est une double faute : de calculabilité et de raisonnement. Retiens : on « divise par $A$ à gauche », car $A$ est à gauche de $X$ dans $AX = B$.
\end{attention}

\courssub{Interprétation géométrique : intersection de deux droites}

Revenons au plan muni d'un repère. Chaque équation $ax + by = e$ (avec $(a,b) \neq (0,0)$) est l'équation d'une \cle{droite}. Résoudre le système, c'est donc chercher les points d'intersection de deux droites. Or deux droites du plan sont dans l'une des trois situations suivantes : sécantes en un unique point, strictement parallèles, ou confondues. Cela explique géométriquement la discussion algébrique.

\begin{propriete}[Discussion selon le déterminant]
Soit le système $AX = B$ associé aux droites $(D_1) : ax + by = e$ et $(D_2) : cx + dy = f$.
\begin{itemize}
\item Si $\det(A) = ad - bc \neq 0$ : les droites sont \cle{sécantes} ; le système admet une \textbf{unique solution}, les coordonnées du point d'intersection.
\item Si $\det(A) = 0$ : les vecteurs colonnes $\begin{pmatrix} a \\ c \end{pmatrix}$ et $\begin{pmatrix} b \\ d \end{pmatrix}$ sont colinéaires, les droites sont \cle{parallèles}. Deux cas se présentent (résultat admis) :
\begin{itemize}
\item droites \textbf{strictement parallèles} : \textbf{aucune solution} ;
\item droites \textbf{confondues} : \textbf{une infinité de solutions} (tous les points de la droite commune).
\end{itemize}
\end{itemize}
\end{propriete}

\begin{popfigure}
\begin{center}
\begin{tikzpicture}
\begin{axis}[
  width=0.72\linewidth, height=6.2cm,
  axis lines=middle, xlabel={$x$}, ylabel={$y$},
  xmin=-1, xmax=5.5, ymin=-1, ymax=6.5,
  xtick={1,2,3,4,5}, ytick={1,2,3,4,5,6},
  grid=both, grid style={popline!40},
  legend style={at={(0.02,0.98)}, anchor=north west, font=\small}
]
\addplot[popcurve, domain=0:4.2, samples=50] {11 - 3*x};
\addlegendentry{$(D_1) : 3x + y = 11$}
\addplot[popcurve, poporange, domain=0:4.2, samples=50] {(19 - 5*x)/2};
\addlegendentry{$(D_2) : 5x + 2y = 19$}
\addplot[only marks, mark=*, mark size=2.5pt, popink] coordinates {(3,2)};
\node[popink, font=\small, anchor=south west] at (axis cs:3.1,2.1) {$(3\,;\,2)$};
\end{axis}
\end{tikzpicture}
\end{center}
\legende{Cas $\det(A) \neq 0$ : les droites $(D_1)$ et $(D_2)$ sont sécantes en un unique point, dont les coordonnées $(3\,;\,2)$ forment l'unique solution du système.}
\end{popfigure}

\begin{popfigure}
\begin{center}
\begin{tikzpicture}
\begin{axis}[
  width=0.46\linewidth, height=5.6cm,
  axis lines=middle, xlabel={$x$}, ylabel={$y$},
  xmin=-0.5, xmax=4, ymin=-0.5, ymax=5,
  xtick={1,2,3}, ytick={1,2,3,4},
  title={Parallèles distinctes}, title style={font=\small}
]
\addplot[popcurve, domain=0:3.5, samples=50] {3.5 - 0.5*x};
\addplot[popcurve, poporange, domain=0:3.5, samples=50] {2 - 0.5*x};
\end{axis}
\end{tikzpicture}
\hfill
\begin{tikzpicture}
\begin{axis}[
  width=0.46\linewidth, height=5.6cm,
  axis lines=middle, xlabel={$x$}, ylabel={$y$},
  xmin=-0.5, xmax=4, ymin=-0.5, ymax=5,
  xtick={1,2,3}, ytick={1,2,3,4},
  title={Droites confondues}, title style={font=\small}
]
\addplot[popcurve, domain=0:3.5, samples=50] {3 - 0.75*x};
\addplot[popcurve, poporange, dashed, domain=0:3.5, samples=50] {3 - 0.75*x};
\end{axis}
\end{tikzpicture}
\end{center}
\legende{Cas $\det(A) = 0$. À gauche : $x + 2y = 7$ et $x + 2y = 4$, droites strictement parallèles, aucune solution. À droite : $3x + 4y = 12$ et $6x + 8y = 24$ (la deuxième équation est le double de la première), droites confondues, infinité de solutions.}
\end{popfigure}

\begin{exemplebox}[Reconnaître les deux cas où det(A) = 0]
\textbf{Cas 1 : aucune solution.} Considérons $\begin{cases} x + 2y = 4 \\ x + 2y = 7 \end{cases}$. Ici $\det(A) = 1 \times 2 - 2 \times 1 = 0$. Les deux équations sont incompatibles : un même nombre $x + 2y$ ne peut valoir à la fois 4 et 7. Géométriquement, les droites sont strictement parallèles.

\textbf{Cas 2 : infinité de solutions.} Considérons $\begin{cases} 3x + 4y = 12 \\ 6x + 8y = 24 \end{cases}$. Ici $\det(A) = 3 \times 8 - 4 \times 6 = 24 - 24 = 0$. La deuxième équation est exactement le double de la première : elle n'apporte aucune information nouvelle. Tout couple $\left(x\,;\, \dfrac{12 - 3x}{4}\right)$, avec $x \in \mathbb{R}$, est solution : il y en a une infinité.
\end{exemplebox}

\courssub{Application concrète : problème de stock et de mélange}

Les systèmes $2\times 2$ apparaissent naturellement dès qu'une situation mélange deux grandeurs inconnues liées par deux contraintes : prix, stocks, mélanges, effectifs. La méthode matricielle donne alors un procédé de résolution systématique, particulièrement apprécié lorsque les coefficients sont « laids ».

\begin{exoresolu}[Gestion de stock dans une station-service]
Une station-service de Douala vend deux types de carburant : super ($x$, en litres) et gasoil ($y$, en litres). En une journée, la vente totale est de 1450 litres pour une recette de \num{1015000} FCFA, le super valant 750 FCFA le litre et le gasoil 650 FCFA le litre. Déterminer les quantités vendues de chaque carburant en résolvant le système par la matrice inverse.
\tcblower
\textbf{Mise en équations.} Soit $x$ le nombre de litres de super et $y$ celui de gasoil.
\[
\begin{cases} x + y = 1450 & \text{(volume total)} \\ 750x + 650y = \num{1015000} & \text{(recette totale)} \end{cases}
\]
\textbf{Étape 1.} Forme matricielle : $A = \begin{pmatrix} 1 & 1 \\ 750 & 650 \end{pmatrix}$, $B = \begin{pmatrix} 1450 \\ \num{1015000} \end{pmatrix}$.

\textbf{Étape 2.} Déterminant : $\det(A) = 1 \times 650 - 1 \times 750 = -100 \neq 0$. Solution unique.

\textbf{Étape 3.} Inverse :
\[
A^{-1} = \dfrac{1}{-100}\begin{pmatrix} 650 & -1 \\ -750 & 1 \end{pmatrix} = \begin{pmatrix} -6{,}5 & 0{,}01 \\ 7{,}5 & -0{,}01 \end{pmatrix}.
\]

\textbf{Étape 4.} Solution :
\begin{align*}
X &= A^{-1}B = \dfrac{1}{-100}\begin{pmatrix} 650 & -1 \\ -750 & 1 \end{pmatrix}\begin{pmatrix} 1450 \\ \num{1015000} \end{pmatrix}\\
&= \dfrac{1}{-100}\begin{pmatrix} 650 \times 1450 - \num{1015000} \\ -750 \times 1450 + \num{1015000} \end{pmatrix}
= \dfrac{1}{-100}\begin{pmatrix} 942500 - \num{1015000} \\ -\num{1087500} + \num{1015000} \end{pmatrix}\\
&= \dfrac{1}{-100}\begin{pmatrix} -72500 \\ -72500 \end{pmatrix} = \begin{pmatrix} 725 \\ 725 \end{pmatrix}.
\end{align*}
\textbf{Conclusion.} La station a vendu 725 litres de super et 725 litres de gasoil.

\textbf{Vérification :} $725 + 725 = 1450$ \checkmark\ et $750 \times 725 + 650 \times 725 = 1400 \times 725 = \num{1015000}$ \checkmark.
\end{exoresolu}

\begin{savaistu}[La résolution de systèmes par matrice inverse : un classique du Bac TI]
À l'examen, la résolution d'un système $2\times 2$ par la matrice inverse est l'une des questions les plus fréquentes des exercices de calcul matriciel. Le schéma type est le suivant : on te donne une matrice $A$, on te demande de calculer $\det(A)$ puis $A^{-1}$, puis d'\emph{en déduire} la résolution d'un système. Le mot « en déduire » est un signal : l'examinateur attend que tu utilises la formule $X = A^{-1}B$, et non une résolution par substitution qui te ferait perdre des points de méthode. Autre fait amusant : cette idée de « résoudre en inversant » est à la base de méthodes industrielles colossales. Les ingénieurs télécoms (chez MTN ou Orange, par exemple) traitent des signaux qui se mélangent : les séparer revient à résoudre des systèmes linéaires à des milliers d'inconnues, exactement par le même principe matriciel que celui que tu viens d'apprendre, appliqué à des matrices géantes inversées par ordinateur.
\end{savaistu}

\begin{exercice}[À toi de jouer]
Un cybercafé de Bafoussam propose deux formules : connexion simple à 200 FCFA l'heure et connexion avec impression à 500 FCFA l'heure. En une journée, 60 heures de connexion ont été vendues pour une recette de \num{21000} FCFA.
\begin{enumerate}
\item Écrire le système correspondant sous la forme $AX = B$.
\item Calculer $\det(A)$ et justifier l'existence d'une unique solution.
\item Calculer $A^{-1}$ puis résoudre le système par la formule $X = A^{-1}B$.
\end{enumerate}
\end{exercice}

\begin{corrige}[Exercice : cybercafé de Bafoussam]
\textbf{1.} Soit $x$ le nombre d'heures simples et $y$ le nombre d'heures avec impression :
\[
\begin{cases} x + y = 60 \\ 200x + 500y = \num{21000} \end{cases}
\iff AX = B, \quad A = \begin{pmatrix} 1 & 1 \\ 200 & 500 \end{pmatrix},\ B = \begin{pmatrix} 60 \\ \num{21000} \end{pmatrix}.
\]
\textbf{2.} $\det(A) = 500 - 200 = 300 \neq 0$ : $A$ est inversible, le système admet une unique solution.

\textbf{3.} $A^{-1} = \dfrac{1}{300}\begin{pmatrix} 500 & -1 \\ -200 & 1 \end{pmatrix}$, puis
\[
X = A^{-1}B = \dfrac{1}{300}\begin{pmatrix} 500 \times 60 - \num{21000} \\ -200 \times 60 + \num{21000} \end{pmatrix} = \dfrac{1}{300}\begin{pmatrix} 9000 \\ 9000 \end{pmatrix} = \begin{pmatrix} 30 \\ 30 \end{pmatrix}.
\]
Le cybercafé a vendu 30 heures de connexion simple et 30 heures avec impression. Vérification : $200 \times 30 + 500 \times 30 = 6000 + 15000 = \num{21000}$ \checkmark.
\end{corrige}

\begin{aretenir}[L'essentiel de la section]
\begin{itemize}
\item Tout système $\begin{cases} ax + by = e \\ cx + dy = f \end{cases}$ s'écrit $AX = B$ avec $A = \begin{pmatrix} a & b \\ c & d \end{pmatrix}$.
\item Si $\det(A) \neq 0$ : \textbf{solution unique} $X = A^{-1}B$ (on multiplie $B$ \emph{à gauche} par $A^{-1}$).
\item Si $\det(A) = 0$ : \textbf{aucune solution} (droites strictement parallèles) ou \textbf{infinité de solutions} (droites confondues).
\item Géométriquement, résoudre le système revient à chercher l'intersection de deux droites du plan.
\item La méthode en quatre étapes : forme matricielle, déterminant, inverse, produit $A^{-1}B$ — suivie d'une vérification.
\end{itemize}
\end{aretenir}

Cette section t'a montré que les matrices ne sont pas qu'un jeu de calcul : elles \emph{pensent} les systèmes d'équations. Dans la section suivante, tu découvriras que cette même matrice $2\times 2$ cache encore un autre trésor : elle décrit les \emph{applications linéaires du plan}, c'est-à-dire les transformations qui préservent les droites et l'origine. Le lien entre matrices et géométrie ne fait que commencer.

\coursec{Matrice d'une application linéaire du plan}

Dans la section précédente, nous avons vu qu'une matrice $2\times 2$ inversible permet de résoudre un système linéaire de deux équations à deux inconnues : le système $\begin{cases} ax+by=e \\ cx+dy=f \end{cases}$ s'écrit $A\begin{pmatrix} x \\ y \end{pmatrix} = \begin{pmatrix} e \\ f \end{pmatrix}$, et sa solution est $\begin{pmatrix} x \\ y \end{pmatrix} = A^{-1}\begin{pmatrix} e \\ f \end{pmatrix}$. Nous allons maintenant donner une \emph{signification géométrique} profonde aux matrices : chaque matrice $2\times 2$ décrit une \cle{transformation du plan}, et réciproquement, toute transformation « linéaire » du plan (homothétie, symétrie, projection, rotation) peut être codée par une matrice. C'est ce pont entre l'algèbre (les calculs matriciels) et la géométrie (les transformations) qui fait toute la puissance de cette leçon.

\courssub{Rappel : les applications linéaires du plan}

\begin{definition}[Application linéaire du plan]
Soit $f$ une application qui, à tout vecteur $\vec{u}$ du plan, associe un vecteur $f(\vec{u})$ du plan. On dit que $f$ est une \cle{application linéaire} du plan si elle vérifie les deux conditions suivantes, pour tous vecteurs $\vec{u}$, $\vec{v}$ et tout réel $k$ :
\begin{enumerate}
\item \textbf{Additivité :} $f(\vec{u}+\vec{v}) = f(\vec{u}) + f(\vec{v})$ ;
\item \textbf{Homogénéité :} $f(k\vec{u}) = k\,f(\vec{u})$.
\end{enumerate}
Une application linéaire \emph{bijective} est appelée un \cle{automorphisme} du plan.
\end{definition}

\begin{attention}[Conséquences immédiates]
Si $f$ est linéaire, alors $f(\vec{0}) = \vec{0}$ (prendre $k=0$ dans l'homogénéité) et $f(-\vec{u}) = -f(\vec{u})$. Une application qui ne fixe pas l'origine, comme une translation de vecteur non nul, n'est donc \emph{jamais} linéaire. De même, une application du type $f(\vec{u}) = \vec{u} + \vec{a}$ avec $\vec{a}\neq\vec{0}$ n'est pas linéaire : c'est une erreur classique au Bac.
\end{attention}

\begin{exemplebox}[Les transformations usuelles sont linéaires]
Les transformations suivantes du plan sont des applications linéaires :
\begin{itemize}
\item les \textbf{homothéties} de centre $O$ (rapport $k$) ;
\item les \textbf{symétries} par rapport à une droite passant par $O$ (symétries orthogonales ou obliques) et les symétries centrales de centre $O$ ;
\item les \textbf{projections} sur une droite passant par $O$, parallèlement à une autre droite ;
\item les \textbf{rotations} de centre $O$.
\end{itemize}
\textbf{Vérification sur un cas :} soit $h$ l'homothétie de centre $O$ et de rapport $k$. Pour tous vecteurs $\vec{u}, \vec{v}$ et tout réel $\lambda$ :
\[ h(\vec{u}+\vec{v}) = k(\vec{u}+\vec{v}) = k\vec{u}+k\vec{v} = h(\vec{u})+h(\vec{v}), \qquad h(\lambda\vec{u}) = k\lambda\vec{u} = \lambda\, h(\vec{u}). \]
L'homothétie $h$ est donc bien linéaire. Pour une symétrie axiale, l'additivité se lit géométriquement : l'image d'un parallélogramme est un parallélogramme, donc l'image de la somme est la somme des images.
\end{exemplebox}

\courssub{Théorème fondamental : une application linéaire est déterminée par les images d'une base}

Voici l'idée qui change tout. Pour décrire une application quelconque du plan, il faudrait a priori donner l'image de chacun des vecteurs du plan, c'est-à-dire une infinité de données. Mais pour une application \emph{linéaire}, il suffit de connaître les images de \textbf{deux} vecteurs bien choisis : ceux d'une base. Tout le reste s'en déduit par calcul.

\begin{propriete}[Théorème fondamental — démonstration exigible]
Soit $(\vec{i},\vec{j})$ une base du plan. Une application linéaire $f$ est \textbf{entièrement déterminée} par les images $f(\vec{i})$ et $f(\vec{j})$. Autrement dit : deux applications linéaires qui coïncident sur $\vec{i}$ et sur $\vec{j}$ sont égales partout.
\end{propriete}

\begin{experience}[Démonstration du théorème fondamental]
Soit $f$ une application linéaire et $\vec{u}$ un vecteur quelconque du plan.
\begin{enumerate}
\item Comme $(\vec{i},\vec{j})$ est une base, il existe un \textbf{unique} couple de réels $(x,y)$ tel que
\[ \vec{u} = x\,\vec{i} + y\,\vec{j}. \]
\item Appliquons $f$ et utilisons ses deux propriétés de linéarité :
\begin{align*}
f(\vec{u}) &= f(x\,\vec{i} + y\,\vec{j}) \\
&= f(x\,\vec{i}) + f(y\,\vec{j}) & \text{(additivité)}\\
&= x\,f(\vec{i}) + y\,f(\vec{j}) & \text{(homogénéité, deux fois)}.
\end{align*}
\item Conclusion : dès que l'on connaît $f(\vec{i})$ et $f(\vec{j})$, on peut calculer $f(\vec{u})$ pour \emph{n'importe quel} vecteur $\vec{u}$, par la formule
\[ \boxed{\;f(\vec{u}) = x\,f(\vec{i}) + y\,f(\vec{j})\;} \]
L'application $f$ est donc entièrement déterminée par les deux vecteurs $f(\vec{i})$ et $f(\vec{j})$. $\blacksquare$
\end{enumerate}
\end{experience}

\begin{exemplebox}[Application numérique immédiate]
Soit $f$ l'application linéaire telle que $f(\vec{i}) = 2\vec{i}+\vec{j}$ et $f(\vec{j}) = -\vec{i}+3\vec{j}$. Calculons l'image de $\vec{u} = 2\vec{i}+\vec{j}$ :
\[ f(\vec{u}) = 2\,f(\vec{i}) + f(\vec{j}) = 2(2\vec{i}+\vec{j}) + (-\vec{i}+3\vec{j}) = 3\vec{i}+5\vec{j}. \]
On n'a eu besoin que des images de $\vec{i}$ et $\vec{j}$ !
\end{exemplebox}

\begin{popfigure}
\begin{tikzpicture}[scale=1.15, >=Stealth]
\draw[popline, ->] (-1.6,0) -- (4.6,0);
\draw[popline, ->] (0,-1.2) -- (0,4.2);
\draw[popdark, very thick, ->] (0,0) -- (1,0) node[midway, below] {$\vec{i}$};
\draw[popdark, very thick, ->] (0,0) -- (0,1) node[midway, left] {$\vec{j}$};
\draw[popink, very thick, ->] (0,0) -- (2,1) node[midway, above left] {$f(\vec{i})$};
\draw[poporange, very thick, ->] (0,0) -- (-1,3) node[midway, left] {$f(\vec{j})$};
\draw[popmuted, dashed] (2,1) -- (1,4) -- (-1,3);
\draw[popblue, very thick, ->] (0,0) -- (1,4) node[midway, right] {$f(\vec{i})+f(\vec{j})$};
\fill[popdark] (0,0) circle (1.5pt) node[below left] {$O$};
\fill[popink] (2,1) circle (1.5pt) node[below right] {$(2\,;\,1)$};
\fill[poporange] (-1,3) circle (1.5pt) node[above left] {$(-1\,;\,3)$};
\node[popdark, align=left] at (3.4,3.4) {$M = \begin{pmatrix} 2 & -1 \\ 1 & 3 \end{pmatrix}$};
\node[popmuted, align=left] at (3.4,2.3) {\footnotesize colonne 1 : $f(\vec{i})$};
\node[popmuted, align=left] at (3.4,1.8) {\footnotesize colonne 2 : $f(\vec{j})$};
\end{tikzpicture}
\legende{Les images $f(\vec{i})$ et $f(\vec{j})$ suffisent à décrire $f$ ; leurs coordonnées forment les colonnes de la matrice.}
\end{popfigure}

\courssub{Matrice d'une application linéaire dans une base}

\begin{definition}[Matrice d'une application linéaire]
Soit $f$ une application linéaire du plan et $(\vec{i},\vec{j})$ une base. La \cle{matrice de $f$ dans la base $(\vec{i},\vec{j})$} est la matrice $2\times 2$ dont :
\begin{itemize}
\item la \textbf{première colonne} est formée des coordonnées de $f(\vec{i})$ dans la base $(\vec{i},\vec{j})$ ;
\item la \textbf{deuxième colonne} est formée des coordonnées de $f(\vec{j})$ dans la base $(\vec{i},\vec{j})$.
\end{itemize}
Autrement dit, si $f(\vec{i}) = a\vec{i}+c\vec{j}$ et $f(\vec{j}) = b\vec{i}+d\vec{j}$, alors
\[ \operatorname{Mat}(f) = \begin{pmatrix} a & b \\ c & d \end{pmatrix}. \]
\end{definition}

\begin{attention}[Le piège n°1 : colonnes, pas lignes !]
L'erreur la plus fréquente au Bac consiste à écrire les coordonnées de $f(\vec{i})$ et $f(\vec{j})$ \textbf{en lignes}. Retiens bien : \textbf{les images des vecteurs de base s'écrivent en COLONNES}. Si $f(\vec{i}) = 2\vec{i}+\vec{j}$ et $f(\vec{j}) = -\vec{i}+3\vec{j}$, la matrice est
\[ \begin{pmatrix} 2 & -1 \\ 1 & 3 \end{pmatrix} \quad \text{et non pas} \quad \begin{pmatrix} 2 & 1 \\ -1 & 3 \end{pmatrix}. \]
Moyen mnémotechnique : « les vecteurs sont des colonnes, leurs images restent des colonnes ».
\end{attention}

\begin{methode}[Construire la matrice d'une application linéaire]
Pour déterminer la matrice de $f$ dans la base $(\vec{i},\vec{j})$ :
\begin{enumerate}
\item Calculer (ou lire sur une figure) l'image $f(\vec{i})$ et l'exprimer dans la base : $f(\vec{i}) = a\vec{i}+c\vec{j}$.
\item Calculer de même $f(\vec{j}) = b\vec{i}+d\vec{j}$.
\item Écrire la matrice en plaçant ces coordonnées \textbf{en colonnes} : $M = \begin{pmatrix} a & b \\ c & d \end{pmatrix}$.
\item Vérifier sur un vecteur simple : calculer $f(\vec{i}+\vec{j})$ directement et par le produit matriciel.
\end{enumerate}
\end{methode}

\courssub{Coordonnées de l'image d'un vecteur : le produit matriciel}

\begin{propriete}[Théorème — démonstration exigible]
Soit $f$ une application linéaire de matrice $M = \begin{pmatrix} a & b \\ c & d \end{pmatrix}$ dans une base $(\vec{i},\vec{j})$. Si $\vec{u}$ a pour coordonnées $\begin{pmatrix} x \\ y \end{pmatrix}$, alors $f(\vec{u})$ a pour coordonnées
\[ M\begin{pmatrix} x \\ y \end{pmatrix} = \begin{pmatrix} ax+by \\ cx+dy \end{pmatrix}. \]
Autrement dit : \textbf{les coordonnées de l'image s'obtiennent par le produit de la matrice par le vecteur-colonne des coordonnées}.
\end{propriete}

\begin{experience}[Démonstration par la linéarité]
Soit $\vec{u} = x\vec{i}+y\vec{j}$. Par le théorème fondamental :
\[ f(\vec{u}) = x\,f(\vec{i}) + y\,f(\vec{j}). \]
Or $f(\vec{i}) = a\vec{i}+c\vec{j}$ et $f(\vec{j}) = b\vec{i}+d\vec{j}$ (définition de la matrice). Donc :
\begin{align*}
f(\vec{u}) &= x(a\vec{i}+c\vec{j}) + y(b\vec{i}+d\vec{j}) \\
&= (ax+by)\,\vec{i} + (cx+dy)\,\vec{j}.
\end{align*}
Les coordonnées de $f(\vec{u})$ sont donc $\begin{pmatrix} ax+by \\ cx+dy \end{pmatrix}$, qui est exactement le produit $M\begin{pmatrix} x \\ y \end{pmatrix}$. $\blacksquare$
\end{experience}

\begin{exoresolu}[Image d'un vecteur par produit matriciel]
Soit $f$ l'application linéaire de matrice $M = \begin{pmatrix} 2 & -1 \\ 1 & 3 \end{pmatrix}$ dans la base $(\vec{i},\vec{j})$. Déterminer les coordonnées de l'image du vecteur $\vec{u} = 2\vec{i}+\vec{j}$.
\tcblower
\textbf{Solution.} Les coordonnées de $\vec{u}$ sont $\begin{pmatrix} 2 \\ 1 \end{pmatrix}$. On calcule le produit :
\[ M\begin{pmatrix} 2 \\ 1 \end{pmatrix} = \begin{pmatrix} 2 & -1 \\ 1 & 3 \end{pmatrix}\begin{pmatrix} 2 \\ 1 \end{pmatrix} = \begin{pmatrix} 2\times 2 + (-1)\times 1 \\ 1\times 2 + 3\times 1 \end{pmatrix} = \begin{pmatrix} 3 \\ 5 \end{pmatrix}. \]
Donc $f(\vec{u}) = 3\vec{i}+5\vec{j}$. On retrouve exactement le résultat obtenu plus haut par linéarité : $2f(\vec{i})+f(\vec{j}) = 2(2\vec{i}+\vec{j})+(-\vec{i}+3\vec{j}) = 3\vec{i}+5\vec{j}$. Les deux méthodes coïncident, ce qui conforte le théorème.
\end{exoresolu}

\begin{popfigure}
\begin{tikzpicture}[scale=1.0, >=Stealth]
\draw[popline, ->] (-1.2,0) -- (5.4,0);
\draw[popline, ->] (0,-1.2) -- (0,5.6);
\foreach \x in {1,2,3,4,5} \draw[popline!50] (\x,-0.06) -- (\x,0.06);
\foreach \y in {1,2,3,4,5} \draw[popline!50] (-0.06,\y) -- (0.06,\y);
\draw[popdark, thick, ->] (0,0) -- (1,0) node[midway, below] {$\vec{i}$};
\draw[popdark, thick, ->] (0,0) -- (0,1) node[midway, left] {$\vec{j}$};
\draw[popblue, very thick, ->] (0,0) -- (2,1) node[midway, below right] {$\vec{u}=2\vec{i}+\vec{j}$};
\draw[popmuted, dashed] (2,0) -- (2,1);
\draw[popink, very thick, ->] (0,0) -- (3,5) node[midway, above right] {$f(\vec{u})=3\vec{i}+5\vec{j}$};
\draw[popink!60, thick, ->] (0,0) -- (2,1);
\draw[poporange, thick, ->] (0,0) -- (-1,3) node[midway, left] {$f(\vec{j})$};
\draw[popmuted, dashed] (2,1) -- (4,2);
\draw[popmuted, dashed] (2,1) -- (1,4) -- (3,5);
\draw[popmuted, dashed] (4,2) -- (3,5);
\fill[popdark] (0,0) circle (1.5pt) node[below left] {$O$};
\fill[popblue] (2,1) circle (1.5pt);
\fill[popink] (3,5) circle (1.5pt);
\node[popmuted, align=left] at (4.3,4.3) {\footnotesize $f(\vec{u}) = 2f(\vec{i})+f(\vec{j})$};
\node[popmuted, align=left] at (4.3,3.7) {\footnotesize parallélogramme image};
\end{tikzpicture}
\legende{L'image de $\vec{u}=2\vec{i}+\vec{j}$ se construit géométriquement : le parallélogramme construit sur $\vec{i},\vec{j}$ a pour image le parallélogramme construit sur $f(\vec{i}),f(\vec{j})$.}
\end{popfigure}

\courssub{Matrices des transformations usuelles}

\begin{exemplebox}[Un tableau à connaître par cœur]
Dans une base $(\vec{i},\vec{j})$ orthonormée directe :
\begin{center}
\begin{tabularx}{\linewidth}{|X|c|}
\hline
\rowcolor{popblueL} \textbf{Application linéaire} & \textbf{Matrice} \\
\hline
Identité : $f(\vec{u})=\vec{u}$ & $\begin{pmatrix} 1 & 0 \\ 0 & 1 \end{pmatrix}$ \\
\hline
Homothétie de rapport $k$ & $\begin{pmatrix} k & 0 \\ 0 & k \end{pmatrix}$ \\
\hline
Symétrie d'axe $(O,\vec{i})$ & $\begin{pmatrix} 1 & 0 \\ 0 & -1 \end{pmatrix}$ \\
\hline
Symétrie d'axe $(O,\vec{j})$ & $\begin{pmatrix} -1 & 0 \\ 0 & 1 \end{pmatrix}$ \\
\hline
Symétrie centrale de centre $O$ & $\begin{pmatrix} -1 & 0 \\ 0 & -1 \end{pmatrix}$ \\
\hline
Projection sur $(O,\vec{i})$ parallèlement à $(O,\vec{j})$ & $\begin{pmatrix} 1 & 0 \\ 0 & 0 \end{pmatrix}$ \\
\hline
\end{tabularx}
\end{center}
\textbf{Vérification pour la symétrie d'axe $(O,\vec{i})$ :} l'axe des abscisses est invariant donc $f(\vec{i}) = \vec{i} = 1\vec{i}+0\vec{j}$ (colonne 1 : $\begin{pmatrix} 1 \\ 0 \end{pmatrix}$), et $\vec{j}$ est envoyé sur son opposé : $f(\vec{j}) = -\vec{j} = 0\vec{i}-1\vec{j}$ (colonne 2 : $\begin{pmatrix} 0 \\ -1 \end{pmatrix}$). D'où la matrice annoncée. Pour la projection sur $(O,\vec{i})$ : $f(\vec{i})=\vec{i}$ et $f(\vec{j})=\vec{0}$, d'où les colonnes $\begin{pmatrix} 1 \\ 0 \end{pmatrix}$ et $\begin{pmatrix} 0 \\ 0 \end{pmatrix}$.
\end{exemplebox}

\begin{propriete}[Complément Bac : matrice d'une rotation — admise]
Dans une base orthonormée \textbf{directe} $(\vec{i},\vec{j})$, la rotation de centre $O$ et d'angle $\theta$ est une application linéaire de matrice
\[ R_\theta = \begin{pmatrix} \cos\theta & -\sin\theta \\ \sin\theta & \cos\theta \end{pmatrix}. \]
Ce résultat est admis ; il s'obtient en lisant les coordonnées trigonométriques des images de $\vec{i}$ et $\vec{j}$ sur le cercle trigonométrique : $f(\vec{i}) = \cos\theta\,\vec{i}+\sin\theta\,\vec{j}$ et $f(\vec{j}) = -\sin\theta\,\vec{i}+\cos\theta\,\vec{j}$.
\end{propriete}

\begin{popfigure}
\begin{tikzpicture}[scale=2.1, >=Stealth]
\draw[popline, ->] (-1.5,0) -- (1.6,0);
\draw[popline, ->] (0,-0.6) -- (0,1.5);
\draw[popline!60] (0,0) circle (1);
\draw[popdark, very thick, ->] (0,0) -- (1,0) node[midway, below] {$\vec{i}$};
\draw[popdark, very thick, ->] (0,0) -- (0,1) node[midway, left] {$\vec{j}$};
\draw[popink, very thick, ->] (0,0) -- (50:1) node[midway, above right] {\footnotesize $f(\vec{i})$};
\draw[poporange, very thick, ->] (0,0) -- (140:1) node[midway, above left] {\footnotesize $f(\vec{j})$};
\draw[popink, ->] (0.45,0) arc (0:50:0.45) node[midway, right] {\footnotesize $\theta$};
\draw[poporange, ->] (0,0.55) arc (90:140:0.55) node[midway, above] {\footnotesize $\theta$};
\draw[popmuted, dashed] (50:1) -- ++(0,{-sin(50)});
\fill[popink] (50:1) circle (0.8pt);
\fill[poporange] (140:1) circle (0.8pt);
\node[popink, align=center] at (1.15,-0.45) {\footnotesize $(\cos\theta\,;\,\sin\theta)$};
\node[poporange, align=center] at (-1.25,0.95) {\footnotesize $(-\sin\theta\,;\,\cos\theta)$};
\fill[popdark] (0,0) circle (0.8pt) node[below left] {$O$};
\end{tikzpicture}
\legende{Rotation d'angle $\theta$ : les coordonnées de $f(\vec{i})$ et $f(\vec{j})$ se lisent sur le cercle trigonométrique et forment les colonnes de $R_\theta$.}
\end{popfigure}

\begin{exemplebox}[Un cas particulier utile]
Pour $\theta = \dfrac{\pi}{2}$ : $\cos\theta = 0$ et $\sin\theta = 1$, donc
\[ R_{\pi/2} = \begin{pmatrix} 0 & -1 \\ 1 & 0 \end{pmatrix}. \]
Vérifions : le quart de tour direct envoie $\vec{i}$ sur $\vec{j}$ (colonne $\begin{pmatrix} 0 \\ 1 \end{pmatrix}$) et $\vec{j}$ sur $-\vec{i}$ (colonne $\begin{pmatrix} -1 \\ 0 \end{pmatrix}$). C'est exactement la matrice obtenue.
\end{exemplebox}

\courssub{Lecture inverse : retrouver l'application à partir de sa matrice}

La correspondance fonctionne dans les deux sens : toute matrice $2\times 2$ définit une unique application linéaire.

\begin{methode}[De la matrice à l'application linéaire]
Soit $M = \begin{pmatrix} a & b \\ c & d \end{pmatrix}$ une matrice donnée. Pour décrire l'application linéaire $f$ associée :
\begin{enumerate}
\item Lire les images des vecteurs de base dans les colonnes : $f(\vec{i}) = a\vec{i}+c\vec{j}$ et $f(\vec{j}) = b\vec{i}+d\vec{j}$.
\item Écrire l'expression générale : l'image de $\vec{u}(x\,;\,y)$ est $f(\vec{u}) = (ax+by)\,\vec{i} + (cx+dy)\,\vec{j}$.
\item Si possible, identifier géométriquement $f$ (homothétie, symétrie, projection, rotation...) en comparant avec le tableau des matrices usuelles.
\end{enumerate}
\end{methode}

\begin{exoresolu}[Identifier une transformation à partir de sa matrice]
On donne $M = \begin{pmatrix} 0 & 1 \\ 1 & 0 \end{pmatrix}$. Déterminer l'application linéaire $f$ de matrice $M$ dans une base orthonormée $(\vec{i},\vec{j})$ et l'interpréter géométriquement.
\tcblower
\textbf{Solution.} \textbf{Étape 1 :} lecture des colonnes : $f(\vec{i}) = 0\vec{i}+1\vec{j} = \vec{j}$ et $f(\vec{j}) = 1\vec{i}+0\vec{j} = \vec{i}$. $f$ échange les vecteurs de base.
\textbf{Étape 2 :} pour $\vec{u}(x\,;\,y)$ : $M\begin{pmatrix} x \\ y \end{pmatrix} = \begin{pmatrix} y \\ x \end{pmatrix}$, donc $f(\vec{u}) = y\,\vec{i}+x\,\vec{j}$ : $f$ échange les coordonnées.
\textbf{Étape 3 :} interprétation. Un point $M(x\,;\,y)$ a pour image $M'(y\,;\,x)$. On reconnaît la \textbf{symétrie orthogonale par rapport à la première bissectrice} (la droite d'équation $y=x$). Vérification : les points de cette droite $(x\,;\,x)$ sont invariants, et le milieu de $[MM']$ est sur la droite avec $(MM')$ perpendiculaire à celle-ci.
\end{exoresolu}

\begin{attention}[Conditions d'application à ne pas oublier]
\begin{itemize}
\item La matrice d'une application linéaire \textbf{dépend de la base choisie} : la même transformation peut avoir des matrices différentes dans deux bases différentes. Précise toujours la base dans ta rédaction.
\item La matrice de la rotation $\begin{pmatrix} \cos\theta & -\sin\theta \\ \sin\theta & \cos\theta \end{pmatrix}$ n'est valable que dans une base orthonormée \textbf{directe}. Dans une base orthonormée indirecte, la matrice de la même rotation géométrique n'a pas cette forme (elle s'obtient par changement de base).
\item Ne confonds pas les coordonnées d'un point $M(x\,;\,y)$ et celles d'un vecteur : le produit matriciel s'applique au vecteur-colonne des coordonnées, écrit verticalement.
\end{itemize}
\end{attention}

\begin{aretenir}[L'essentiel de la section]
\begin{itemize}
\item Une application linéaire vérifie $f(\vec{u}+\vec{v}) = f(\vec{u})+f(\vec{v})$ et $f(k\vec{u}) = kf(\vec{u})$ ; en particulier $f(\vec{0})=\vec{0}$.
\item Une application linéaire est entièrement déterminée par $f(\vec{i})$ et $f(\vec{j})$ : pour $\vec{u}=x\vec{i}+y\vec{j}$, on a $f(\vec{u}) = x\,f(\vec{i})+y\,f(\vec{j})$.
\item La matrice de $f$ dans la base $(\vec{i},\vec{j})$ a pour \textbf{colonnes} les coordonnées de $f(\vec{i})$ et $f(\vec{j})$.
\item Les coordonnées de l'image s'obtiennent par produit matriciel : $f(\vec{u}) \longleftrightarrow M\begin{pmatrix} x \\ y \end{pmatrix}$.
\item Matrices à connaître : identité, homothétie $kI$, symétries d'axes, projection $\begin{pmatrix} 1 & 0 \\ 0 & 0 \end{pmatrix}$, rotation $\begin{pmatrix} \cos\theta & -\sin\theta \\ \sin\theta & \cos\theta \end{pmatrix}$.
\end{itemize}
\end{aretenir}

\begin{savaistu}[Des matrices aux images numériques]
Les applications linéaires du plan sont au cœur des images numériques : quand tu fais pivoter, rétrécir ou déformer une photo sur ton téléphone, le logiciel applique en réalité des matrices $2\times 2$ aux coordonnées de chaque pixel ! Les rotations et symétries que tu viens d'étudier sont exactement celles des boutons « rotation » et « miroir » de ta galerie photo. En infographie et en robotique, on enchaîne ces transformations : c'est le produit de matrices, que tu étudieras dans la prochaine section, qui permet de composer ces effets.
\end{savaistu}

\begin{exercice}[Pour t'entraîner]
Soit $f$ l'application linéaire du plan telle que $f(\vec{i}) = 3\vec{i}-\vec{j}$ et $f(\vec{j}) = 2\vec{i}+\vec{j}$.
\begin{enumerate}
\item Écrire la matrice $M$ de $f$ dans la base $(\vec{i},\vec{j})$.
\item Déterminer les coordonnées de $f(\vec{u})$ pour $\vec{u} = -\vec{i}+4\vec{j}$.
\item Déterminer, si possible, un vecteur non nul $\vec{v}$ tel que $f(\vec{v}) = \vec{0}$. Que peut-on en conjecturer sur $f$ ?
\end{enumerate}
\end{exercice}

\begin{corrige}[Exercice 1]
\textbf{1.} Les images des vecteurs de base s'écrivent en colonnes :
\[ M = \begin{pmatrix} 3 & 2 \\ -1 & 1 \end{pmatrix}. \]
\textbf{2.} Les coordonnées de $\vec{u}$ sont $\begin{pmatrix} -1 \\ 4 \end{pmatrix}$. Produit matriciel :
\[ M\begin{pmatrix} -1 \\ 4 \end{pmatrix} = \begin{pmatrix} 3\times(-1)+2\times 4 \\ -1\times(-1)+1\times 4 \end{pmatrix} = \begin{pmatrix} 5 \\ 5 \end{pmatrix}. \]
Donc $f(\vec{u}) = 5\vec{i}+5\vec{j}$.
\textbf{3.} On cherche $\vec{v}(x\,;\,y)$ tel que $M\begin{pmatrix} x \\ y \end{pmatrix} = \begin{pmatrix} 0 \\ 0 \end{pmatrix}$, c'est-à-dire $\begin{cases} 3x+2y=0 \\ -x+y=0 \end{cases}$. La deuxième équation donne $y=x$ ; en reportant : $3x+2x = 5x = 0$, donc $x=y=0$. Le seul vecteur d'image nulle est $\vec{0}$. On conjecture que $f$ est bijective (un automorphisme) ; cela sera confirmé par le fait que $\det M = 3\times 1 - 2\times(-1) = 5 \neq 0$, donc $M$ est inversible.
\end{corrige}

\coursec{Opérations sur les applications linéaires et automorphismes}

Dans la section précédente, nous avons appris à associer à toute application linéaire $f$ du plan une matrice $M$ dans une base donnée $(\vec{i},\vec{j})$, et à calculer les coordonnées de l'image d'un vecteur par un simple produit matriciel $X' = MX$. Cette correspondance est un pont entre deux mondes : le monde \emph{géométrique} des transformations du plan, et le monde \emph{algébrique} des matrices. La question naturelle qui se pose alors est la suivante : si l'on \emph{combine} des applications linéaires — en les ajoutant, en les multipliant par un réel, en les composant, en les inversant — que devient le pont ? La réponse est magnifique : \textbf{chaque opération sur les applications linéaires se traduit par l'opération correspondante sur les matrices}. C'est précisément ce que nous allons établir dans cette section, et c'est ce qui fait toute la puissance du calcul matriciel.

Dans toute cette section, le plan vectoriel est rapporté à une base fixée $(\vec{i},\vec{j})$, et toutes les matrices sont écrites dans cette base.

\courssub{Somme de deux applications linéaires et produit par un réel}

\textbf{Intuition.} Imaginons deux transformations du plan : $f$ qui dilate et $g$ qui cisaille. Appliquer « $f$ plus $g$ » à un vecteur $\vec{u}$, c'est tout simplement additionner les deux images $f(\vec{u})$ et $g(\vec{u})$. De même, « $2f$ » consiste à doubler chaque image. Ces constructions sont naturelles ; la bonne surprise est qu'elles restent linéaires et que leurs matrices se devinent immédiatement.

\begin{definition}[Somme et produit par un réel]
Soient $f$ et $g$ deux applications linéaires du plan et $k$ un réel.
\begin{itemize}
\item L'application \cle{somme} $f + g$ est définie pour tout vecteur $\vec{u}$ par :
\[ (f+g)(\vec{u}) = f(\vec{u}) + g(\vec{u}). \]
\item L'application \cle{produit par le réel} $k$ est définie pour tout vecteur $\vec{u}$ par :
\[ (k\,f)(\vec{u}) = k\,f(\vec{u}). \]
\end{itemize}
On vérifie sans peine que $f+g$ et $kf$ sont encore des applications linéaires.
\end{definition}

\begin{propriete}[Matrice d'une somme et d'un produit par un réel]
Soient $f$ et $g$ deux applications linéaires de matrices respectives $M$ et $N$ dans la base $(\vec{i},\vec{j})$, et $k$ un réel. Alors :
\[ \mathrm{Mat}(f+g) = M + N \qquad \text{et} \qquad \mathrm{Mat}(k\,f) = k\,M. \]
\end{propriete}

\begin{experience}[Démonstration guidée]
Nous allons utiliser l'\textbf{unicité de la matrice} d'une application linéaire dans une base donnée : si une matrice $A$ vérifie $X' = AX$ pour les coordonnées de tout vecteur et de son image, alors $A$ est \emph{la} matrice de l'application.

Soit $\vec{u}$ un vecteur de coordonnées $X = \begin{pmatrix} x \\ y \end{pmatrix}$.
\begin{itemize}
\item[\textbf{Étape 1.}] Par définition du produit matriciel, $f(\vec{u})$ a pour coordonnées $MX$ et $g(\vec{u})$ a pour coordonnées $NX$.
\item[\textbf{Étape 2.}] Les coordonnées de $(f+g)(\vec{u}) = f(\vec{u}) + g(\vec{u})$ s'obtiennent en additionnant les coordonnées : c'est la colonne $MX + NX$.
\item[\textbf{Étape 3.}] Or, par distributivité du produit matriciel sur l'addition : $MX + NX = (M+N)X$.
\item[\textbf{Étape 4.}] Ainsi, pour tout vecteur de coordonnées $X$, l'image par $f+g$ a pour coordonnées $(M+N)X$. Par unicité de la matrice, $\mathrm{Mat}(f+g) = M+N$.
\end{itemize}
Pour le produit par un réel : les coordonnées de $(kf)(\vec{u}) = k\,f(\vec{u})$ sont $k(MX) = (kM)X$, donc par unicité $\mathrm{Mat}(kf) = kM$. \hfill$\blacksquare$
\end{experience}

\begin{exemplebox}[Somme et produit par un réel]
Soient $f$ et $g$ de matrices $M = \begin{pmatrix} 1 & 2 \\ 0 & 3 \end{pmatrix}$ et $N = \begin{pmatrix} 4 & -1 \\ 2 & 1 \end{pmatrix}$. Alors :
\[ \mathrm{Mat}(f+g) = \begin{pmatrix} 1+4 & 2-1 \\ 0+2 & 3+1 \end{pmatrix} = \begin{pmatrix} 5 & 1 \\ 2 & 4 \end{pmatrix}, \qquad \mathrm{Mat}(3f) = \begin{pmatrix} 3 & 6 \\ 0 & 9 \end{pmatrix}. \]
Vérifions sur un vecteur : pour $\vec{u}$ de coordonnées $X = \begin{pmatrix} 1 \\ 1 \end{pmatrix}$, on a $MX = \begin{pmatrix} 3 \\ 3 \end{pmatrix}$ et $NX = \begin{pmatrix} 3 \\ 3 \end{pmatrix}$, donc $(f+g)(\vec{u})$ a pour coordonnées $\begin{pmatrix} 6 \\ 6 \end{pmatrix}$, ce qui coïncide bien avec $(M+N)X = \begin{pmatrix} 5 & 1 \\ 2 & 4 \end{pmatrix}\begin{pmatrix} 1 \\ 1 \end{pmatrix} = \begin{pmatrix} 6 \\ 6 \end{pmatrix}$.
\end{exemplebox}

\courssub{Composée de deux applications linéaires}

\textbf{Intuition.} Composer, c'est \emph{enchaîner} : on applique d'abord $f$, puis $g$ au résultat. C'est ce que fait un opérateur qui applique d'abord une rotation puis un agrandissement sur une image. Si $f$ est codée par la matrice $M$ et $g$ par la matrice $N$, quel est le code de l'enchaînement $g \circ f$ ? Le schéma suivant résume la situation.

\begin{popfigure}
\begin{center}
\begin{tikzpicture}[scale=1, every node/.style={font=\small}]
\node[draw=popblue, fill=popblueL, rounded corners, thick, minimum width=2.2cm, minimum height=0.9cm] (u) at (0,0) {$\vec{u}$, coord. $X$};
\node[draw=poporange, fill=poporangeL, rounded corners, thick, minimum width=2.2cm, minimum height=0.9cm] (fu) at (4.6,0) {$f(\vec{u})$, coord. $MX$};
\node[draw=poppurple, fill=poppurpleL, rounded corners, thick, minimum width=2.6cm, minimum height=0.9cm] (gfu) at (9.4,0) {$g(f(\vec{u}))$, coord. $(NM)X$};
\draw[-{Stealth[length=3mm]}, very thick, popdark] (u) -- (fu) node[midway, above=1pt]{$f$} node[midway, below=1pt]{matrice $M$};
\draw[-{Stealth[length=3mm]}, very thick, popdark] (fu) -- (gfu) node[midway, above=1pt]{$g$} node[midway, below=1pt]{matrice $N$};
\draw[-{Stealth[length=3mm]}, very thick, popgreen] (u) to[bend right=35] node[midway, below]{$g\circ f$ : matrice $NM$} (gfu);
\end{tikzpicture}
\end{center}
\legende{Composer $f$ puis $g$ revient à multiplier les matrices, dans le même ordre : $N \times M$.}
\end{popfigure}

\begin{propriete}[Théorème — Matrice d'une composée]
Soient $f$ et $g$ deux applications linéaires de matrices respectives $M$ et $N$ dans la base $(\vec{i},\vec{j})$. Alors $g \circ f$ est une application linéaire et :
\[ \mathrm{Mat}(g \circ f) = N \times M. \]
\emph{La matrice de la composée est le produit des matrices, dans le même ordre que la composition.}
\end{propriete}

\begin{experience}[Démonstration guidée (exigible)]
Soit $\vec{u}$ un vecteur de coordonnées $X$.
\begin{itemize}
\item[\textbf{Étape 1.}] Puisque $f$ a pour matrice $M$, le vecteur $f(\vec{u})$ a pour coordonnées $X' = MX$.
\item[\textbf{Étape 2.}] Appliquons $g$ au vecteur $f(\vec{u})$ : puisque $g$ a pour matrice $N$, le vecteur $g(f(\vec{u}))$ a pour coordonnées $NX'$.
\item[\textbf{Étape 3.}] En substituant $X' = MX$ : les coordonnées de $(g\circ f)(\vec{u})$ sont $N(MX)$.
\item[\textbf{Étape 4.}] Par associativité du produit matriciel : $N(MX) = (NM)X$.
\item[\textbf{Étape 5.}] Ainsi, pour tout vecteur de coordonnées $X$, son image par $g\circ f$ a pour coordonnées $(NM)X$. Par unicité de la matrice, $\mathrm{Mat}(g\circ f) = NM$. \hfill$\blacksquare$
\end{itemize}
\end{experience}

\begin{attention}[Erreur fréquente : l'ordre du produit]
Le produit matriciel \textbf{n'est pas commutatif} : en général $NM \neq MN$. Or $g \circ f$ signifie « $f$ d'abord, $g$ ensuite », et sa matrice est $N \times M$ : la matrice de la \emph{première} application apparaît à \emph{droite}. Écrire $\mathrm{Mat}(g\circ f) = MN$ est l'erreur la plus fréquente au Baccalauréat. Moyen mnémotechnique : l'ordre des lettres se conserve, $g \circ f \longrightarrow N \times M$. Et retenez qu'en général $f \circ g \neq g \circ f$ : tourner puis cisailler ne donne pas le même résultat que cisailler puis tourner !
\end{attention}

\begin{exoresolu}[Calcul de la matrice d'une composée]
\textbf{Énoncé.} Dans une base $(\vec{i},\vec{j})$, $f$ a pour matrice $M = \begin{pmatrix} 1 & 2 \\ 0 & 3 \end{pmatrix}$ et $g$ a pour matrice $N = \begin{pmatrix} 4 & -1 \\ 2 & 1 \end{pmatrix}$. Déterminer la matrice de $g \circ f$, puis celle de $f \circ g$, et comparer.
\tcblower
\textbf{Solution.} La matrice de $g \circ f$ est $NM$ :
\begin{align*}
NM &= \begin{pmatrix} 4 & -1 \\ 2 & 1 \end{pmatrix}\begin{pmatrix} 1 & 2 \\ 0 & 3 \end{pmatrix}
= \begin{pmatrix} 4\times 1 + (-1)\times 0 & 4\times 2 + (-1)\times 3 \\ 2\times 1 + 1\times 0 & 2\times 2 + 1\times 3 \end{pmatrix}
= \begin{pmatrix} 4 & 5 \\ 2 & 7 \end{pmatrix}.
\end{align*}
La matrice de $f \circ g$ est $MN$ :
\[ MN = \begin{pmatrix} 1 & 2 \\ 0 & 3 \end{pmatrix}\begin{pmatrix} 4 & -1 \\ 2 & 1 \end{pmatrix} = \begin{pmatrix} 1\times 4 + 2\times 2 & 1\times(-1) + 2\times 1 \\ 0\times 4 + 3\times 2 & 0\times(-1) + 3\times 1 \end{pmatrix} = \begin{pmatrix} 8 & 1 \\ 6 & 3 \end{pmatrix}. \]
On constate que $NM \neq MN$ : les composées $g\circ f$ et $f\circ g$ sont bien deux applications différentes.
\end{exoresolu}

\courssub{Automorphismes et application réciproque}

\textbf{Intuition.} Certaines transformations peuvent être « défaites » : si l'on double toutes les distances (homothétie de rapport $2$), on peut revenir en arrière en divisant par $2$ (homothétie de rapport $\tfrac{1}{2}$). D'autres sont irréversibles : une projection sur une droite « écrase » le plan, et l'information perdue ne peut être récupérée. Les transformations réversibles sont exactement les applications linéaires bijectives.

\begin{definition}[Automorphisme]
On appelle \cle{automorphisme} du plan toute application linéaire \textbf{bijective} du plan sur lui-même. Son application réciproque, notée $f^{-1}$, vérifie :
\[ f \circ f^{-1} = f^{-1} \circ f = \mathrm{id}, \]
où $\mathrm{id}$ est l'application identique, de matrice $I_2 = \begin{pmatrix} 1 & 0 \\ 0 & 1 \end{pmatrix}$.
\end{definition}

\begin{propriete}[Théorème — Caractérisation et matrice de la réciproque]
Soit $f$ une application linéaire de matrice $M$ dans la base $(\vec{i},\vec{j})$.
\begin{enumerate}
\item $f$ est un automorphisme si et seulement si $\det(M) \neq 0$.
\item Dans ce cas, $f^{-1}$ est une application linéaire et :
\[ \mathrm{Mat}(f^{-1}) = M^{-1}. \]
\end{enumerate}
\end{propriete}

\begin{experience}[Démonstration guidée (exigible)]
\textbf{Point 1.} Dire que $f$ est bijective, c'est dire que pour tout vecteur $\vec{u'}$ de coordonnées $X' = \begin{pmatrix} x' \\ y' \end{pmatrix}$, l'équation $f(\vec{u}) = \vec{u'}$, c'est-à-dire $MX = X'$, possède une \textbf{unique} solution $X$. Or nous avons vu dans la section sur les systèmes que le système $MX = X'$ admet une unique solution si et seulement si $\det(M) \neq 0$. D'où l'équivalence.

\textbf{Point 2.} Supposons $f$ bijective. Par définition de la réciproque :
\[ f \circ f^{-1} = \mathrm{id}. \]
D'après le théorème sur la composée, en passant aux matrices :
\[ M \times \mathrm{Mat}(f^{-1}) = \mathrm{Mat}(\mathrm{id}) = I_2. \]
Ainsi $\mathrm{Mat}(f^{-1})$ est une matrice qui, multipliée par $M$, donne $I_2$ : c'est exactement la définition de l'inverse $M^{-1}$. Par unicité de l'inverse, $\mathrm{Mat}(f^{-1}) = M^{-1}$. \hfill$\blacksquare$
\end{experience}

\begin{aretenir}[Le réflexe à avoir]
Pour montrer qu'une application linéaire est un automorphisme, un seul calcul suffit : celui du \textbf{déterminant} de sa matrice. S'il est non nul, l'application est bijective, et la matrice de la réciproque s'obtient par la formule du cours :
\[ M = \begin{pmatrix} a & b \\ c & d \end{pmatrix} \quad\Longrightarrow\quad M^{-1} = \frac{1}{\det(M)}\begin{pmatrix} d & -b \\ -c & a \end{pmatrix}. \]
\end{aretenir}

\begin{exoresolu}[Étude complète : somme, composée, réciproque]
\textbf{Énoncé.} Dans la base $(\vec{i},\vec{j})$, soient $f$ et $g$ les applications linéaires de matrices :
\[ M = \begin{pmatrix} 2 & 1 \\ 1 & 1 \end{pmatrix} \qquad \text{et} \qquad N = \begin{pmatrix} 1 & 0 \\ 1 & 2 \end{pmatrix}. \]
\begin{enumerate}
\item Déterminer la matrice de $f + g$.
\item Déterminer la matrice de $g \circ f$.
\item Montrer que $f$ est un automorphisme et déterminer la matrice de $f^{-1}$.
\item Vérifier que $M \times M^{-1} = I_2$.
\end{enumerate}
\tcblower
\textbf{Solution.}
\begin{enumerate}
\item $\mathrm{Mat}(f+g) = M + N = \begin{pmatrix} 2+1 & 1+0 \\ 1+1 & 1+2 \end{pmatrix} = \begin{pmatrix} 3 & 1 \\ 2 & 3 \end{pmatrix}$.
\item $\mathrm{Mat}(g\circ f) = NM = \begin{pmatrix} 1 & 0 \\ 1 & 2 \end{pmatrix}\begin{pmatrix} 2 & 1 \\ 1 & 1 \end{pmatrix} = \begin{pmatrix} 2 & 1 \\ 4 & 3 \end{pmatrix}$.
\item $\det(M) = 2\times 1 - 1\times 1 = 1 \neq 0$, donc $f$ est un automorphisme. De plus :
\[ \mathrm{Mat}(f^{-1}) = M^{-1} = \frac{1}{1}\begin{pmatrix} 1 & -1 \\ -1 & 2 \end{pmatrix} = \begin{pmatrix} 1 & -1 \\ -1 & 2 \end{pmatrix}. \]
\item Vérification : $M \times M^{-1} = \begin{pmatrix} 2 & 1 \\ 1 & 1 \end{pmatrix}\begin{pmatrix} 1 & -1 \\ -1 & 2 \end{pmatrix} = \begin{pmatrix} 2-1 & -2+2 \\ 1-1 & -1+2 \end{pmatrix} = \begin{pmatrix} 1 & 0 \\ 0 & 1 \end{pmatrix} = I_2$. \hfill$\blacksquare$
\end{enumerate}
\end{exoresolu}

\courssub{Interprétation géométrique : défaire une transformation}

Appliquer $f^{-1}$ après $f$, c'est \textbf{ramener chaque point à sa position initiale}. La figure ci-dessous illustre cette idée avec une homothétie de rapport $2$ (matrice $M = \begin{pmatrix} 2 & 0 \\ 0 & 2 \end{pmatrix}$) suivie de sa réciproque, l'homothétie de rapport $\tfrac{1}{2}$ (matrice $M^{-1} = \begin{pmatrix} 1/2 & 0 \\ 0 & 1/2 \end{pmatrix}$) : le triangle final se superpose exactement au triangle initial.

\begin{popfigure}
\begin{center}
\begin{tikzpicture}[scale=0.85]
% Triangle initial
\fill[popblueL] (0,0) -- (1.2,0) -- (0.4,0.9) -- cycle;
\draw[popblue, very thick] (0,0) -- (1.2,0) -- (0.4,0.9) -- cycle;
\node[popblue, below] at (0.6,-0.05) {\small figure initiale};
% Flèche f
\draw[-{Stealth[length=3mm]}, very thick, popdark] (1.7,0.45) -- (3.1,0.45) node[midway, above]{\small $f$ : rapport $2$};
% Triangle agrandi
\fill[poporangeL] (3.6,-0.45) -- (6.0,-0.45) -- (4.4,1.35) -- cycle;
\draw[poporange, very thick] (3.6,-0.45) -- (6.0,-0.45) -- (4.4,1.35) -- cycle;
\node[poporange, below] at (4.8,-0.5) {\small image par $f$};
% Flèche f^{-1}
\draw[-{Stealth[length=3mm]}, very thick, popdark] (6.5,0.45) -- (7.9,0.45) node[midway, above]{\small $f^{-1}$ : rapport $\tfrac12$};
% Triangle final
\fill[popgreenL] (8.4,0) -- (9.6,0) -- (8.8,0.9) -- cycle;
\draw[popgreen, very thick] (8.4,0) -- (9.6,0) -- (8.8,0.9) -- cycle;
\node[popgreen, below] at (9.0,-0.05) {\small figure retrouvée};
\end{tikzpicture}
\end{center}
\legende{Composer $f$ puis $f^{-1}$ redonne l'identité : $M^{-1}M = I_2$, la figure finale coïncide avec la figure initiale.}
\end{popfigure}

\begin{attention}[Conditions d'application]
Deux pièges à éviter :
\begin{itemize}
\item On ne peut parler de $f^{-1}$ que si $\det(M) \neq 0$. Avant tout calcul d'inverse, \textbf{calcule d'abord le déterminant} et annonce clairement qu'il est non nul.
\item Ne confonds pas $\mathrm{Mat}(f^{-1}) = M^{-1}$ avec $\mathrm{Mat}(f \circ f) = M^2$ : la notation $f^{-1}$ désigne la \emph{réciproque} (qui n'existe que si $f$ est bijective), pas une puissance.
\end{itemize}
\end{attention}

\begin{savaistu}[Les matrices dans ton téléphone]
Quand tu retouches une photo sur ton smartphone — rotation de $30^\circ$, puis agrandissement de $150\%$, puis symétrie — le logiciel ne transforme pas l'image trois fois : il \textbf{multiplie les matrices} des trois transformations pour n'obtenir qu'\emph{une seule} matrice, appliquée en une passe à chacun des millions de pixels. C'est exactement le théorème $\mathrm{Mat}(g\circ f) = NM$ à l'œuvre ! Et quand tu appuies sur « annuler », le logiciel applique la matrice inverse $M^{-1}$. Les jeux vidéo, les effets spéciaux du cinéma et la cartographie numérique utilisée par les ingénieurs de CAMTEL ou de l'IGN reposent sur ce même principe.
\end{savaistu}

\courssub{Synthèse : le dictionnaire applications linéaires $\leftrightarrow$ matrices}

Toute cette section peut se résumer en un tableau de correspondance, véritable « dictionnaire » entre le langage des applications et celui des matrices. Apprends-le par cœur : il condense l'essentiel de ce qui est exigible au Baccalauréat.

\begin{center}
\begin{tabularx}{\linewidth}{|>{\columncolor{popblueL}}l|>{\columncolor{poporangeL}}X|}
\hline
\rowcolor{popblue!40} \textbf{Application linéaire} & \textbf{Matrice (dans la base $(\vec{i},\vec{j})$)} \\
\hline
Image de $\vec{u}$ (coordonnées $X$) par $f$ & $X' = MX$ \\
\hline
Somme $f + g$ & $M + N$ \\
\hline
Produit par un réel $k\,f$ & $k\,M$ \\
\hline
Composée $g \circ f$ ($f$ d'abord, $g$ ensuite) & $N \times M$ (même ordre) \\
\hline
Application identique $\mathrm{id}$ & $I_2 = \begin{pmatrix} 1 & 0 \\ 0 & 1 \end{pmatrix}$ \\
\hline
$f$ est un automorphisme (bijective) & $\det(M) \neq 0$ \\
\hline
Réciproque $f^{-1}$ & $M^{-1}$ \\
\hline
$f \circ f^{-1} = f^{-1} \circ f = \mathrm{id}$ & $M\,M^{-1} = M^{-1}M = I_2$ \\
\hline
\end{tabularx}
\end{center}

\begin{methode}[Méthode type au Baccalauréat]
Face à un exercice mêlant applications linéaires et matrices :
\begin{enumerate}
\item \textbf{Identifier} les matrices $M$ et $N$ des applications en présence (colonnes = images des vecteurs de base).
\item \textbf{Traduire} la question en langage matriciel à l'aide du dictionnaire ci-dessus : somme $\to$ $M+N$, composée $g\circ f$ $\to$ $NM$, réciproque $\to$ $M^{-1}$.
\item \textbf{Calculer} soigneusement, en respectant l'ordre des facteurs pour les produits.
\item \textbf{Conclure} en revenant au langage des applications : « donc la matrice de $g\circ f$ est ... », « donc $f$ est un automorphisme et $f^{-1}$ a pour matrice ... ».
\end{enumerate}
\end{methode}

\begin{exercice}[Pour t'entraîner]
Dans une base $(\vec{i},\vec{j})$, soient $f$ et $g$ les applications linéaires de matrices $M = \begin{pmatrix} 3 & -1 \\ 2 & 1 \end{pmatrix}$ et $N = \begin{pmatrix} 0 & 2 \\ 1 & 1 \end{pmatrix}$.
\begin{enumerate}
\item Déterminer les matrices de $2f - g$, de $g \circ f$ et de $f \circ g$.
\item $f$ est-elle un automorphisme ? Si oui, déterminer la matrice de $f^{-1}$.
\item Soit $\vec{u}$ de coordonnées $\begin{pmatrix} 1 \\ 2 \end{pmatrix}$. Calculer les coordonnées de $(g\circ f)(\vec{u})$ de deux façons différentes.
\end{enumerate}
\end{exercice}

\begin{corrige}[Exercice]
\begin{enumerate}
\item $2M - N = \begin{pmatrix} 6 & -2 \\ 4 & 2 \end{pmatrix} - \begin{pmatrix} 0 & 2 \\ 1 & 1 \end{pmatrix} = \begin{pmatrix} 6 & -4 \\ 3 & 1 \end{pmatrix}$. \quad $NM = \begin{pmatrix} 0 & 2 \\ 1 & 1 \end{pmatrix}\begin{pmatrix} 3 & -1 \\ 2 & 1 \end{pmatrix} = \begin{pmatrix} 4 & 2 \\ 5 & 0 \end{pmatrix}$. \quad $MN = \begin{pmatrix} 3 & -1 \\ 2 & 1 \end{pmatrix}\begin{pmatrix} 0 & 2 \\ 1 & 1 \end{pmatrix} = \begin{pmatrix} -1 & 5 \\ 1 & 5 \end{pmatrix}$.
\item $\det(M) = 3\times 1 - (-1)\times 2 = 5 \neq 0$ : $f$ est un automorphisme et $M^{-1} = \dfrac{1}{5}\begin{pmatrix} 1 & 1 \\ -2 & 3 \end{pmatrix}$.
\item Première façon : $MX = \begin{pmatrix} 3-2 \\ 2+2 \end{pmatrix} = \begin{pmatrix} 1 \\ 4 \end{pmatrix}$, puis $N(MX) = \begin{pmatrix} 8 \\ 5 \end{pmatrix}$. Deuxième façon : $(NM)X = \begin{pmatrix} 4 & 2 \\ 5 & 0 \end{pmatrix}\begin{pmatrix} 1 \\ 2 \end{pmatrix} = \begin{pmatrix} 8 \\ 5 \end{pmatrix}$. Les deux résultats coïncident, comme le garantit le théorème.
\end{enumerate}
\end{corrige}

\coursec{Applications géométriques et synthèse}

Dans la section précédente, tu as appris à composer des applications linéaires et à inverser un automorphisme en manipulant leurs matrices : la composée correspond au produit des matrices, la réciproque à la matrice inverse. Tu disposes donc désormais de toute la machinerie algébrique. Il reste une question fondamentale, celle qui donne tout son sens à cette leçon : \emph{que voit-on géométriquement quand on regarde une matrice ?} Une matrice $2\times 2$ n'est pas qu'un tableau de nombres : c'est le portrait robot d'une transformation du plan. Dans cette dernière section, nous allons apprendre à \cle{lire} une matrice comme on lit une carte : identifier la transformation qu'elle code, mesurer son effet sur les aires, trouver ses points fixes, puis résoudre un problème de synthèse complet, comme au Baccalauréat.

\courssub{Transformations usuelles et leurs matrices}

\courssub{Lecture géométrique des colonnes : l'effet sur le carré unité}

Voici l'idée directrice, simple et puissante. Soit $f$ l'application linéaire de matrice $M = \begin{pmatrix} a & b \\ c & d \end{pmatrix}$ dans la base $(\vec{i}, \vec{j})$. Par définition même de la matrice d'une application linéaire :
\begin{itemize}
\item la \cle{première colonne} $\begin{pmatrix} a \\ c \end{pmatrix}$ contient les coordonnées de $f(\vec{i})$ ;
\item la \cle{deuxième colonne} $\begin{pmatrix} b \\ d \end{pmatrix}$ contient les coordonnées de $f(\vec{j})$.
\end{itemize}

Or le carré unité, de sommets $O(0;0)$, $I(1;0)$, $K(1;1)$, $J(0;1)$, est entièrement déterminé par $\vec{i} = \overrightarrow{OI}$ et $\vec{j} = \overrightarrow{OJ}$. Son image par $f$ est le \cle{parallélogramme} construit sur les vecteurs $f(\vec{i})$ et $f(\vec{j})$, c'est-à-dire sur les deux colonnes de $M$. \textbf{Regarder les colonnes de $M$, c'est regarder ce que devient le carré unité.}

\begin{popfigure}
\begin{center}
\begin{tikzpicture}[scale=0.95, >=Stealth]
% --- Styles communs ---
\tikzset{
  axes/.style={->, popmuted},
  grid/.style={popmuted, very thin},
  vec/.style={->, very thick},
  square/.style={thick, fill opacity=0.3},
  labelnode/.style={below, popdark, font=\scriptsize}
}
% --- 1. Homothétie k=1,5 (Haut-Gauche) ---
\begin{scope}
\draw[axes] (-1.5,0) -- (2.0,0);
\draw[axes] (0,-1.5) -- (0,2.0);
\fill[popblueL] (0,0) -- (1.5,0) -- (1.5,1.5) -- (0,1.5) -- cycle;
\draw[popblue, square] (0,0) -- (1.5,0) -- (1.5,1.5) -- (0,1.5) -- cycle;
\draw[vec, popblue] (0,0) -- (1.5,0) node[midway, below] {$\vec{i}'$};
\draw[vec, popblue] (0,0) -- (0,1.5) node[midway, left] {$\vec{j}'$};
\node[labelnode, align=center] at (0.75,-0.7) {$\begin{pmatrix} 1,5 & 0 \\ 0 & 1,5 \end{pmatrix}$ : homothétie $k=1,5$};
\end{scope}
% --- 2. Symétrie axe (OI) (Haut-Droite) ---
\begin{scope}[xshift=5.8cm]
\draw[axes] (-1.5,0) -- (2.0,0);
\draw[axes] (0,-2.0) -- (0,1.5);
\fill[poporangeL] (0,0) -- (1,0) -- (1,-1) -- (0,-1) -- cycle;
\draw[poporange, square] (0,0) -- (1,0) -- (1,-1) -- (0,-1) -- cycle;
\draw[vec, poporange] (0,0) -- (1,0) node[midway, below] {$\vec{i}'$};
\draw[vec, poporange] (0,0) -- (0,-1) node[midway, right] {$\vec{j}'$};
\node[labelnode, align=center] at (0.5,-1.3) {$\begin{pmatrix} 1 & 0 \\ 0 & -1 \end{pmatrix}$ : symétrie axiale $(O,\vec{i})$};
\end{scope}
% --- 3. Cisaillement (Bas-Gauche) ---
\begin{scope}[yshift=-4.8cm]
\draw[axes] (-1.5,0) -- (2.5,0);
\draw[axes] (0,-1.5) -- (0,2.0);
\fill[poppurpleL] (0,0) -- (1,0) -- (2,1) -- (1,1) -- cycle;
\draw[poppurple, square] (0,0) -- (1,0) -- (2,1) -- (1,1) -- cycle;
\draw[vec, poppurple] (0,0) -- (1,0) node[midway, below] {$\vec{i}'$};
\draw[vec, poppurple] (0,0) -- (1,1) node[midway, left] {$\vec{j}'$};
\node[labelnode, align=center] at (1.0,-0.7) {$\begin{pmatrix} 1 & 1 \\ 0 & 1 \end{pmatrix}$ : cisaillement axe $(O,\vec{i})$};
\end{scope}
% --- 4. Rotation $\pi/2$ (Bas-Droite) ---
\begin{scope}[xshift=5.8cm, yshift=-4.8cm]
\draw[axes] (-2.0,0) -- (1.5,0);
\draw[axes] (0,-1.5) -- (0,2.0);
\fill[popgreenL] (0,0) -- (0,1) -- (-1,1) -- (-1,0) -- cycle;
\draw[popgreen, square] (0,0) -- (0,1) -- (-1,1) -- (-1,0) -- cycle;
\draw[vec, popgreen] (0,0) -- (0,1) node[midway, right] {$\vec{i}'$};
\draw[vec, popgreen] (0,0) -- (-1,0) node[midway, below] {$\vec{j}'$};
\node[labelnode, align=center] at (-0.5,-0.7) {$\begin{pmatrix} 0 & -1 \\ 1 & 0 \end{pmatrix}$ : rotation $+\pi/2$};
\end{scope}
\end{tikzpicture}
\end{center}
\legende{Le carré unité (en pointillés implicites aux axes) et son image par quatre matrices : homothétie (agrandissement), symétrie axiale (retournement), cisaillement (glissement) et rotation (quart de tour direct). Les flèches épaisses colorées sont les colonnes des matrices, images de $\vec{i}$ et $\vec{j}$.}
\end{popfigure}

\begin{propriete}[Matrices des transformations usuelles]
Dans une base orthonormée directe $(\vec{i}, \vec{j})$ du plan :
\begin{itemize}
\item l'\cle{homothétie} de centre $O$ et de rapport $k$ a pour matrice $\begin{pmatrix} k & 0 \\ 0 & k \end{pmatrix}$ ;
\item la \cle{rotation} de centre $O$ et d'angle $\alpha$ (orienté) a pour matrice $\begin{pmatrix} \cos\alpha & -\sin\alpha \\ \sin\alpha & \cos\alpha \end{pmatrix}$ ;
\item la \cle{symétrie} d'axe $(O,\vec{i})$ a pour matrice $\begin{pmatrix} 1 & 0 \\ 0 & -1 \end{pmatrix}$, celle d'axe $(O,\vec{j})$ a pour matrice $\begin{pmatrix} -1 & 0 \\ 0 & 1 \end{pmatrix}$, celle d'axe la première bissectrice a pour matrice $\begin{pmatrix} 0 & 1 \\ 1 & 0 \end{pmatrix}$ ;
\item une \cle{affinité orthogonale} d'axe $(O,\vec{i})$, de direction $(O,\vec{j})$ et de rapport $k$ a pour matrice $\begin{pmatrix} 1 & 0 \\ 0 & k \end{pmatrix}$ (elle conserve les abscisses et multiplie les ordonnées par $k$) ;
\item un \cle{cisaillement} d'axe $(O,\vec{i})$ a pour matrice $\begin{pmatrix} 1 & \lambda \\ 0 & 1 \end{pmatrix}$.
\end{itemize}
\end{propriete}

\begin{methode}[Identifier géométriquement une transformation à partir de sa matrice]
Pour reconnaître la transformation associée à une matrice $M = \begin{pmatrix} a & b \\ c & d \end{pmatrix}$ dans une base orthonormée :
\begin{enumerate}
\item \textbf{Calculer les images des vecteurs de base} : la colonne 1 donne $f(\vec{i})$, la colonne 2 donne $f(\vec{j})$. Placer ces vecteurs sur une figure mentale (ou réelle !).
\item \textbf{Comparer avec les matrices usuelles} : matrice scalaire $\begin{pmatrix} k & 0 \\ 0 & k \end{pmatrix}$ (homothétie), matrice de rotation (colonnes orthonormées, $\det=1$), matrice diagonale avec des $1$ et $-1$ (symétrie axiale), matrice triangulaire avec des $1$ sur la diagonale (cisaillement), etc.
\item \textbf{Calculer le déterminant} : son signe indique si l'orientation est conservée ($\det > 0$) ou renversée ($\det < 0$, cas des symétries axiales) ; sa valeur absolue donne le facteur d'agrandissement des aires.
\item \textbf{Chercher les points fixes} (voir plus loin) : une droite de points fixes révèle une symétrie axiale, une affinité ou un cisaillement ; le seul point fixe $O$ oriente vers une homothétie ($k\neq 1$) ou une rotation ($\alpha\not\equiv 0[2\pi]$).
\item \textbf{Conclure en donnant les éléments caractéristiques} : centre, axe, rapport, angle, direction.
\end{enumerate}
\end{methode}

\begin{exemplebox}[Étude complète de la matrice $M = \begin{pmatrix} 0 & -1 \\ 1 & 0 \end{pmatrix}$]
Appliquons la méthode pas à pas.
\begin{enumerate}
\item \textbf{Images des vecteurs de base} : $f(\vec{i}) = \begin{pmatrix} 0 \\ 1 \end{pmatrix} = \vec{j}$ et $f(\vec{j}) = \begin{pmatrix} -1 \\ 0 \end{pmatrix} = -\vec{i}$. Le vecteur $\vec{i}$ « monte » sur $\vec{j}$, et $\vec{j}$ « recule » sur $-\vec{i}$ : tout le plan a tourné d'un quart de tour dans le sens direct.
\item \textbf{Déterminant} : $\det(M) = 0\times 0 - (-1)\times 1 = 1 > 0$ : l'orientation et les aires sont conservées.
\item \textbf{Points fixes} : $MX = X$ s'écrit $\begin{cases} -y = x \\ x = y \end{cases}$, d'où $x = y = 0$ : seul $O$ est fixe.
\item \textbf{Conclusion} : $f$ est la \cle{rotation} de centre $O$ et d'angle $\dfrac{\pi}{2}$ (quart de tour direct). Vérification sur un point : l'image de $A(2;1)$ est $A'(-1;2)$, et on a bien $OA = OA' = \sqrt{5}$ avec $\left(\overrightarrow{OA}, \overrightarrow{OA'}\right) = \dfrac{\pi}{2}$.
\end{enumerate}
\end{exemplebox}

\courssub{Effet sur les aires : le rôle du déterminant}

Reprenons le carré unité, d'aire $1$. Son image est le parallélogramme construit sur les vecteurs colonnes $f(\vec{i}) = (a;c)$ et $f(\vec{j}) = (b;d)$. L'aire de ce parallélogramme vaut $|ad - bc|$, c'est-à-dire exactement $|\det(M)|$. Ce résultat, observé sur le carré unité, est en fait général.

\begin{propriete}[Aire de l'image d'une figure — admis]
Soit $f$ une application linéaire du plan de matrice $M$, et $\mathcal{F}$ une figure plane d'aire $\mathcal{A}$. Alors l'image $f(\mathcal{F})$ a pour aire :
\[ \mathcal{A}' = |\det(M)| \times \mathcal{A}. \]
Autrement dit, \textbf{toutes les aires sont multipliées par la valeur absolue du déterminant}. En particulier, $f$ conserve les aires si et seulement si $|\det(M)| = 1$.
\end{propriete}

\begin{exemplebox}[Illustration numérique]
Reprenons nos exemples :
\begin{itemize}
\item homothétie de rapport $1,5$ : $\det = 1,5^2 = 2,25$. Un champ de \SI{400}{m^2} vu à travers cette transformation devient un champ de $2,25 \times 400 = \SI{900}{m^2}$. C'est le résultat bien connu : les aires sont multipliées par le \emph{carré} du rapport ;
\item symétrie axiale : $\det = -1$, donc $|\det| = 1$ : les aires sont conservées (heureusement, un miroir ne grossit pas !) ;
\item cisaillement $\begin{pmatrix} 1 & 1 \\ 0 & 1 \end{pmatrix}$ : $\det = 1$. Le carré unité devient un parallélogramme \emph{de même aire} : c'est le principe du « paquet de feuilles qu'on fait glisser » — le volume ne change pas ;
\item rotation : $\det = 1$. Les aires et l'orientation sont conservées.
\end{itemize}
\end{exemplebox}

\begin{attention}[$\det(M) = 1$ ne signifie pas $M = I_2$ !]
C'est l'une des confusions les plus fréquentes au Bac. Dire que $\det(M) = 1$ signifie seulement que la transformation \textbf{conserve les aires et l'orientation}. Cela ne veut \emph{pas} dire que $M$ est la matrice identité ni que la transformation est l'identité. Contre-exemples éclatants : le cisaillement $\begin{pmatrix} 1 & 1 \\ 0 & 1 \end{pmatrix}$ et la rotation $\begin{pmatrix} 0 & -1 \\ 1 & 0 \end{pmatrix}$ ont tous deux un déterminant égal à $1$, et pourtant aucun n'est l'identité — le premier déforme le carré en parallélogramme, la seconde fait tourner tout le plan. Réciproquement, $M = I_2$ implique bien $\det(M) = 1$, mais l'implication ne fonctionne que dans ce sens. Retiens : \textbf{le déterminant mesure un effet (sur les aires), il ne caractérise pas la transformation}.
\end{attention}

\courssub{Points fixes d'une application linéaire}

\begin{definition}[Point fixe]
Soit $f$ une application linéaire du plan de matrice $M$. Un point $P(x;y)$, de vecteur colonne $X = \begin{pmatrix} x \\ y \end{pmatrix}$, est dit \cle{fixe} par $f$ lorsque $f(P) = P$, c'est-à-dire :
\[ MX = X \quad \Longleftrightarrow \quad (M - I_2)X = \begin{pmatrix} 0 \\ 0 \end{pmatrix}. \]
L'origine $O$ est toujours un point fixe d'une application linéaire (car $f(\vec{0}) = \vec{0}$). La question intéressante est : \emph{y en a-t-il d'autres ?}
\end{definition}

\begin{methode}[Trouver les points fixes d'une application linéaire]
\begin{enumerate}
\item Écrire la matrice $M - I_2 = \begin{pmatrix} a-1 & b \\ c & d-1 \end{pmatrix}$.
\item Résoudre le système homogène $(M - I_2)X = 0$, c'est-à-dire $\begin{cases} (a-1)x + by = 0 \\ cx + (d-1)y = 0 \end{cases}$.
\item Interpréter géométriquement :
\begin{itemize}
\item si $\det(M - I_2) \neq 0$ : le système admet l'unique solution $(0;0)$ — \textbf{$O$ est le seul point fixe} (cas des homothéties de rapport $k \neq 1$, des rotations d'angle non nul modulo $2\pi$) ;
\item si $\det(M - I_2) = 0$ : les deux équations sont proportionnelles, l'ensemble des solutions est une \textbf{droite de points fixes} passant par $O$ (cas des symétries axiales, des affinités, des cisaillements) — ou le plan entier si $M = I_2$.
\end{itemize}
\end{enumerate}
\end{methode}

\begin{exoresolu}[Points fixes d'une affinité]
On considère l'application linéaire $f$ de matrice $M = \begin{pmatrix} 1 & 0 \\ 0 & 3 \end{pmatrix}$ dans une base orthonormée $(\vec{i},\vec{j})$. Déterminer l'ensemble des points fixes de $f$ et en déduire la nature géométrique de $f$.
\tcblower
\textbf{Solution détaillée.}
\begin{enumerate}
\item On calcule $M - I_2 = \begin{pmatrix} 0 & 0 \\ 0 & 2 \end{pmatrix}$.
\item Le système $(M - I_2)X = 0$ s'écrit :
\[ \begin{cases} 0x + 0y = 0 \\ 2y = 0 \end{cases} \Longleftrightarrow y = 0. \]
\item L'ensemble des points fixes est donc la droite d'équation $y = 0$, c'est-à-dire l'axe $(O,\vec{i})$.
\end{enumerate}
\textbf{Interprétation} : $f$ conserve l'abscisse de chaque point et multiplie son ordonnée par $3$ ; les points de l'axe des abscisses ne bougent pas. $f$ est l'\cle{affinité orthogonale} d'axe $(O,\vec{i})$, de direction $(O,\vec{j})$ et de rapport $3$. Vérification : l'image de $P(2;1)$ est $P'(2;3)$ — même abscisse, ordonnée triplée, et $P'$ s'obtient bien en « s'éloignant » de l'axe des abscisses dans la direction verticale.
\end{exoresolu}

\courssub{Problème de synthèse type Bac}

\begin{exoresolu}[Modélisation économique — production d'une coopérative]
Une coopérative agricole de la région de l'Ouest transforme chaque semaine du maïs et du manioc. On note $x_n$ et $y_n$ les quantités (en sacs) de maïs et de manioc produites la semaine $n$. D'une semaine à l'autre, la production évolue selon le schéma :
\[ \begin{cases} x_{n+1} = {0,8}\,x_n + {0,3}\,y_n \\ y_{n+1} = {0,2}\,x_n + {0,7}\,y_n \end{cases} \]
On pose $X_n = \begin{pmatrix} x_n \\ y_n \end{pmatrix}$ et $A = \begin{pmatrix} {0,8} & {0,3} \\ {0,2} & {0,7} \end{pmatrix}$.
\begin{enumerate}
\item Écrire la relation de récurrence sous forme matricielle, puis exprimer $X_n$ en fonction de $A$, $n$ et $X_0$.
\item La première semaine, la coopérative produit $x_0 = 100$ sacs de maïs et $y_0 = 50$ sacs de manioc. Calculer la production de la semaine $1$.
\item Calculer $\det(A)$, montrer que $A$ est inversible et déterminer $A^{-1}$.
\item La semaine $n+1$, la coopérative a produit $x_{n+1} = 110$ sacs de maïs et $y_{n+1} = 90$ sacs de manioc. Retrouver la production de la semaine $n$.
\end{enumerate}
\tcblower
\textbf{Solution détaillée.}
\begin{enumerate}
\item La relation s'écrit $X_{n+1} = A X_n$. Par récurrence immédiate (itération du produit matriciel) : $X_n = A^n X_0$.
\item On calcule :
\[ X_1 = A X_0 = \begin{pmatrix} {0,8} & {0,3} \\ {0,2} & {0,7} \end{pmatrix} \begin{pmatrix} 100 \\ 50 \end{pmatrix} = \begin{pmatrix} {0,8} \times 100 + {0,3} \times 50 \\ {0,2} \times 100 + {0,7} \times 50 \end{pmatrix} = \begin{pmatrix} 95 \\ 55 \end{pmatrix}. \]
La semaine $1$ : $95$ sacs de maïs et $55$ sacs de manioc.
\item $\det(A) = {0,8} \times {0,7} - {0,3} \times {0,2} = 0,56 - 0,06 = 0,5 \neq 0$, donc $A$ est inversible et :
\[ A^{-1} = \frac{1}{0,5}\begin{pmatrix} {0,7} & -{0,3} \\ -{0,2} & {0,8} \end{pmatrix} = \begin{pmatrix} 1,4 & -0,6 \\ -0,4 & 1,6 \end{pmatrix}. \]
\item De $X_{n+1} = A X_n$, on tire $X_n = A^{-1} X_{n+1}$ :
\[ X_n = \begin{pmatrix} 1,4 & -0,6 \\ -0,4 & 1,6 \end{pmatrix} \begin{pmatrix} 110 \\ 90 \end{pmatrix} = \begin{pmatrix} 1,4 \times 110 - 0,6 \times 90 \\ -0,4 \times 110 + 1,6 \times 90 \end{pmatrix} = \begin{pmatrix} 100 \\ 100 \end{pmatrix}. \]
La semaine $n$, la coopérative avait produit $100$ sacs de maïs et $100$ sacs de manioc. Vérification : $A\begin{pmatrix} 100 \\ 100 \end{pmatrix} = \begin{pmatrix} 110 \\ 90 \end{pmatrix}$. C'est correct.
\end{enumerate}
\textbf{Remarque de synthèse} : ce problème mobilise toute la leçon — écriture matricielle d'un système, produit matrice-vecteur, déterminant, inverse, et résolution par matrice inverse. C'est exactement l'esprit des problèmes du Baccalauréat TI.
\end{exoresolu}

\begin{exercice}[Pour t'entraîner]
Soit $M = \begin{pmatrix} 2 & 1 \\ 1 & 2 \end{pmatrix}$ la matrice d'une application linéaire $f$ dans une base orthonormée.
\begin{enumerate}
\item Calculer $\det(M)$. $f$ est-elle un automorphisme ? Conserve-t-elle les aires ?
\item Déterminer les points fixes de $f$.
\item Un triangle a pour aire $6$ cm$^2$. Quelle est l'aire de son image par $f$ ?
\item Calculer $M^{-1}$ et en déduire l'antécédent du point $B(5;4)$.
\end{enumerate}
\end{exercice}

\begin{corrige}[Exercice 1]
\begin{enumerate}
\item $\det(M) = 4 - 1 = 3 \neq 0$ : $f$ est un automorphisme. Comme $|\det(M)| = 3 \neq 1$, $f$ ne conserve pas les aires : elle les multiplie par $3$.
\item $M - I_2 = \begin{pmatrix} 1 & 1 \\ 1 & 1 \end{pmatrix}$. Le système $(M - I_2)X = 0$ se réduit à $x + y = 0$ : les points fixes forment la droite d'équation $y = -x$ (la deuxième bissectrice).
\item $\mathcal{A}' = 3 \times 6 = 18$ cm$^2$.
\item $M^{-1} = \dfrac{1}{3}\begin{pmatrix} 2 & -1 \\ -1 & 2 \end{pmatrix}$. L'antécédent de $B$ est $M^{-1}\begin{pmatrix} 5 \\ 4 \end{pmatrix} = \dfrac{1}{3}\begin{pmatrix} 6 \\ 3 \end{pmatrix} = \begin{pmatrix} 2 \\ 1 \end{pmatrix}$ : c'est le point $A(2;1)$. Vérification : $M\begin{pmatrix} 2 \\ 1 \end{pmatrix} = \begin{pmatrix} 5 \\ 4 \end{pmatrix}$.
\end{enumerate}
\end{corrige}

\begin{savaistu}[Les matrices au service des secrets : le chiffrement de Hill]
En 1929, le mathématicien américain Lester Hill a inventé un procédé de cryptographie fondé... sur les matrices inversibles ! Le principe : on remplace chaque lettre par un nombre (A $= 0$, B $= 1$, ..., Z $= 25$), on regroupe les nombres du message par paires dans des vecteurs colonnes, puis on multiplie chaque vecteur par une matrice secrète $M$ inversible (la \emph{clé}). Le message chiffré est la suite des vecteurs obtenus. Pour déchiffrer, le destinataire — qui connaît $M^{-1}$ — multiplie simplement les vecteurs reçus par $M^{-1}$ et retrouve le message initial. Sans la matrice inverse, le message reste illisible ! C'est exactement la technique que tu as utilisée pour retrouver la production de la semaine $n$ dans le problème de la coopérative. Aujourd'hui, des versions très perfectionnées de ces idées (arithmétique modulaire, grandes matrices) protègent tes transactions Mobile Money et tes messages WhatsApp : les mathématiques de Terminale sont au cœur de la sécurité numérique.
\end{savaistu}

\courssub{Bilan de la leçon : carte mentale}

\begin{aretenir}[L'essentiel de la leçon]
\begin{itemize}
\item Une \textbf{matrice $2\times 2$} se manipule par somme, produit par un réel et produit (non commutatif !) : $AB \neq BA$ en général.
\item Le \textbf{déterminant} $\det(M) = ad - bc$ décide de l'inversibilité : $M$ inversible $\Leftrightarrow \det(M) \neq 0$, et alors $M^{-1} = \dfrac{1}{\det(M)}\begin{pmatrix} d & -b \\ -c & a \end{pmatrix}$.
\item Un \textbf{système linéaire} $2\times 2$ s'écrit $MX = B$ et se résout par $X = M^{-1}B$ quand $M$ est inversible.
\item Toute \textbf{application linéaire} du plan est codée par sa matrice dans une base : les \textbf{colonnes} sont les images des vecteurs de base. Composée $\leftrightarrow$ produit ; réciproque $\leftrightarrow$ inverse.
\item Géométriquement : $|\det(M)|$ est le \textbf{facteur des aires}, le signe de $\det(M)$ dit si l'\textbf{orientation} est conservée, et l'équation $MX = X$ donne les \textbf{points fixes}.
\end{itemize}
\end{aretenir}

\begin{popfigure}
\begin{center}
\begin{tikzpicture}[mindmap, grow cyclic, every node/.style={concept, text=white, font=\small\bfseries},
root concept/.append style={concept color=popblue, font=\normalsize\bfseries},
level 1/.append style={level distance=3.4cm, sibling angle=72},
level 2/.append style={level distance=2.2cm, sibling angle=50},
concept color=popblue]
\node[root concept]{Matrices $2\times 2$}
child[concept color=popgreen]{ node {Opérations}
  child[concept color=popgreen!80]{ node[font=\scriptsize\bfseries] {Somme, produit par un réel} }
  child[concept color=popgreen!80]{ node[font=\scriptsize\bfseries] {Produit $AB$ (non commutatif)} }
}
child[concept color=poporange]{ node {Déterminant}
  child[concept color=poporange!80]{ node[font=\scriptsize\bfseries] {$\det = ad - bc$} }
  child[concept color=poporange!80]{ node[font=\scriptsize\bfseries] {Aires $\times\,|\det|$} }
}
child[concept color=poppurple]{ node {Inverse}
  child[concept color=poppurple!80]{ node[font=\scriptsize\bfseries] {Existe ssi $\det \neq 0$} }
  child[concept color=poppurple!80]{ node[font=\scriptsize\bfseries, align=center] {$M^{-1} = \frac{1}{\det}\begin{pmatrix} d & -b \\ -c & a \end{pmatrix}$} }
}
child[concept color=popink]{ node {Systèmes}
  child[concept color=popink!80]{ node[font=\scriptsize\bfseries] {$MX = B \Rightarrow X = M^{-1}B$} }
}
child[concept color=popgold]{ node {Applications linéaires}
  child[concept color=popgold!80]{ node[font=\scriptsize\bfseries] {Colonnes $=$ images de $\vec{i}, \vec{j}$} }
  child[concept color=popgold!80]{ node[font=\scriptsize\bfseries] {Composée $\leftrightarrow$ produit ; points fixes $MX = X$} }
};
\end{tikzpicture}
\end{center}
\legende{Carte mentale de synthèse : les cinq piliers de la leçon et leurs liens. Le déterminant est le nœud central : il commande l'inversibilité, la résolution des systèmes et l'effet sur les aires.}
\end{popfigure}

Tu es maintenant au bout de la leçon. Retiens le fil conducteur : \textbf{une matrice est un langage} — elle traduit en nombres une transformation géométrique ou un système d'équations, et l'algèbre des matrices (produit, déterminant, inverse) traduit fidèlement la géométrie (composition, aires, réciprocité). Au Baccalauréat, les questions se succèdent presque toujours dans le même ordre : calculer, inverser, résoudre, interpréter. Si tu as compris chaque étape de cette section, tu peux aborder sereinement n'importe quel sujet sur les matrices. Entraîne-toi, refais les exemples sans regarder les solutions, et surtout : dessine toujours le carré unité et son image — une figure vaut mille calculs !

\newpage

\coursec{Activité d'intégration}

\courssub{Objectifs de l'activité}

À l'issue de cette activité d'intégration, tu seras capable de :

\begin{itemize}
\item modéliser une situation économique concrète (gestion de stock) à l'aide d'une matrice carrée d'ordre 2 ;
\item effectuer un produit matrice--vecteur pour calculer l'évolution d'un stock, et un produit de deux matrices pour composer deux transformations ;
\item calculer le déterminant d'une matrice $2\times 2$, justifier son inversibilité et déterminer sa matrice inverse ;
\item résoudre un système linéaire $2\times 2$ par la méthode de la matrice inverse, puis vérifier la solution ;
\item mettre en évidence, sur un exemple, que le produit matriciel n'est pas commutatif ($AB \neq BA$) et interpréter ce résultat dans le contexte.
\end{itemize}

Cette activité mobilise donc l'ensemble des savoir-faire de la leçon : opérations sur les matrices, déterminant, inverse, résolution de système et composition d'applications linéaires du plan.

\courssub{Situation}

\textbf{Le dépôt de boissons de Maman Ngo Bell.}

Maman Ngo Bell tient un dépôt de boissons bien connu au marché Mokolo à Yaoundé. Elle y vend essentiellement deux produits : des \textbf{casiers de jus} (notés $J$) et des \textbf{casiers d'eau minérale} (notés $E$). Chaque semaine, son stock évolue à cause des ventes, des retours de casiers vides et des réassorts partiels qu'elle reçoit de ses fournisseurs.

Après plusieurs mois d'observation, son neveu Étienne, élève en Terminale TI, a remarqué que le stock se transforme chaque semaine de façon régulière. Si l'on note $\begin{pmatrix} j \\ e \end{pmatrix}$ le stock un dimanche soir (nombre de casiers de jus et de casiers d'eau), alors le stock le dimanche suivant est $\begin{pmatrix} j' \\ e' \end{pmatrix}$ avec :
\[
\begin{pmatrix} j' \\ e' \end{pmatrix} = A \begin{pmatrix} j \\ e \end{pmatrix}, \qquad \text{où } A = \begin{pmatrix} 0,8 & 0,1 \\ 0,2 & 0,9 \end{pmatrix}.
\]
Autrement dit : chaque semaine, Maman Ngo Bell conserve $80\,\%$ de ses casiers de jus et récupère l'équivalent de $10\,\%$ du stock d'eau en jus (promotions, échanges avec une voisine commerçante), tandis que $20\,\%$ du stock de jus est converti en eau (réassort du fournisseur) et que $90\,\%$ du stock d'eau est conservé.

Le dimanche 1\textsuperscript{er} septembre, son stock est de $100$ casiers de jus et $50$ casiers d'eau, c'est-à-dire le vecteur $X_0 = \begin{pmatrix} 100 \\ 50 \end{pmatrix}$.

Un lundi matin, en ouvrant le dépôt, elle constate un stock de $55$ casiers de jus et $75$ casiers d'eau, mais elle a oublié de noter le stock du lundi précédent. Étienne veut l'aider à le retrouver.

Par ailleurs, un concurrent du quartier, M. Fotso, gère son propre dépôt selon une transformation différente, de matrice :
\[
B = \begin{pmatrix} 0,9 & 0 \\ 0,1 & 0,85 \end{pmatrix}.
\]
Les deux commerçants se demandent ce qui se passerait si un stock subissait d'abord la transformation de Maman Ngo Bell puis celle de M. Fotso, ou l'inverse : le résultat serait-il le même ?

\courssub{Consignes}

Travail en groupes de 3 ou 4 élèves, avec calculatrice. Chaque groupe rédige ses réponses sur une feuille, puis un rapporteur présente les résultats au tableau.

\textbf{Partie 1 : Évolution du stock.}
\begin{enumerate}
\item Calculer le stock $X_1$ de Maman Ngo Bell après une semaine, en effectuant le produit $A X_0$.
\item Calculer la matrice $A^2 = A \times A$, puis en déduire le stock $X_2$ après deux semaines de deux façons : en calculant $A X_1$ et en calculant $A^2 X_0$. Comparer les résultats.
\item Interpréter concrètement le coefficient $0,2$ de la matrice $A$ (deuxième ligne, première colonne).
\end{enumerate}

\textbf{Partie 2 : Retrouver le stock perdu.}
\begin{enumerate}
\setcounter{enumi}{3}
\item Traduire le problème de Maman Ngo Bell par une équation matricielle $A X = Y$, où $Y = \begin{pmatrix} 55 \\ 75 \end{pmatrix}$ est le stock constaté et $X$ le stock inconnu de la semaine précédente.
\item Calculer le déterminant de $A$. Justifier que $A$ est inversible et calculer $A^{-1}$.
\item En déduire $X$, puis vérifier le résultat en calculant $A X$.
\item Écrire le système linéaire de deux équations à deux inconnues associé à l'équation $A X = Y$, et retrouver la solution par la méthode de votre choix (substitution ou combinaison). Comparer avec la question 6.
\end{enumerate}

\textbf{Partie 3 : Composition des deux transformations.}
\begin{enumerate}
\setcounter{enumi}{7}
\item Calculer les matrices $AB$ et $BA$. Que constate-t-on ?
\item En partant du stock $X_0 = \begin{pmatrix} 100 \\ 50 \end{pmatrix}$, calculer le stock obtenu en appliquant d'abord la transformation de Maman Ngo Bell puis celle de M. Fotso, puis celui obtenu dans l'ordre inverse. Conclure sur l'importance de l'ordre des opérations.
\end{enumerate}

\courssub{Ressources à mobiliser}

\begin{aretenir}[Boîte à outils de la leçon]
\begin{itemize}
\item \textbf{Produit matrice--vecteur} : si $A = \begin{pmatrix} a & b \\ c & d \end{pmatrix}$ et $X = \begin{pmatrix} x \\ y \end{pmatrix}$, alors $A X = \begin{pmatrix} ax + by \\ cx + dy \end{pmatrix}$.
\item \textbf{Produit de deux matrices} : le coefficient de la ligne $i$, colonne $j$ de $AB$ s'obtient en multipliant la ligne $i$ de $A$ par la colonne $j$ de $B$.
\item \textbf{Déterminant} : $\det(A) = ad - bc$. La matrice $A$ est inversible si et seulement si $\det(A) \neq 0$.
\item \textbf{Inverse} : si $\det(A) \neq 0$, alors $A^{-1} = \dfrac{1}{ad - bc}\begin{pmatrix} d & -b \\ -c & a \end{pmatrix}$, et l'on a $A A^{-1} = A^{-1} A = I_2$.
\item \textbf{Système linéaire} : l'équation $A X = Y$ équivaut à un système linéaire $2\times 2$ ; si $A$ est inversible, l'unique solution est $X = A^{-1} Y$.
\item \textbf{Composition} : la matrice de la composée $g \circ f$ est le produit des matrices, dans l'ordre $B \times A$ si $f$ a pour matrice $A$ et $g$ pour matrice $B$. Le produit matriciel n'est pas commutatif en général.
\end{itemize}
\end{aretenir}

\courssub{Démarche et corrigé commenté}

\begin{exoresolu}[Le dépôt de boissons de Maman Ngo Bell]
On reprend la situation : $A = \begin{pmatrix} 0,8 & 0,1 \\ 0,2 & 0,9 \end{pmatrix}$, $X_0 = \begin{pmatrix} 100 \\ 50 \end{pmatrix}$, $Y = \begin{pmatrix} 55 \\ 75 \end{pmatrix}$, $B = \begin{pmatrix} 0,9 & 0 \\ 0,1 & 0,85 \end{pmatrix}$.
\tcblower
\textbf{Partie 1 : Évolution du stock.}

\textbf{1.} On calcule le produit matrice--vecteur :
\[
X_1 = A X_0 = \begin{pmatrix} 0,8 & 0,1 \\ 0,2 & 0,9 \end{pmatrix} \begin{pmatrix} 100 \\ 50 \end{pmatrix}
= \begin{pmatrix} 0,8 \times 100 + 0,1 \times 50 \\ 0,2 \times 100 + 0,9 \times 50 \end{pmatrix}
= \begin{pmatrix} 80 + 5 \\ 20 + 45 \end{pmatrix} = \begin{pmatrix} 85 \\ 65 \end{pmatrix}.
\]
Après une semaine, le dépôt compte $85$ casiers de jus et $65$ casiers d'eau.

\textbf{2.} Calculons d'abord $A^2$ :
\begin{align*}
A^2 &= \begin{pmatrix} 0,8 & 0,1 \\ 0,2 & 0,9 \end{pmatrix} \begin{pmatrix} 0,8 & 0,1 \\ 0,2 & 0,9 \end{pmatrix} \\
&= \begin{pmatrix} 0,8 \times 0,8 + 0,1 \times 0,2 & 0,8 \times 0,1 + 0,1 \times 0,9 \\ 0,2 \times 0,8 + 0,9 \times 0,2 & 0,2 \times 0,1 + 0,9 \times 0,9 \end{pmatrix}
= \begin{pmatrix} 0,66 & 0,17 \\ 0,34 & 0,83 \end{pmatrix}.
\end{align*}
Première méthode :
\[
X_2 = A X_1 = \begin{pmatrix} 0,8 & 0,1 \\ 0,2 & 0,9 \end{pmatrix} \begin{pmatrix} 85 \\ 65 \end{pmatrix}
= \begin{pmatrix} 68 + 6,5 \\ 17 + 58,5 \end{pmatrix} = \begin{pmatrix} 74,5 \\ 75,5 \end{pmatrix}.
\]
Deuxième méthode :
\[
X_2 = A^2 X_0 = \begin{pmatrix} 0,66 & 0,17 \\ 0,34 & 0,83 \end{pmatrix} \begin{pmatrix} 100 \\ 50 \end{pmatrix}
= \begin{pmatrix} 66 + 8,5 \\ 34 + 41,5 \end{pmatrix} = \begin{pmatrix} 74,5 \\ 75,5 \end{pmatrix}.
\]
Les deux calculs donnent le même résultat : c'est l'illustration de la relation $A(AX_0) = A^2 X_0$ (associativité du produit). Après deux semaines, on prévoit environ $74$ ou $75$ casiers de jus et $75$ ou $76$ casiers d'eau (on arrondit, un casier ne se coupe pas !).

\textbf{3.} Le coefficient $0,2$ (ligne 2, colonne 1) signifie que chaque semaine, $20\,\%$ du stock de jus contribue au stock d'eau : concrètement, Maman Ngo Bell revend une partie de ses casiers de jus et réassort en eau minérale, dont la demande augmente en saison sèche.

\textbf{Partie 2 : Retrouver le stock perdu.}

\textbf{4.} Le stock constaté $Y$ est l'image du stock inconnu $X$ par la transformation hebdomadaire : $A X = Y$, c'est-à-dire
\[
\begin{pmatrix} 0,8 & 0,1 \\ 0,2 & 0,9 \end{pmatrix} \begin{pmatrix} j \\ e \end{pmatrix} = \begin{pmatrix} 55 \\ 75 \end{pmatrix}.
\]

\textbf{5.} Déterminant :
\[
\det(A) = 0,8 \times 0,9 - 0,1 \times 0,2 = 0,72 - 0,02 = 0,7.
\]
Comme $\det(A) = 0,7 \neq 0$, la matrice $A$ est inversible et :
\[
A^{-1} = \frac{1}{0,7} \begin{pmatrix} 0,9 & -0,1 \\ -0,2 & 0,8 \end{pmatrix}.
\]

\textbf{6.} On en déduit :
\begin{align*}
X &= A^{-1} Y = \frac{1}{0,7} \begin{pmatrix} 0,9 & -0,1 \\ -0,2 & 0,8 \end{pmatrix} \begin{pmatrix} 55 \\ 75 \end{pmatrix} \\
&= \frac{1}{0,7} \begin{pmatrix} 0,9 \times 55 - 0,1 \times 75 \\ -0,2 \times 55 + 0,8 \times 75 \end{pmatrix}
= \frac{1}{0,7} \begin{pmatrix} 49,5 - 7,5 \\ -11 + 60 \end{pmatrix}
= \frac{1}{0,7} \begin{pmatrix} 42 \\ 49 \end{pmatrix}
= \begin{pmatrix} 60 \\ 70 \end{pmatrix}.
\end{align*}
Le stock de la semaine précédente était donc de $60$ casiers de jus et $70$ casiers d'eau. \textbf{Vérification} :
\[
A X = \begin{pmatrix} 0,8 & 0,1 \\ 0,2 & 0,9 \end{pmatrix} \begin{pmatrix} 60 \\ 70 \end{pmatrix}
= \begin{pmatrix} 48 + 7 \\ 12 + 63 \end{pmatrix} = \begin{pmatrix} 55 \\ 75 \end{pmatrix} = Y. \quad \checkmark
\]

\textbf{7.} Le système associé est $\begin{cases} 0,8 j + 0,1 e = 55 \\ 0,2 j + 0,9 e = 75 \end{cases}$. En multipliant la première équation par $9$ puis en soustrayant la deuxième : $7,2 j - 0,2 j = 495 - 75$, soit $7 j = 420$, d'où $j = 60$, puis $e = \dfrac{55 - 0,8 \times 60}{0,1} = 70$. On retrouve bien la même solution : la méthode matricielle est un raccourci puissant, surtout lorsqu'on doit résoudre plusieurs systèmes avec la même matrice $A$.

\textbf{Partie 3 : Composition des transformations.}

\textbf{8.}
\[
AB = \begin{pmatrix} 0,8 & 0,1 \\ 0,2 & 0,9 \end{pmatrix} \begin{pmatrix} 0,9 & 0 \\ 0,1 & 0,85 \end{pmatrix}
= \begin{pmatrix} 0,8\times0,9 + 0,1\times0,1 & 0,8\times0 + 0,1\times0,85 \\ 0,2\times0,9 + 0,9\times0,1 & 0,2\times0 + 0,9\times0,85 \end{pmatrix}
= \begin{pmatrix} 0,72 + 0,01 & 0 + 0,085 \\ 0,18 + 0,09 & 0 + 0,765 \end{pmatrix}
= \begin{pmatrix} 0,73 & 0,085 \\ 0,27 & 0,765 \end{pmatrix}.
\]
\[
BA = \begin{pmatrix} 0,9 & 0 \\ 0,1 & 0,85 \end{pmatrix} \begin{pmatrix} 0,8 & 0,1 \\ 0,2 & 0,9 \end{pmatrix}
= \begin{pmatrix} 0,72 & 0,09 \\ 0,08 + 0,17 & 0,01 + 0,765 \end{pmatrix}
= \begin{pmatrix} 0,72 & 0,09 \\ 0,25 & 0,775 \end{pmatrix}.
\]
On constate que $AB \neq BA$ : \textbf{le produit matriciel n'est pas commutatif}.

\textbf{9.} Appliquer d'abord $A$ puis $B$ revient à calculer $B(AX_0) = (BA)X_0$ :
\[
(BA) X_0 = \begin{pmatrix} 0,72 & 0,09 \\ 0,25 & 0,775 \end{pmatrix} \begin{pmatrix} 100 \\ 50 \end{pmatrix}
= \begin{pmatrix} 72 + 4,5 \\ 25 + 38,75 \end{pmatrix} = \begin{pmatrix} 76,5 \\ 63,75 \end{pmatrix}.
\]
Dans l'ordre inverse ($B$ puis $A$) :
\[
(AB) X_0 = \begin{pmatrix} 0,73 & 0,085 \\ 0,27 & 0,765 \end{pmatrix} \begin{pmatrix} 100 \\ 50 \end{pmatrix}
= \begin{pmatrix} 73 + 4,25 \\ 27 + 38,25 \end{pmatrix} = \begin{pmatrix} 77,25 \\ 65,25 \end{pmatrix}.
\]
Les stocks finaux sont différents : l'ordre dans lequel on applique les transformations a une réelle importance économique. Maman Ngo Bell et M. Fotso ne peuvent donc pas « échanger » l'ordre de leurs politiques de gestion sans modifier le résultat final.
\end{exoresolu}

\begin{attention}[Pièges fréquents rencontrés pendant l'activité]
\begin{itemize}
\item Dans le produit $A X_0$, on multiplie les \textbf{lignes} de $A$ par la \textbf{colonne} $X_0$, et non les colonnes entre elles.
\item Ne pas oublier le facteur $\dfrac{1}{\det(A)}$ dans le calcul de $A^{-1}$ : écrire $A^{-1} = \begin{pmatrix} 0,9 & -0,1 \\ -0,2 & 0,8 \end{pmatrix}$ est faux.
\item Dans la formule de l'inverse, on échange $a$ et $d$ et l'on change les signes de $b$ et $c$ : ne pas inverser les deux gestes.
\item Pour la composée « $A$ puis $B$ », la matrice est $BA$ et non $AB$ : la matrice de droite agit en premier.
\item Vérifier chaque coefficient du produit de matrices : une erreur de signe ou d'addition (comme $0,2\times0 + 0,9\times0,85 = 0,765$ et non $0,785$) change le résultat final.
\end{itemize}
\end{attention}

\courssub{Bilan de l'activité}

Cette activité a permis de mettre en évidence les points essentiels de la leçon :

\begin{itemize}
\item Une \textbf{matrice $2\times 2$} est un outil efficace pour modéliser une transformation régulière d'une grandeur à deux composantes (ici, un stock) : le produit matrice--vecteur code l'évolution hebdomadaire.
\item La \textbf{puissance $A^2$} permet de prévoir l'évolution sur deux périodes en un seul calcul, grâce à l'associativité du produit matriciel.
\item Le \textbf{déterminant} joue le rôle de test d'inversibilité : $\det(A) = 0,7 \neq 0$ garantissait que le problème « remonter le temps » avait une unique solution.
\item La \textbf{matrice inverse} transforme la résolution d'un système linéaire en un simple produit matriciel $X = A^{-1} Y$ : c'est la méthode la plus rapide dès que la matrice est connue.
\item Enfin, l'activité a montré de façon concrète que \textbf{le produit matriciel n'est pas commutatif} : $AB \neq BA$, et cette non-commutativité a une interprétation économique claire — l'ordre des décisions de gestion modifie le stock final.
\end{itemize}

\begin{savaistu}[Des matrices partout autour de nous]
Les matrices comme celle de Maman Ngo Bell, dont les coefficients représentent des proportions transférées d'un état à un autre, sont appelées \emph{matrices de transition}. Elles sont au cœur des \emph{chaînes de Markov}, utilisées pour prévoir l'évolution des parts de marché entre opérateurs téléphoniques (MTN, Orange, Camtel), modéliser la propagation d'une épidémie ou classer les pages web par le moteur de recherche Google. Un simple tableau de quatre nombres peut ainsi décrire des phénomènes économiques ou sociaux de grande ampleur !
\end{savaistu}

\newpage

\coursec{Méthodes et exercices résolus}

\coursec{Matrices 2×2 et applications linéaires du plan}

\emph{Famille de situations : Représentations et transformations des configurations planes dans l'environnement.}  
Cette leçon constitue l'outil algébrique fondamental pour modéliser et résoudre les problèmes de géométrie plane (rotations, homothéties, symétries, similitudes) et pour traiter les systèmes linéaires 2×2. Elle s'inscrit dans la continuité de l'étude des vecteurs et des applications linéaires en Première.

\courssub{Matrices 2×2 : définition et opérations}

\begin{definition}[Matrice carrée d'ordre 2]
Une \cle{matrice carrée d'ordre 2} est un tableau carré à 2 lignes et 2 colonnes de nombres réels :
\[ A=\begin{pmatrix} a & b \\ c & d \end{pmatrix} \quad \text{avec } a,b,c,d\in\mathbb{R}. \]
L'ensemble de ces matrices est noté $M_2(\mathbb{R})$. Deux matrices sont égales si et seulement si leurs coefficients respectifs sont égaux.
\end{definition}

\begin{propriete}[Opérations dans $M_2(\mathbb{R})$]
Soient $A=\begin{pmatrix} a & b \\ c & d \end{pmatrix}$, $B=\begin{pmatrix} a' & b' \\ c' & d' \end{pmatrix}$ et $k\in\mathbb{R}$.
\begin{enumerate}
\item \textbf{Somme} : $A+B=\begin{pmatrix} a+a' & b+b' \\ c+c' & d+d' \end{pmatrix}$.
\item \textbf{Produit par un scalaire} : $kA=\begin{pmatrix} ka & kb \\ kc & kd \end{pmatrix}$.
\item \textbf{Produit matriciel} : $AB=\begin{pmatrix} aa'+bc' & ab'+bd' \\ ca'+dc' & cb'+dd' \end{pmatrix}$.
  La règle est \cle{ligne $\times$ colonne} : le coefficient $(i,j)$ de $AB$ est le produit scalaire de la ligne $i$ de $A$ par la colonne $j$ de $B$.
\end{enumerate}
\textbf{Propriétés algébriques} : La somme et le produit sont associatifs ; le produit est distributif par rapport à la somme ; le produit n'est \textbf{pas commutatif} en général ($AB\neq BA$) ; $I_2=\begin{pmatrix} 1 & 0 \\ 0 & 1 \end{pmatrix}$ est l'élément neutre du produit.
\end{propriete}

\begin{methode}[Calculer avec les matrices $2\times 2$]
Pour manipuler des matrices $A=\begin{pmatrix} a & b \\ c & d \end{pmatrix}$ et $B=\begin{pmatrix} a' & b' \\ c' & d' \end{pmatrix}$ :
\begin{enumerate}
\item \textbf{Somme} : on additionne \cle{coefficient par coefficient} : $A+B=\begin{pmatrix} a+a' & b+b' \\ c+c' & d+d' \end{pmatrix}$.
\item \textbf{Produit par un réel} $k$ : on multiplie \emph{tous} les coefficients par $k$ : $kA=\begin{pmatrix} ka & kb \\ kc & kd \end{pmatrix}$.
\item \textbf{Produit de deux matrices} : règle \cle{ligne $\times$ colonne}. Le coefficient en ligne $i$, colonne $j$ de $AB$ s'obtient en multipliant la ligne $i$ de $A$ par la colonne $j$ de $B$ :
\[ AB=\begin{pmatrix} aa'+bc' & ab'+bd' \\ ca'+dc' & cb'+dd' \end{pmatrix}. \]
\item \textbf{Vérifier} : le produit $AB$ n'est \emph{pas} commutatif en général : $AB\neq BA$. Ne jamais conclure sans calculer les deux.
\end{enumerate}
\textbf{Réflexe} : pour le produit, pose ton doigt sur la ligne de $A$ et fais-le glisser le long de la colonne de $B$ en multipliant terme à terme.
\end{methode}

\begin{exoresolu}[Opérations sur les matrices]
On considère les matrices $A=\begin{pmatrix} 2 & -1 \\ 3 & 4 \end{pmatrix}$ et $B=\begin{pmatrix} 1 & 5 \\ -2 & 0 \end{pmatrix}$.
\begin{enumerate}
\item Calculer $A+B$ et $3A-2B$.
\item Calculer $AB$ et $BA$. Que constate-t-on ?
\item Calculer $A^2$.
\end{enumerate}
\tcblower
\textbf{1.} La somme se fait coefficient par coefficient :
\[ A+B=\begin{pmatrix} 2+1 & -1+5 \\ 3+(-2) & 4+0 \end{pmatrix}=\begin{pmatrix} 3 & 4 \\ 1 & 4 \end{pmatrix}. \]
Pour $3A-2B$, on multiplie d'abord chaque matrice par son scalaire :
\[ 3A=\begin{pmatrix} 6 & -3 \\ 9 & 12 \end{pmatrix}, \qquad 2B=\begin{pmatrix} 2 & 10 \\ -4 & 0 \end{pmatrix}, \]
\[ 3A-2B=\begin{pmatrix} 6-2 & -3-10 \\ 9-(-4) & 12-0 \end{pmatrix}=\begin{pmatrix} 4 & -13 \\ 13 & 12 \end{pmatrix}. \]

\textbf{2.} Produit $AB$ par la règle ligne $\times$ colonne :
\begin{align*}
AB&=\begin{pmatrix} 2 & -1 \\ 3 & 4 \end{pmatrix}\begin{pmatrix} 1 & 5 \\ -2 & 0 \end{pmatrix}\\
&=\begin{pmatrix} 2\times 1+(-1)\times(-2) & 2\times 5+(-1)\times 0 \\ 3\times 1+4\times(-2) & 3\times 5+4\times 0 \end{pmatrix}
=\begin{pmatrix} 4 & 10 \\ -5 & 15 \end{pmatrix}.
\end{align*}
Produit $BA$ :
\[ BA=\begin{pmatrix} 1\times 2+5\times 3 & 1\times(-1)+5\times 4 \\ -2\times 2+0\times 3 & -2\times(-1)+0\times 4 \end{pmatrix}=\begin{pmatrix} 17 & 19 \\ -4 & 2 \end{pmatrix}. \]
\textbf{Constatation} : $AB\neq BA$ : le produit matriciel n'est pas commutatif.

\textbf{3.} On calcule $A^2=A\times A$ :
\[ A^2=\begin{pmatrix} 2\times 2+(-1)\times 3 & 2\times(-1)+(-1)\times 4 \\ 3\times 2+4\times 3 & 3\times(-1)+4\times 4 \end{pmatrix}=\begin{pmatrix} 1 & -6 \\ 18 & 13 \end{pmatrix}. \]
\textbf{Conclusion} : $A+B=\begin{pmatrix} 3 & 4 \\ 1 & 4 \end{pmatrix}$, $3A-2B=\begin{pmatrix} 4 & -13 \\ 13 & 12 \end{pmatrix}$, $AB=\begin{pmatrix} 4 & 10 \\ -5 & 15 \end{pmatrix}$, $BA=\begin{pmatrix} 17 & 19 \\ -4 & 2 \end{pmatrix}$ et $A^2=\begin{pmatrix} 1 & -6 \\ 18 & 13 \end{pmatrix}$.
\end{exoresolu}

\courssub{Déterminant et inverse d'une matrice 2×2}

\begin{definition}[Déterminant d'une matrice 2×2]
Soit $A=\begin{pmatrix} a & b \\ c & d \end{pmatrix}\in M_2(\mathbb{R})$. Le \cle{déterminant} de $A$ est le réel noté $\det(A)$ ou $|A|$ défini par :
\[ \det(A)=ad-bc. \]
C'est le produit des coefficients de la diagonale principale \emph{moins} le produit des coefficients de l'autre diagonale.
\end{definition}

\begin{propriete}[Inversibilité et matrice inverse]
\begin{enumerate}
\item Une matrice $A\in M_2(\mathbb{R})$ est \cle{inversible} (ou \cle{régulière}) si et seulement si $\det(A)\neq 0$.
\item Si $\det(A)\neq 0$, l'\cle{inverse} de $A$ est l'unique matrice $A^{-1}$ telle que $AA^{-1}=A^{-1}A=I_2$. Elle est donnée par la formule :
\[ A^{-1}=\frac{1}{\det(A)}\begin{pmatrix} d & -b \\ -c & a \end{pmatrix}. \]
\item \textbf{Propriétés utiles (démonstration exigible au Bac)} :
  $\det(AB)=\det(A)\det(B)$ ; $\det(A^{-1})=\frac{1}{\det(A)}$ ; $(AB)^{-1}=B^{-1}A^{-1}$ ; $(A^{-1})^{-1}=A$.
\end{enumerate}
\end{propriete}

\begin{methode}[Déterminant et inverse]
Soit $A=\begin{pmatrix} a & b \\ c & d \end{pmatrix}$.
\begin{enumerate}
\item \textbf{Calculer le déterminant} : $\det(A)=ad-bc$.
\item \textbf{Tester l'inversibilité} : $A$ est inversible \cle{si et seulement si} $\det(A)\neq 0$. Cette vérification est \emph{obligatoire} avant tout calcul d'inverse.
\item \textbf{Appliquer la formule de l'inverse} : si $\det(A)\neq 0$,
\[ A^{-1}=\frac{1}{\det(A)}\begin{pmatrix} d & -b \\ -c & a \end{pmatrix}. \]
On échange $a$ et $d$, on change les signes de $b$ et $c$, et on divise tout par le déterminant.
\item \textbf{Contrôler} : vérifier que $A\times A^{-1}=I_2=\begin{pmatrix} 1 & 0 \\ 0 & 1 \end{pmatrix}$.
\end{enumerate}
\end{methode}

\begin{exoresolu}[Déterminant et inverse]
On considère les matrices $M=\begin{pmatrix} 3 & 2 \\ 1 & 4 \end{pmatrix}$ et $N=\begin{pmatrix} 2 & -4 \\ -1 & 2 \end{pmatrix}$.
\begin{enumerate}
\item Calculer $\det(M)$ et $\det(N)$.
\item Les matrices $M$ et $N$ sont-elles inversibles ? Justifier.
\item Le cas échéant, déterminer l'inverse et vérifier le résultat.
\end{enumerate}
\tcblower
\textbf{1.} Par définition du déterminant :
\[ \det(M)=3\times 4-2\times 1=12-2=10 ; \qquad \det(N)=2\times 2-(-4)\times(-1)=4-4=0. \]

\textbf{2.} Une matrice $2\times 2$ est inversible si et seulement si son déterminant est non nul.
\begin{itemize}
\item $\det(M)=10\neq 0$, donc $M$ est \cle{inversible}.
\item $\det(N)=0$, donc $N$ n'est \textbf{pas} inversible.
\end{itemize}

\textbf{3.} Comme $\det(M)=10$, la formule de l'inverse donne :
\[ M^{-1}=\frac{1}{10}\begin{pmatrix} 4 & -2 \\ -1 & 3 \end{pmatrix}=\begin{pmatrix} \dfrac{2}{5} & -\dfrac{1}{5} \\[6pt] -\dfrac{1}{10} & \dfrac{3}{10} \end{pmatrix}. \]
\textbf{Vérification} :
\[ M\times M^{-1}=\frac{1}{10}\begin{pmatrix} 3 & 2 \\ 1 & 4 \end{pmatrix}\begin{pmatrix} 4 & -2 \\ -1 & 3 \end{pmatrix}=\frac{1}{10}\begin{pmatrix} 12-2 & -6+6 \\ 4-4 & -2+12 \end{pmatrix}=\frac{1}{10}\begin{pmatrix} 10 & 0 \\ 0 & 10 \end{pmatrix}=I_2. \]
\textbf{Conclusion} : $M$ est inversible et $M^{-1}=\dfrac{1}{10}\begin{pmatrix} 4 & -2 \\ -1 & 3 \end{pmatrix}$ ; $N$ n'est pas inversible car $\det(N)=0$.
\end{exoresolu}

\courssub{Matrices et systèmes linéaires 2×2}

\begin{propriete}[Écriture matricielle et résolution]
Tout système linéaire à deux inconnues
\[ \begin{cases} ax+by=e \\ cx+dy=f \end{cases} \]
s'écrit sous la forme matricielle $AX=B$ avec
$A=\begin{pmatrix} a & b \\ c & d \end{pmatrix}$, $X=\begin{pmatrix} x \\ y \end{pmatrix}$, $B=\begin{pmatrix} e \\ f \end{pmatrix}$.
\begin{itemize}
\item Si $\det(A)\neq 0$, $A$ est inversible et le système admet un \cle{unique couple solution} donné par $X=A^{-1}B$.
\item Si $\det(A)=0$, le système n'a pas de solution unique (soit aucune, soit une infinité) ; on ne peut pas utiliser la matrice inverse, il faut raisonner par combinaison linéaire (méthode de substitution ou d'élimination).
\end{itemize}
\end{propriete}

\begin{methode}[Système linéaire et matrice inverse]
Pour résoudre le système $\begin{cases} ax+by=e \\ cx+dy=f \end{cases}$ :
\begin{enumerate}
\item \textbf{Écrire le système sous forme matricielle} $AX=B$ avec
$A=\begin{pmatrix} a & b \\ c & d \end{pmatrix}$, $X=\begin{pmatrix} x \\ y \end{pmatrix}$, $B=\begin{pmatrix} e \\ f \end{pmatrix}$. Attention à l'\cle{ordre des inconnues} et aux signes des coefficients.
\item \textbf{Calculer $\det(A)$}. Si $\det(A)\neq 0$, le système admet un \cle{unique couple solution}.
\item \textbf{Calculer $A^{-1}$} puis conclure : $X=A^{-1}B$.
\item \textbf{Vérifier} en substituant la solution trouvée dans les deux équations de départ.
\end{enumerate}
\textbf{Réflexe} : si $\det(A)=0$, la méthode échoue : le système a soit aucune solution, soit une infinité ; on revient alors à la combinaison linéaire.
\end{methode}

\begin{exoresolu}[Résolution d'un système par la matrice inverse]
Un commerçant du marché Mokolo vend deux types de paniers. Le premier jour, la vente de $3$ paniers de type A et $2$ paniers de type B rapporte $\num{17000}$ FCFA. Le second jour, la vente de $2$ paniers de type A et $5$ paniers de type B rapporte $\num{26000}$ FCFA.
\begin{enumerate}
\item Modéliser la situation par un système de deux équations à deux inconnues.
\item Écrire ce système sous forme matricielle $AX=B$.
\item Résoudre le système à l'aide de la matrice inverse et conclure.
\end{enumerate}
\tcblower
\textbf{1.} Soit $x$ le prix (en milliers de FCFA) d'un panier de type A et $y$ celui d'un panier de type B. Les données se traduisent par :
\[ \begin{cases} 3x+2y=17 \\ 2x+5y=26 \end{cases} \]

\textbf{2.} Forme matricielle : $AX=B$ avec
\[ A=\begin{pmatrix} 3 & 2 \\ 2 & 5 \end{pmatrix}, \qquad X=\begin{pmatrix} x \\ y \end{pmatrix}, \qquad B=\begin{pmatrix} 17 \\ 26 \end{pmatrix}. \]

\textbf{3.} \textbf{Étape 1 : déterminant.}
\[ \det(A)=3\times 5-2\times 2=15-4=11\neq 0. \]
Donc $A$ est inversible et le système admet un unique couple solution.

\textbf{Étape 2 : inverse.}
\[ A^{-1}=\frac{1}{11}\begin{pmatrix} 5 & -2 \\ -2 & 3 \end{pmatrix}. \]

\textbf{Étape 3 : résolution.} De $AX=B$ on tire $X=A^{-1}B$ :
\[ X=\frac{1}{11}\begin{pmatrix} 5 & -2 \\ -2 & 3 \end{pmatrix}\begin{pmatrix} 17 \\ 26 \end{pmatrix}=\frac{1}{11}\begin{pmatrix} 5\times 17-2\times 26 \\ -2\times 17+3\times 26 \end{pmatrix}=\frac{1}{11}\begin{pmatrix} 85-52 \\ -34+78 \end{pmatrix}=\frac{1}{11}\begin{pmatrix} 33 \\ 44 \end{pmatrix}=\begin{pmatrix} 3 \\ 4 \end{pmatrix}. \]

\textbf{Vérification} : $3\times 3+2\times 4=9+8=17$ \checkmark{} et $2\times 3+5\times 4=6+20=26$ \checkmark.

\textbf{Conclusion} : le panier de type A coûte $\num{3000}$ FCFA et le panier de type B coûte $\num{4000}$ FCFA.
\end{exoresolu}

\courssub{Matrice d'une application linéaire du plan}

\begin{definition}[Application linéaire du plan]
Une application $f : \mathbb{R}^2 \to \mathbb{R}^2$ est \cle{linéaire} si pour tous vecteurs $\vec{u},\vec{v}$ et tout réel $k$ :
\[ f(\vec{u}+\vec{v})=f(\vec{u})+f(\vec{v}) \quad \text{et} \quad f(k\vec{u})=k\,f(\vec{u}). \]
Equivalentement : $f(a\vec{u}+b\vec{v})=a f(\vec{u})+b f(\vec{v})$.
\end{definition}

\begin{propriete}[Matrice d'une application linéaire dans une base]
Soit $f$ une application linéaire du plan et $\mathcal{B}=(\vec{\imath},\vec{\jmath})$ une base. Les images des vecteurs de la base s'écrivent de façon unique :
\[ f(\vec{\imath})=a\vec{\imath}+c\vec{\jmath}, \qquad f(\vec{\jmath})=b\vec{\imath}+d\vec{\jmath}. \]
La \cle{matrice de $f$ dans la base $\mathcal{B}$} est la matrice $M=\begin{pmatrix} a & b \\ c & d \end{pmatrix}$ dont les \cle{colonnes} sont les coordonnées de $f(\vec{\imath})$ et $f(\vec{\jmath})$.
Pour tout vecteur $\vec{u}=x\vec{\imath}+y\vec{\jmath}$ de coordonnées $\begin{pmatrix} x \\ y \end{pmatrix}$, les coordonnées de $f(\vec{u})$ sont obtenues par le produit matriciel :
\[ \begin{pmatrix} x' \\ y' \end{pmatrix}=M\begin{pmatrix} x \\ y \end{pmatrix}. \]
\end{propriete}

\begin{methode}[Matrice d'une application linéaire]
Soit $f$ une application linéaire du plan et $\mathcal{B}=(\vec{\imath},\vec{\jmath})$ une base. Pour construire la matrice $M$ de $f$ dans $\mathcal{B}$ :
\begin{enumerate}
\item \textbf{Calculer les images} $f(\vec{\imath})$ et $f(\vec{\jmath})$.
\item \textbf{Décomposer} chaque image dans la base $\mathcal{B}$ : $f(\vec{\imath})=a\vec{\imath}+c\vec{\jmath}$ et $f(\vec{\jmath})=b\vec{\imath}+d\vec{\jmath}$.
\item \textbf{Placer les coordonnées en colonnes} : la première colonne de $M$ est formée des coordonnées de $f(\vec{\imath})$, la deuxième de celles de $f(\vec{\jmath})$ :
\[ M=\begin{pmatrix} a & b \\ c & d \end{pmatrix}. \]
\item \textbf{Image d'un vecteur} : si $\vec{u}$ a pour coordonnées $\begin{pmatrix} x \\ y \end{pmatrix}$, alors les coordonnées de $f(\vec{u})$ sont données par le produit $M\begin{pmatrix} x \\ y \end{pmatrix}$.
\end{enumerate}
\textbf{Piège classique} : les coordonnées des images se rangent en \cle{colonnes}, jamais en lignes.
\end{methode}

\begin{exoresolu}[Matrice d'une application linéaire et image d'un vecteur]
Le plan est muni d'une base $\mathcal{B}=(\vec{\imath},\vec{\jmath})$. Soit $f$ l'application linéaire définie par :
\[ f(\vec{\imath})=2\vec{\imath}-\vec{\jmath} \qquad \text{et} \qquad f(\vec{\jmath})=\vec{\imath}+3\vec{\jmath}. \]
\begin{enumerate}
\item Écrire la matrice $M$ de $f$ dans la base $\mathcal{B}$.
\item Déterminer les coordonnées de $f(\vec{u})$ où $\vec{u}=-2\vec{\imath}+5\vec{\jmath}$.
\item Soit $\vec{v}=x\vec{\imath}+y\vec{\jmath}$. Exprimer les coordonnées $(x';y')$ de $f(\vec{v})$ en fonction de $x$ et $y$.
\end{enumerate}
\tcblower
\textbf{1.} Les coordonnées de $f(\vec{\imath})=2\vec{\imath}-\vec{\jmath}$ sont $\begin{pmatrix} 2 \\ -1 \end{pmatrix}$ : c'est la \emph{première colonne}. Celles de $f(\vec{\jmath})=\vec{\imath}+3\vec{\jmath}$ sont $\begin{pmatrix} 1 \\ 3 \end{pmatrix}$ : c'est la \emph{deuxième colonne}. D'où :
\[ M=\begin{pmatrix} 2 & 1 \\ -1 & 3 \end{pmatrix}. \]

\textbf{2.} Les coordonnées de $\vec{u}$ sont $\begin{pmatrix} -2 \\ 5 \end{pmatrix}$. Par produit matriciel :
\[ M\begin{pmatrix} -2 \\ 5 \end{pmatrix}=\begin{pmatrix} 2 & 1 \\ -1 & 3 \end{pmatrix}\begin{pmatrix} -2 \\ 5 \end{pmatrix}=\begin{pmatrix} 2\times(-2)+1\times 5 \\ -1\times(-2)+3\times 5 \end{pmatrix}=\begin{pmatrix} -4+5 \\ 2+15 \end{pmatrix}=\begin{pmatrix} 1 \\ 17 \end{pmatrix}. \]
Donc $f(\vec{u})=\vec{\imath}+17\vec{\jmath}$.

\textbf{3.} Pour $\vec{v}$ de coordonnées $\begin{pmatrix} x \\ y \end{pmatrix}$ :
\[ \begin{pmatrix} x' \\ y' \end{pmatrix}=M\begin{pmatrix} x \\ y \end{pmatrix}=\begin{pmatrix} 2 & 1 \\ -1 & 3 \end{pmatrix}\begin{pmatrix} x \\ y \end{pmatrix}=\begin{pmatrix} 2x+y \\ -x+3y \end{pmatrix}. \]
\textbf{Conclusion} : $M=\begin{pmatrix} 2 & 1 \\ -1 & 3 \end{pmatrix}$, $f(\vec{u})$ a pour coordonnées $(1\,;\,17)$, et pour tout vecteur $\vec{v}(x;y)$, $f(\vec{v})$ a pour coordonnées $(2x+y\,;\,-x+3y)$.
\end{exoresolu}

\courssub{Opérations sur les applications linéaires et automorphismes}

\begin{propriete}[Matrices de la somme, de la composée et de la réciproque]
Soient $f$ et $g$ deux applications linéaires du plan, de matrices respectives $A$ et $B$ dans une même base $\mathcal{B}$.
\begin{enumerate}
\item \textbf{Somme} : la matrice de $f+g$ est $A+B$.
\item \textbf{Composée} : la matrice de $g\circ f$ est le produit $B\times A$. \cle{Attention à l'ordre} : l'application écrite à droite ($f$) agit en premier, sa matrice est donc à droite dans le produit.
\item \textbf{Réciproque} : $f$ est un \cle{automorphisme} (application linéaire bijective) si et seulement si $\det(A)\neq 0$ ; la matrice de $f^{-1}$ est alors $A^{-1}$.
\item \textbf{Contrôle} : pour une composée, vérifier que $A\times B\neq B\times A$ en général ; pour une réciproque, vérifier $A\times A^{-1}=I_2$.
\end{enumerate}
\end{propriete}

\begin{methode}[Opérations sur les applications linéaires]
Soient $f$ et $g$ deux applications linéaires du plan, de matrices respectives $A$ et $B$ dans une même base $\mathcal{B}$.
\begin{enumerate}
\item \textbf{Somme} : la matrice de $f+g$ est $A+B$.
\item \textbf{Composée} : la matrice de $g\circ f$ est le produit $B\times A$. \cle{Attention à l'ordre} : l'application écrite à droite ($f$) agit en premier, sa matrice est donc à droite dans le produit.
\item \textbf{Réciproque} : $f$ est un automorphisme (bijectif) si et seulement si $\det(A)\neq 0$ ; la matrice de $f^{-1}$ est alors $A^{-1}$.
\item \textbf{Contrôle} : pour une composée, vérifier que $A\times B\neq B\times A$ en général ; pour une réciproque, vérifier $A\times A^{-1}=I_2$.
\end{enumerate}
\end{methode}

\begin{exoresolu}[Composée et automorphisme réciproque]
Dans une base $\mathcal{B}=(\vec{\imath},\vec{\jmath})$, les applications linéaires $f$ et $g$ ont pour matrices :
\[ A=\begin{pmatrix} 1 & 2 \\ 0 & 3 \end{pmatrix} \quad \text{(matrice de } f\text{)} \qquad B=\begin{pmatrix} 2 & -1 \\ 1 & 1 \end{pmatrix} \quad \text{(matrice de } g\text{)}. \]
\begin{enumerate}
\item Déterminer la matrice de $g\circ f$ dans $\mathcal{B}$.
\item Montrer que $f$ est un automorphisme et déterminer la matrice de $f^{-1}$.
\item Soit $\vec{u}$ de coordonnées $(1\,;\,1)$. Calculer les coordonnées de $(g\circ f)(\vec{u})$ de deux manières différentes.
\end{enumerate}
\tcblower
\textbf{1.} La matrice de $g\circ f$ est $B\times A$ (car $f$ agit en premier) :
\[ BA=\begin{pmatrix} 2 & -1 \\ 1 & 1 \end{pmatrix}\begin{pmatrix} 1 & 2 \\ 0 & 3 \end{pmatrix}=\begin{pmatrix} 2\times 1+(-1)\times 0 & 2\times 2+(-1)\times 3 \\ 1\times 1+1\times 0 & 1\times 2+1\times 3 \end{pmatrix}=\begin{pmatrix} 2 & 1 \\ 1 & 5 \end{pmatrix}. \]

\textbf{2.} $\det(A)=1\times 3-2\times 0=3\neq 0$, donc $A$ est inversible et $f$ est un \cle{automorphisme} du plan. La matrice de $f^{-1}$ est :
\[ A^{-1}=\frac{1}{3}\begin{pmatrix} 3 & -2 \\ 0 & 1 \end{pmatrix}. \]
Vérification : $AA^{-1}=\dfrac{1}{3}\begin{pmatrix} 1 & 2 \\ 0 & 3 \end{pmatrix}\begin{pmatrix} 3 & -2 \\ 0 & 1 \end{pmatrix}=\dfrac{1}{3}\begin{pmatrix} 3 & 0 \\ 0 & 3 \end{pmatrix}=I_2$ \checkmark.

\textbf{3.} \emph{Première manière} : avec la matrice de $g\circ f$ :
\[ BA\begin{pmatrix} 1 \\ 1 \end{pmatrix}=\begin{pmatrix} 2 & 1 \\ 1 & 5 \end{pmatrix}\begin{pmatrix} 1 \\ 1 \end{pmatrix}=\begin{pmatrix} 3 \\ 6 \end{pmatrix}. \]
\emph{Deuxième manière} : en appliquant successivement $f$ puis $g$ :
\[ A\begin{pmatrix} 1 \\ 1 \end{pmatrix}=\begin{pmatrix} 1 & 2 \\ 0 & 3 \end{pmatrix}\begin{pmatrix} 1 \\ 1 \end{pmatrix}=\begin{pmatrix} 3 \\ 3 \end{pmatrix}, \qquad B\begin{pmatrix} 3 \\ 3 \end{pmatrix}=\begin{pmatrix} 2 & -1 \\ 1 & 1 \end{pmatrix}\begin{pmatrix} 3 \\ 3 \end{pmatrix}=\begin{pmatrix} 3 \\ 6 \end{pmatrix}. \]
Les deux méthodes donnent le même résultat, ce qui conforte le calcul.
\textbf{Conclusion} : $\mathrm{Mat}(g\circ f)=\begin{pmatrix} 2 & 1 \\ 1 & 5 \end{pmatrix}$, $\mathrm{Mat}(f^{-1})=\dfrac{1}{3}\begin{pmatrix} 3 & -2 \\ 0 & 1 \end{pmatrix}$, et $(g\circ f)(\vec{u})$ a pour coordonnées $(3\,;\,6)$.
\end{exoresolu}

\courssub{Applications géométriques et synthèse}

\begin{aretenir}[Matrices des transformations usuelles (base orthonormée)]
Dans un repère orthonormé $(O;\vec{\imath},\vec{\jmath})$, les transformations planes suivantes sont des applications linéaires dont les matrices sont :
\begin{center}
\begin{tabular}{|l|cc|}
\hline
\textbf{Transformation} & \textbf{Matrice} \\ \hline
Symétrie orthogonale d'axe $(Ox)$ & $\begin{pmatrix} 1 & 0 \\ 0 & -1 \end{pmatrix}$ \\ \hline
Symétrie orthogonale d'axe $(Oy)$ & $\begin{pmatrix} -1 & 0 \\ 0 & 1 \end{pmatrix}$ \\ \hline
Symétrie centrale de centre $O$ & $-I_2=\begin{pmatrix} -1 & 0 \\ 0 & -1 \end{pmatrix}$ \\ \hline
Homothétie de centre $O$, rapport $k$ & $kI_2=\begin{pmatrix} k & 0 \\ 0 & k \end{pmatrix}$ \\ \hline
Rotation de centre $O$, angle $\theta$ & $\begin{pmatrix} \cos\theta & -\sin\theta \\ \sin\theta & \cos\theta \end{pmatrix}$ \\ \hline
\end{tabular}
\end{center}
\textbf{Composée} : la matrice de la composée de deux transformations est le produit de leurs matrices (dans l'ordre inverse de l'application). Une \cle{similitude} (composée d'une rotation et d'une homothétie de même centre) a pour matrice $\begin{pmatrix} a & -b \\ b & a \end{pmatrix}$ avec $a,b\in\mathbb{R}$.
\end{aretenir}

\begin{methode}[Matrices et transformations usuelles]
Dans la base orthonormée $(\vec{\imath},\vec{\jmath})$, reconnaître et utiliser les matrices des transformations de référence :
\begin{enumerate}
\item \textbf{Symétrie orthogonale d'axe des abscisses} : $\begin{pmatrix} 1 & 0 \\ 0 & -1 \end{pmatrix}$ ; \textbf{d'axe des ordonnées} : $\begin{pmatrix} -1 & 0 \\ 0 & 1 \end{pmatrix}$.
\item \textbf{Symétrie centrale de centre $O$} : $-I_2=\begin{pmatrix} -1 & 0 \\ 0 & -1 \end{pmatrix}$.
\item \textbf{Homothétie de centre $O$ et de rapport $k$} : $kI_2=\begin{pmatrix} k & 0 \\ 0 & k \end{pmatrix}$.
\item \textbf{Rotation de centre $O$ et d'angle $\theta$} : $\begin{pmatrix} \cos\theta & -\sin\theta \\ \sin\theta & \cos\theta \end{pmatrix}$.
\item Pour trouver l'image d'un point $M(x;y)$, calculer le produit de la matrice par $\begin{pmatrix} x \\ y \end{pmatrix}$ ; pour composer deux transformations, multiplier leurs matrices dans le bon ordre.
\end{enumerate}
\end{methode}

\begin{exoresolu}[Rotation et homothétie]
Le plan est muni d'un repère orthonormé $(O\,;\vec{\imath},\vec{\jmath})$. On considère la rotation $r$ de centre $O$ et d'angle $\dfrac{\pi}{2}$, et l'homothétie $h$ de centre $O$ et de rapport $2$.
\begin{enumerate}
\item Écrire les matrices $R$ et $H$ de $r$ et $h$ dans la base $(\vec{\imath},\vec{\jmath})$.
\item Déterminer la matrice de $t=h\circ r$ et caractériser $t$.
\item Déterminer les coordonnées de l'image par $t$ du point $A(3\,;\,-1)$.
\end{enumerate}
\tcblower
\textbf{1.} Pour la rotation d'angle $\theta=\dfrac{\pi}{2}$ : $\cos\dfrac{\pi}{2}=0$ et $\sin\dfrac{\pi}{2}=1$, donc
\[ R=\begin{pmatrix} \cos\frac{\pi}{2} & -\sin\frac{\pi}{2} \\ \sin\frac{\pi}{2} & \cos\frac{\pi}{2} \end{pmatrix}=\begin{pmatrix} 0 & -1 \\ 1 & 0 \end{pmatrix}, \qquad H=2I_2=\begin{pmatrix} 2 & 0 \\ 0 & 2 \end{pmatrix}. \]

\textbf{2.} La matrice de $t=h\circ r$ est $H\times R$ :
\[ HR=\begin{pmatrix} 2 & 0 \\ 0 & 2 \end{pmatrix}\begin{pmatrix} 0 & -1 \\ 1 & 0 \end{pmatrix}=\begin{pmatrix} 0 & -2 \\ 2 & 0 \end{pmatrix}. \]
On reconnaît $HR=2R$ : c'est la matrice d'une \cle{similitude} (composée d'une rotation d'angle $\dfrac{\pi}{2}$ et d'une homothétie de rapport $2$).

\textbf{3.} Image de $A(3\,;\,-1)$ :
\[ HR\begin{pmatrix} 3 \\ -1 \end{pmatrix}=\begin{pmatrix} 0 & -2 \\ 2 & 0 \end{pmatrix}\begin{pmatrix} 3 \\ -1 \end{pmatrix}=\begin{pmatrix} 0\times 3+(-2)\times(-1) \\ 2\times 3+0\times(-1) \end{pmatrix}=\begin{pmatrix} 2 \\ 6 \end{pmatrix}. \]
\textbf{Conclusion} : $t$ est la similitude de centre $O$, de rapport $2$ et d'angle $\dfrac{\pi}{2}$, et l'image de $A$ est le point $A'(2\,;\,6)$.
\end{exoresolu}

\begin{popfigure}
\begin{tikzpicture}[scale=1.2, >=Stealth]
\draw[popmuted, thin] (-3.5,-3.5) grid (3.5,3.5);
\draw[->, thick, popdark] (-4,0) -- (4,0) node[right] {$x$};
\draw[->, thick, popdark] (0,-4) -- (0,4) node[above] {$y$};
\foreach \x in {-3,-2,-1,1,2,3} \draw (\x,2pt) -- (\x,-2pt) node[below, font=\tiny] {\x};
\foreach \y in {-3,-2,-1,1,2,3} \draw (2pt,\y) -- (-2pt,\y) node[left, font=\tiny] {\y};

% Point A
\coordinate (A) at (3,-1);
\fill[poporange] (A) circle (2.5pt);
\node[below right, poporange] at (A) {$A(3;-1)$};

% Point A' image
\coordinate (Ap) at (2,6);
\fill[popblue] (Ap) circle (2.5pt);
\node[above right, popblue] at (Ap) {$A'(2;6)$};

% Arc de rotation
\draw[poporange, thick, ->] (1.5,0) arc (0:90:1.5);
\node[poporange] at (1.2,1.2) {$\frac{\pi}{2}$};

% Vecteurs
\draw[->, thick, poporange] (0,0) -- (A) node[midway, below right, font=\small] {$\vec{u}$};
\draw[->, thick, popblue] (0,0) -- (Ap) node[midway, above left, font=\small] {$\vec{u}'$};

% Cercle homothétie
\draw[popgreen, dashed] (0,0) circle (3.16);
\draw[popgreen, dashed] (0,0) circle (6.32);
\node[popgreen, font=\small] at (3.5, -3) {$\times 2$};
\end{tikzpicture}
\legende{Illustration de la similitude $t = h \circ r$ : rotation de $\pi/2$ puis homothétie de rapport 2.}
\end{popfigure}

\begin{aretenir}[Synthèse des formules clés — M20]
\begin{itemize}
\item \textbf{Opérations} : $A+B$ (coeff. par coeff.), $kA$, $AB$ (ligne $\times$ colonne). \textbf{Non-commutativité} : $AB\neq BA$.
\item \textbf{Déterminant} : $\det\begin{pmatrix} a & b \\ c & d \end{pmatrix}=ad-bc$.
\item \textbf{Inverse} : si $\det\neq 0$, $A^{-1}=\frac{1}{\det}\begin{pmatrix} d & -b \\ -c & a \end{pmatrix}$. Vérifier $AA^{-1}=I_2$.
\item \textbf{Système $AX=B$} : si $\det(A)\neq 0$, solution unique $X=A^{-1}B$.
\item \textbf{Matrice de $f$ lin.} : colonnes = coordonnées de $f(\vec{\imath}), f(\vec{\jmath})$. Image : $M\binom{x}{y}$.
\item \textbf{Composée} : $\mathrm{Mat}(g\circ f) = \mathrm{Mat}(g) \times \mathrm{Mat}(f)$.
\item \textbf{Automorphisme} : $f$ bijectif $\iff \det\neq 0$ ; $\mathrm{Mat}(f^{-1}) = (\mathrm{Mat}(f))^{-1}$.
\item \textbf{Géométrie (base orthonormée)} : Rotation $\begin{pmatrix} \cos\theta & -\sin\theta \\ \sin\theta & \cos\theta \end{pmatrix}$, Homothétie $kI_2$, Symétries $\pm I_2$ ou diag$(\pm 1, \pm 1)$.
\end{itemize}
\end{aretenir}

\begin{attention}[Les 5 erreurs qui coûtent des points au Bac]
\begin{enumerate}
\item \textbf{Croire que $AB=BA$.} Le produit matriciel n'est pas commutatif : calculer $AB$ quand on demande $BA$ est une erreur éliminatoire. De même, la matrice de $g\circ f$ est $B\times A$, pas $A\times B$ : l'application de droite agit en premier.
\item \textbf{Calculer $A^{-1}$ sans vérifier que $\det(A)\neq 0$.} La condition d'inversibilité doit être \emph{écrite explicitement} : « $\det(A)=\ldots\neq 0$, donc $A$ est inversible ». Sans cette phrase, la justification est incomplète et coûte des points.
\item \textbf{Se tromper dans la formule de l'inverse.} On échange $a$ et $d$, on change les signes de $b$ et $c$ \emph{uniquement}, et on n'oublie pas le facteur $\dfrac{1}{\det(A)}$ devant toute la matrice. Une vérification $AA^{-1}=I_2$ permet de détecter l'erreur.
\item \textbf{Ranger les images en lignes au lieu de colonnes.} Dans la matrice d'une application linéaire, les coordonnées de $f(\vec{\imath})$ forment la \emph{première colonne}, celles de $f(\vec{\jmath})$ la \emph{deuxième colonne}. Transposer la matrice fausse tous les calculs d'images.
\item \textbf{Mal écrire la forme matricielle d'un système.} Dans $AX=B$, la colonne $X=\begin{pmatrix} x \\ y \end{pmatrix}$ impose l'ordre : la première colonne de $A$ contient les coefficients de $x$, la deuxième ceux de $y$, en respectant les signes. Conclure sans vérifier la solution dans les équations de départ est aussi une rédaction incomplète.
\end{enumerate}
\end{attention}

\courssub{Exercices d'entraînement}

\begin{exercice}[Exercice 1 : Algèbre matricielle]
Soient les matrices $A=\begin{pmatrix} 1 & 2 \\ -1 & 3 \end{pmatrix}$, $B=\begin{pmatrix} 0 & 1 \\ 2 & -1 \end{pmatrix}$ et $C=\begin{pmatrix} 2 & 0 \\ 1 & 1 \end{pmatrix}$.
Calculer, dans l'ordre, les matrices suivantes (si elles ont un sens) :
$A+B$, $2A-B$, $AB$, $BA$, $A^2$, $A(BC)$, $(AB)C$.
Vérifier l'associativité du produit sur cet exemple.
\end{exercice}

\begin{exercice}[Exercice 2 : Inversibilité et système]
On considère la matrice $M=\begin{pmatrix} k & 2 \\ 3 & k \end{pmatrix}$ où $k$ est un réel.
\begin{enumerate}
\item Calculer $\det(M)$ en fonction de $k$.
\item Déterminer les valeurs de $k$ pour lesquelles $M$ est inversible.
\item Pour $k=4$, calculer $M^{-1}$ et résoudre le système $MX=\begin{pmatrix} 10 \\ 11 \end{pmatrix}$.
\end{enumerate}
\end{exercice}

\begin{exercice}[Exercice 3 : Application linéaire et géométrie]
Dans un repère orthonormé $(O;\vec{\imath},\vec{\jmath})$, on considère l'application linéaire $f$ de matrice $A=\begin{pmatrix} 0 & -2 \\ 2 & 0 \end{pmatrix}$.
\begin{enumerate}
\item Montrer que $f$ est une similitude. Préciser son rapport et son angle.
\item Déterminer l'image par $f$ du point $M(1;2)$.
\item Calculer $A^2$ et interpréter géométriquement la transformation $f\circ f$.
\end{enumerate}
\end{exercice}

\begin{exercice}[Exercice 4 : Problème concret — Téléphonie mobile]
Deux forfaits MTN sont proposés.
Forfait \textbf{Go} : un abonnement mensuel fixe + un prix par Go de data consommé.
Forfait \textbf{Illimité} : un abonnement mensuel plus élevé, data illimité.
Un client consomme $2$ Go le premier mois et paie $\num{3500}$ FCFA. Le mois suivant, il consomme $5$ Go et paie $\num{5000}$ FCFA.
\begin{enumerate}
\item Modéliser par un système linéaire pour trouver l'abonnement fixe et le prix du Go.
\item Résoudre par la méthode de la matrice inverse.
\item À partir de quelle consommation le forfait Illimité (abonnement $\num{8000}$ FCFA) devient-il plus avantageux ?
\end{enumerate}
\end{exercice}

\courssub{Corrigés des exercices d'entraînement}

\begin{corrige}[Exercice 1]
\begin{align*}
A+B &= \begin{pmatrix} 1 & 3 \\ 1 & 2 \end{pmatrix}, \quad
2A-B = \begin{pmatrix} 2 & 3 \\ -4 & 7 \end{pmatrix}, \\
AB &= \begin{pmatrix} 1\times0+2\times2 & 1\times1+2\times(-1) \\ -1\times0+3\times2 & -1\times1+3\times(-1) \end{pmatrix} = \begin{pmatrix} 4 & -1 \\ 6 & -4 \end{pmatrix}, \\
BA &= \begin{pmatrix} 0\times1+1\times(-1) & 0\times2+1\times3 \\ 2\times1+(-1)\times(-1) & 2\times2+(-1)\times3 \end{pmatrix} = \begin{pmatrix} -1 & 3 \\ 3 & 1 \end{pmatrix}, \\
A^2 &= \begin{pmatrix} -1 & 8 \\ -4 & 7 \end{pmatrix}, \\
BC &= \begin{pmatrix} 0\times2+1\times1 & 0\times0+1\times1 \\ 2\times2+(-1)\times1 & 2\times0+(-1)\times1 \end{pmatrix} = \begin{pmatrix} 1 & 1 \\ 3 & -1 \end{pmatrix}, \\
A(BC) &= \begin{pmatrix} 1 & 2 \\ -1 & 3 \end{pmatrix}\begin{pmatrix} 1 & 1 \\ 3 & -1 \end{pmatrix} = \begin{pmatrix} 7 & -1 \\ 8 & -4 \end{pmatrix}, \\
(AB)C &= \begin{pmatrix} 4 & -1 \\ 6 & -4 \end{pmatrix}\begin{pmatrix} 2 & 0 \\ 1 & 1 \end{pmatrix} = \begin{pmatrix} 7 & -1 \\ 8 & -4 \end{pmatrix}.
\end{align*}
On a bien $A(BC)=(AB)C$ : associativité vérifiée. Noter que $AB\neq BA$.
\end{corrige}

\begin{corrige}[Exercice 2]
\begin{enumerate}
\item $\det(M)=k\times k - 2\times 3 = k^2-6$.
\item $M$ inversible $\iff \det(M)\neq 0 \iff k^2-6\neq 0 \iff k\neq \sqrt{6}$ et $k\neq -\sqrt{6}$.
\item Pour $k=4$ : $M=\begin{pmatrix} 4 & 2 \\ 3 & 4 \end{pmatrix}$, $\det(M)=16-6=10\neq 0$.
$M^{-1}=\frac{1}{10}\begin{pmatrix} 4 & -2 \\ -3 & 4 \end{pmatrix}=\begin{pmatrix} 0,4 & -0,2 \\ -0,3 & 0,4 \end{pmatrix}$.
$X=M^{-1}\begin{pmatrix} 10 \\ 11 \end{pmatrix}=\frac{1}{10}\begin{pmatrix} 40-22 \\ -30+44 \end{pmatrix}=\frac{1}{10}\begin{pmatrix} 18 \\ 14 \end{pmatrix}=\begin{pmatrix} 1,8 \\ 1,4 \end{pmatrix}$.
Solution : $x=1,8$, $y=1,4$.
\end{enumerate}
\end{corrige}

\begin{corrige}[Exercice 3]
\begin{enumerate}
\item $A=\begin{pmatrix} 0 & -2 \\ 2 & 0 \end{pmatrix}=2\begin{pmatrix} 0 & -1 \\ 1 & 0 \end{pmatrix}$.
La matrice $\begin{pmatrix} 0 & -1 \\ 1 & 0 \end{pmatrix}$ est la rotation d'angle $\frac{\pi}{2}$.
Donc $f$ est la similitude de centre $O$, de rapport $2$ et d'angle $\frac{\pi}{2}$.
\item Image de $M(1;2)$ : $A\begin{pmatrix} 1 \\ 2 \end{pmatrix}=\begin{pmatrix} -4 \\ 2 \end{pmatrix}$. Donc $M'(-4;2)$.
\item $A^2 = \begin{pmatrix} 0 & -2 \\ 2 & 0 \end{pmatrix}^2 = \begin{pmatrix} -4 & 0 \\ 0 & -4 \end{pmatrix} = -4I_2$.
$f\circ f$ est l'homothétie de centre $O$ et de rapport $-4$ (ou rotation de $\pi$ composée avec homothétie de rapport $4$).
\end{enumerate}
\end{corrige}

\begin{corrige}[Exercice 4]
\begin{enumerate}
\item Soit $a$ l'abonnement fixe (en FCFA) et $p$ le prix du Go.
$\begin{cases} a + 2p = 3500 \\ a + 5p = 5000 \end{cases}$.
\item Forme matricielle : $\begin{pmatrix} 1 & 2 \\ 1 & 5 \end{pmatrix}\begin{pmatrix} a \\ p \end{pmatrix}=\begin{pmatrix} 3500 \\ 5000 \end{pmatrix}$.
$\det = 5-2=3\neq 0$.
$A^{-1}=\frac{1}{3}\begin{pmatrix} 5 & -2 \\ -1 & 1 \end{pmatrix}$.
$X=\frac{1}{3}\begin{pmatrix} 5 & -2 \\ -1 & 1 \end{pmatrix}\begin{pmatrix} 3500 \\ 5000 \end{pmatrix}=\frac{1}{3}\begin{pmatrix} 17500-10000 \\ -3500+5000 \end{pmatrix}=\frac{1}{3}\begin{pmatrix} 7500 \\ 1500 \end{pmatrix}=\begin{pmatrix} 2500 \\ 500 \end{pmatrix}$.
Abonnement = $\num{2500}$ FCFA, Prix du Go = $\num{500}$ FCFA.
Vérification : $2500+2\times500=3500$, $2500+5\times500=5000$.
\item Forfait Go : $C_1(x)=2500+500x$. Forfait Illimité : $C_2=8000$.
$C_1(x) \ge C_2 \iff 2500+500x \ge 8000 \iff 500x \ge 5500 \iff x \ge 11$.
À partir de $11$ Go consommés, le forfait Illimité est plus avantageux.
\end{enumerate}
\end{corrige}

\newpage

\coursec{Exercices}

\courssub{Je vérifie mes connaissances}

\begin{exercice}[QCM : opérations sur les matrices]
Pour chaque question, une seule réponse est exacte. On considère les matrices
$A=\begin{pmatrix} 1 & 2 \\ 3 & 4 \end{pmatrix}$ et $B=\begin{pmatrix} 0 & 1 \\ 1 & 0 \end{pmatrix}$.
\begin{enumerate}
\item La somme $A+B$ est égale à :
\begin{enumerate}
\item $\begin{pmatrix} 1 & 3 \\ 4 & 4 \end{pmatrix}$
\item $\begin{pmatrix} 1 & 2 \\ 3 & 4 \end{pmatrix}$
\item $\begin{pmatrix} 1 & 1 \\ 1 & 4 \end{pmatrix}$
\end{enumerate}
\item Le produit $AB$ est égal à :
\begin{enumerate}
\item $\begin{pmatrix} 2 & 1 \\ 4 & 3 \end{pmatrix}$
\item $\begin{pmatrix} 3 & 4 \\ 1 & 2 \end{pmatrix}$
\item $\begin{pmatrix} 0 & 2 \\ 3 & 0 \end{pmatrix}$
\end{enumerate}
\item Le déterminant de $A$ vaut :
\begin{enumerate}
\item $10$
\item $-2$
\item $2$
\end{enumerate}
\item L'inverse de $B$ est :
\begin{enumerate}
\item $B$ elle-même
\item $-B$
\item $B$ n'est pas inversible
\end{enumerate}
\end{enumerate}
\end{exercice}

\begin{exercice}[Vrai ou faux, avec justification]
Pour chacune des affirmations suivantes, dire si elle est vraie ou fausse en justifiant soigneusement (une démonstration courte ou un contre-exemple est attendu).
\begin{enumerate}
\item « Si $A$ et $B$ sont deux matrices carrées d'ordre 2 telles que $AB=I_2$, alors $BA=I_2$. »
\item « Si $\det(A)=0$, alors $A$ est la matrice nulle. »
\item « Pour toutes matrices $A$ et $B$ d'ordre 2, on a $(A+B)^2=A^2+2AB+B^2$. »
\item « Toute application linéaire du plan dont la matrice dans la base canonique a un déterminant non nul est un automorphisme. »
\end{enumerate}
\end{exercice}

\begin{exercice}[Définitions et réflexes de cours]
\begin{enumerate}
\item Donner la définition du déterminant d'une matrice $M=\begin{pmatrix} a & b \\ c & d \end{pmatrix}$.
\item À quelle condition une matrice carrée d'ordre 2 est-elle inversible ? Donner alors la formule de son inverse.
\item Rappeler comment on construit la matrice d'une application linéaire $f$ du plan dans une base $(\vec{\imath},\vec{\jmath})$ à partir des images $f(\vec{\imath})$ et $f(\vec{\jmath})$.
\item Soient $f$ et $g$ deux applications linéaires du plan de matrices respectives $M$ et $N$ dans une même base. Quelle est la matrice de $g\circ f$ dans cette base ?
\end{enumerate}
\end{exercice}

\begin{exercice}[Calculs directs]
On donne $A=\begin{pmatrix} 3 & -1 \\ 2 & 1 \end{pmatrix}$ et $B=\begin{pmatrix} -2 & 4 \\ 1 & 3 \end{pmatrix}$.
\begin{enumerate}
\item Calculer $A+B$, $2A-3B$, $AB$ et $BA$. Que constate-t-on pour les deux derniers résultats ?
\item Calculer $\det(A)$ et $\det(B)$. Ces matrices sont-elles inversibles ?
\item Calculer $\det(AB)$ et vérifier que $\det(AB)=\det(A)\times\det(B)$.
\end{enumerate}
\end{exercice}

\courssub{Je m'entraîne}

\begin{exercice}[Inverse et système linéaire]
On considère la matrice $A=\begin{pmatrix} 2 & 1 \\ 1 & 2 \end{pmatrix}$.
\begin{enumerate}
\item Calculer $A^2$, puis $\det(A)$.
\item Justifier que $A$ est inversible et déterminer $A^{-1}$.
\item En déduire la résolution du système $\begin{cases} 2x+y=7 \\ x+2y=8 \end{cases}$ à l'aide de la matrice inverse.
\end{enumerate}
\end{exercice}

\begin{exercice}[Matrice d'une application linéaire]
Le plan est muni de la base canonique $(\vec{\imath},\vec{\jmath})$. Soit $f$ l'application linéaire du plan définie par
$f(\vec{\imath})=2\vec{\imath}-\vec{\jmath}$ et $f(\vec{\jmath})=\vec{\imath}+3\vec{\jmath}$.
\begin{enumerate}
\item Déterminer la matrice $M$ de $f$ dans la base canonique.
\item Calculer les coordonnées de l'image par $f$ du vecteur $\vec{u}(4\,;-2)$.
\item $f$ est-elle un automorphisme du plan ? Justifier.
\item Si oui, déterminer la matrice de l'application réciproque $f^{-1}$.
\end{enumerate}
\end{exercice}

\begin{exercice}[Somme et composée d'applications linéaires]
Soient $f$ et $g$ deux applications linéaires du plan dont les matrices dans la base canonique sont
$M=\begin{pmatrix} 1 & 2 \\ 0 & 1 \end{pmatrix}$ et $N=\begin{pmatrix} 2 & 0 \\ 1 & -1 \end{pmatrix}$.
\begin{enumerate}
\item Déterminer la matrice de $f+g$.
\item Déterminer la matrice de $g\circ f$, puis celle de $f\circ g$. Comparer.
\item Calculer les coordonnées de l'image du vecteur $\vec{u}(3\,;-1)$ par $g\circ f$.
\item Montrer que $f$ est un automorphisme et déterminer la matrice de $f^{-1}$.
\end{enumerate}
\end{exercice}

\begin{exercice}[Une transformation géométrique]
Le plan est muni d'un repère orthonormé $(O\,;\vec{\imath},\vec{\jmath})$. On considère la transformation $s$ du plan qui, à tout point $M(x\,;y)$, associe le point $M'(x'\,;y')$ tel que
$\begin{pmatrix} x' \\ y' \end{pmatrix}=\begin{pmatrix} 1 & 0 \\ 0 & -1 \end{pmatrix}\begin{pmatrix} x \\ y \end{pmatrix}$.
\begin{enumerate}
\item Déterminer l'ensemble des points fixes de $s$ (points tels que $M'=M$).
\item Déterminer les images des points $O(0\,;0)$, $A(1\,;0)$ et $B(0\,;1)$.
\item Calculer l'aire du triangle $OAB$, puis celle de son image. Que constate-t-on ?
\item Quelle est la nature géométrique de la transformation $s$ ?
\end{enumerate}
\end{exercice}

\begin{popfigure}
\begin{tikzpicture}[scale=1.5, >=Stealth]
\draw[popline, thin] (-1.5, -1.5) grid (1.5, 1.5);
\draw[->, thick] (-1.5,0) -- (1.5,0) node[right] {$\vec{\imath}$};
\draw[->, thick] (0,-1.5) -- (0,1.5) node[above] {$\vec{\jmath}$};
% Original points
\fill[popblue] (0,0) circle (2pt) node[below left] {$O$};
\fill[popblue] (1,0) circle (2pt) node[below] {$A(1;0)$};
\fill[popblue] (0,1) circle (2pt) node[left] {$B(0;1)$};
\draw[popblue, thick] (0,0) -- (1,0) -- (0,1) -- cycle;
% Image points
\fill[poporange] (0,0) circle (2pt);
\fill[poporange] (1,0) circle (2pt);
\fill[poporange] (0,-1) circle (2pt) node[right] {$B'(0;-1)$};
\draw[poporange, thick, dashed] (0,0) -- (1,0) -- (0,-1) -- cycle;
% Axis label
\node[popdark, align=center] at (1.2, -1.2) {Axe des abscisses\\(points fixes)};
\end{tikzpicture}
\legende{Symétrie axiale par rapport à l'axe des abscisses : $O$ et $A$ sont fixes, $B$ va en $B'$.}
\end{popfigure}

\courssub{Je me prépare au Bac}

\begin{exercice}[Type Bac : matrice, inverse et système]
\textbf{Partie A.} On considère la matrice $A=\begin{pmatrix} 3 & 2 \\ 1 & 1 \end{pmatrix}$.
\begin{enumerate}
\item Calculer $\det(A)$ et en déduire que $A$ est inversible.
\item Déterminer $A^{-1}$.
\item Calculer $A^2$ et vérifier que $A^2=4A-I_2$, où $I_2$ désigne la matrice unité d'ordre 2.
\item En déduire une expression de $A^{-1}$ en fonction de $A$ et $I_2$, et retrouver le résultat de la question 2.
\end{enumerate}
\textbf{Partie B.} Un commerçant du marché Mokolo à Yaoundé vend deux lots de produits. Le premier lot contient 3 sacs de riz et 2 bidons d'huile pour un total de \SI{74000}{F} ; le deuxième lot contient 1 sac de riz et 1 bidon d'huile pour un total de \SI{27000}{F}. On note $x$ le prix d'un sac de riz et $y$ celui d'un bidon d'huile, en milliers de francs.
\begin{enumerate}
\item Traduire la situation par un système de deux équations à deux inconnues, puis par une écriture matricielle $A\begin{pmatrix} x \\ y \end{pmatrix}=\begin{pmatrix} 74 \\ 27 \end{pmatrix}$.
\item À l'aide de $A^{-1}$, déterminer le prix d'un sac de riz et celui d'un bidon d'huile.
\end{enumerate}
\end{exercice}

\begin{exercice}[Type Bac : applications linéaires et composition]
Le plan est muni de la base canonique $(\vec{\imath},\vec{\jmath})$. On considère les applications linéaires $f$ et $g$ du plan définies par :
\[ f(x\vec{\imath}+y\vec{\jmath})=(x+y)\vec{\imath}+(x-y)\vec{\jmath} \qquad\text{et}\qquad g(x\vec{\imath}+y\vec{\jmath})=(2x)\vec{\imath}+(x+y)\vec{\jmath}. \]
\begin{enumerate}
\item Déterminer les matrices $M$ et $N$ de $f$ et $g$ dans la base canonique.
\item Montrer que $f$ est un automorphisme du plan et déterminer la matrice de $f^{-1}$.
\item Déterminer la matrice $P$ de $g\circ f$ dans la base canonique.
\item Soit $\vec{u}$ le vecteur de coordonnées $(1\,;3)$. Calculer les coordonnées de $g\circ f(\vec{u})$ de deux manières différentes.
\item Déterminer les coordonnées du vecteur $\vec{v}$ tel que $f(\vec{v})=5\vec{\imath}+\vec{\jmath}$.
\item L'application $g$ est-elle un automorphisme ? Justifier.
\end{enumerate}
\end{exercice}

\begin{exercice}[Problème type Bac : la coopérative agricole]
Une coopérative de la région de l'Ouest transforme chaque mois sa production. On note $m_n$ et $c_n$ les quantités (en tonnes) de maïs et de manioc disponibles au début du mois numéro $n$. Chaque mois, la coopérative vend une partie de son stock et en reçoit de nouvelles quantités, de sorte que les quantités au début du mois suivant sont données par :
\[ \begin{pmatrix} m_{n+1} \\ c_{n+1} \end{pmatrix}=T\begin{pmatrix} m_n \\ c_n \end{pmatrix} \qquad\text{avec}\qquad T=\begin{pmatrix} 2 & 1 \\ 1 & 2 \end{pmatrix}. \]
Au début du premier mois, la coopérative dispose de $m_1=10$ tonnes de maïs et $c_1=4$ tonnes de manioc.
\begin{enumerate}
\item Calculer les quantités de maïs et de manioc au début du deuxième mois.
\item Calculer $T^2$, puis en déduire les quantités prévisibles au début du troisième mois.
\item Calculer $\det(T)$ et justifier que $T$ est inversible. Déterminer $T^{-1}$.
\item À la fin d'un certain mois, le bilan indique $46$ tonnes de maïs et $38$ tonnes de manioc. À l'aide de $T^{-1}$, retrouver les quantités dont disposait la coopérative au début de ce mois.
\item Un agronome affirme : « Si l'on applique deux fois de suite la transformation, la matrice globale est $2T$. » A-t-il raison ? Justifier à l'aide de la question 2.
\end{enumerate}
\end{exercice}

\newpage

\coursec{Corrigés des exercices}

\begin{corrige}[Exercice 1 — QCM : opérations sur les matrices]
\begin{enumerate}
\item On additionne coefficient par coefficient :
\[ A+B=\begin{pmatrix} 1+0 & 2+1 \\ 3+1 & 4+0 \end{pmatrix}=\begin{pmatrix} 1 & 3 \\ 4 & 4 \end{pmatrix}. \]
\textbf{Réponse a.}

\item Produit lignes par colonnes :
\[ AB=\begin{pmatrix} 1\times0+2\times1 & 1\times1+2\times0 \\ 3\times0+4\times1 & 3\times1+4\times0 \end{pmatrix}=\begin{pmatrix} 2 & 1 \\ 4 & 3 \end{pmatrix}. \]
\textbf{Réponse a.}

\item $\det(A)=1\times4-2\times3=4-6=-2$. \textbf{Réponse b.}

\item On calcule $B^2=\begin{pmatrix} 0 & 1 \\ 1 & 0 \end{pmatrix}\begin{pmatrix} 0 & 1 \\ 1 & 0 \end{pmatrix}=\begin{pmatrix} 0+1 & 0+0 \\ 0+0 & 1+0 \end{pmatrix}=\begin{pmatrix} 1 & 0 \\ 0 & 1 \end{pmatrix}=I_2$. Donc $B\times B=I_2$, ce qui prouve que $B$ est inversible et que $B^{-1}=B$. \textbf{Réponse a.}
\end{enumerate}
\textbf{Points de vigilance :} le produit matriciel se fait \emph{lignes} de la première matrice contre \emph{colonnes} de la seconde, et non coefficient par coefficient. Pour le déterminant, attention au signe : $\det(A)=ad-bc$, et non $ad+bc$.
\end{corrige}

\begin{corrige}[Exercice 2 — Vrai ou faux, avec justification]
\begin{enumerate}
\item \textbf{Vrai.} Pour les matrices carrées d'ordre 2, si $AB=I_2$, alors $A$ est inversible et $B=A^{-1}$ ; on a alors aussi $BA=I_2$. Justification par les déterminants : $\det(AB)=\det(A)\det(B)=\det(I_2)=1$, donc $\det(A)\neq0$ et $A$ est inversible ; en multipliant l'égalité $AB=I_2$ à gauche par $A^{-1}$ on obtient $B=A^{-1}$, puis en multipliant à gauche par $A$ : $BA=AA^{-1}=I_2$.

\item \textbf{Faux.} Contre-exemple : $A=\begin{pmatrix} 1 & 1 \\ 1 & 1 \end{pmatrix}$ n'est pas la matrice nulle, pourtant $\det(A)=1\times1-1\times1=0$. La nullité du déterminant signifie seulement que $A$ n'est \emph{pas inversible}.

\item \textbf{Faux.} En développant : $(A+B)^2=(A+B)(A+B)=A^2+AB+BA+B^2$. Or le produit matriciel n'est pas commutatif : en général $BA\neq AB$, donc $AB+BA\neq 2AB$. Contre-exemple avec $A=\begin{pmatrix} 1 & 1 \\ 0 & 1 \end{pmatrix}$ et $B=\begin{pmatrix} 0 & 1 \\ 1 & 0 \end{pmatrix}$ : $AB=\begin{pmatrix} 1 & 1 \\ 1 & 0 \end{pmatrix}$ et $BA=\begin{pmatrix} 0 & 1 \\ 1 & 1 \end{pmatrix}$, donc $AB\neq BA$ et l'identité remarquable usuelle ne s'applique pas.

\item \textbf{Vrai.} Une application linéaire du plan est un automorphisme si et seulement si sa matrice dans une base est inversible, c'est-à-dire si et seulement si son déterminant est non nul. Comme le déterminant est non nul par hypothèse, la matrice est inversible et $f$ est bijective : c'est un automorphisme du plan.
\end{enumerate}
\textbf{Points de vigilance :} pour infirmer une affirmation, un seul contre-exemple suffit, mais il faut le vérifier explicitement. Ne jamais utiliser les identités remarquables sans vérifier la commutativité.
\end{corrige}

\begin{corrige}[Exercice 3 — Définitions et réflexes de cours]
\begin{enumerate}
\item Le \cle{déterminant} de $M=\begin{pmatrix} a & b \\ c & d \end{pmatrix}$ est le réel
\[ \det(M)=ad-bc. \]

\item Une matrice carrée $M$ d'ordre 2 est \cle{inversible} si et seulement si $\det(M)\neq0$. Dans ce cas :
\[ M^{-1}=\frac{1}{ad-bc}\begin{pmatrix} d & -b \\ -c & a \end{pmatrix}. \]

\item La matrice de $f$ dans la base $(\vec{\imath},\vec{\jmath})$ est la matrice dont la \textbf{première colonne} est formée des coordonnées de $f(\vec{\imath})$ et la \textbf{deuxième colonne} des coordonnées de $f(\vec{\jmath})$ dans cette base : si $f(\vec{\imath})=a\vec{\imath}+c\vec{\jmath}$ et $f(\vec{\jmath})=b\vec{\imath}+d\vec{\jmath}$, alors $\mathrm{Mat}(f)=\begin{pmatrix} a & b \\ c & d \end{pmatrix}$.

\item La matrice d'une composée est le produit des matrices dans le même ordre :
\[ \mathrm{Mat}(g\circ f)=N\times M. \]
\end{enumerate}
\textbf{Points de vigilance :} les images des vecteurs de base se rangent \emph{en colonnes}, pas en lignes. Dans $g\circ f$, c'est $f$ qui agit en premier, mais la matrice est $NM$ (attention à l'ordre du produit).
\end{corrige}

\begin{corrige}[Exercice 4 — Calculs directs]
\begin{enumerate}
\item Somme et combinaison linéaire, coefficient par coefficient :
\[ A+B=\begin{pmatrix} 3-2 & -1+4 \\ 2+1 & 1+3 \end{pmatrix}=\begin{pmatrix} 1 & 3 \\ 3 & 4 \end{pmatrix}. \]
\[ 2A-3B=\begin{pmatrix} 6 & -2 \\ 4 & 2 \end{pmatrix}-\begin{pmatrix} -6 & 12 \\ 3 & 9 \end{pmatrix}=\begin{pmatrix} 12 & -14 \\ 1 & -7 \end{pmatrix}. \]
Produits lignes par colonnes :
\[ AB=\begin{pmatrix} 3\times(-2)+(-1)\times1 & 3\times4+(-1)\times3 \\ 2\times(-2)+1\times1 & 2\times4+1\times3 \end{pmatrix}=\begin{pmatrix} -7 & 9 \\ -3 & 11 \end{pmatrix}. \]
\[ BA=\begin{pmatrix} (-2)\times3+4\times2 & (-2)\times(-1)+4\times1 \\ 1\times3+3\times2 & 1\times(-1)+3\times1 \end{pmatrix}=\begin{pmatrix} 2 & 6 \\ 9 & 2 \end{pmatrix}. \]
\textbf{Constat :} $AB\neq BA$ : le produit matriciel n'est \emph{pas commutatif} en général.

\item $\det(A)=3\times1-(-1)\times2=3+2=5\neq0$ : $A$ est inversible. $\det(B)=(-2)\times3-4\times1=-6-4=-10\neq0$ : $B$ est inversible.

\item $\det(AB)=(-7)\times11-9\times(-3)=-77+27=-50$. Vérification :
\[ \det(A)\times\det(B)=5\times(-10)=-50=\det(AB). \]
La propriété $\det(AB)=\det(A)\det(B)$ est bien vérifiée.
\end{enumerate}
\end{corrige}

\begin{corrige}[Exercice 5 — Inverse et système linéaire]
\begin{enumerate}
\item Produit lignes par colonnes :
\[ A^2=\begin{pmatrix} 2 & 1 \\ 1 & 2 \end{pmatrix}\begin{pmatrix} 2 & 1 \\ 1 & 2 \end{pmatrix}=\begin{pmatrix} 4+1 & 2+2 \\ 2+2 & 1+4 \end{pmatrix}=\begin{pmatrix} 5 & 4 \\ 4 & 5 \end{pmatrix}. \]
$\det(A)=2\times2-1\times1=4-1=3$.

\item Comme $\det(A)=3\neq0$, $A$ est inversible et
\[ A^{-1}=\frac{1}{3}\begin{pmatrix} 2 & -1 \\ -1 & 2 \end{pmatrix}=\begin{pmatrix} \tfrac{2}{3} & -\tfrac{1}{3} \\[2pt] -\tfrac{1}{3} & \tfrac{2}{3} \end{pmatrix}. \]
Vérification : $AA^{-1}=\dfrac{1}{3}\begin{pmatrix} 4-1 & -2+2 \\ 2-2 & -1+4 \end{pmatrix}=\dfrac{1}{3}\begin{pmatrix} 3 & 0 \\ 0 & 3 \end{pmatrix}=I_2$.

\item Le système s'écrit matriciellement $A\begin{pmatrix} x \\ y \end{pmatrix}=\begin{pmatrix} 7 \\ 8 \end{pmatrix}$. Comme $A$ est inversible :
\[ \begin{pmatrix} x \\ y \end{pmatrix}=A^{-1}\begin{pmatrix} 7 \\ 8 \end{pmatrix}=\frac{1}{3}\begin{pmatrix} 2 & -1 \\ -1 & 2 \end{pmatrix}\begin{pmatrix} 7 \\ 8 \end{pmatrix}=\frac{1}{3}\begin{pmatrix} 14-8 \\ -7+16 \end{pmatrix}=\frac{1}{3}\begin{pmatrix} 6 \\ 9 \end{pmatrix}=\begin{pmatrix} 2 \\ 3 \end{pmatrix}. \]
Vérification : $2\times2+3=7$ et $2+2\times3=8$. L'unique solution est $(x\,;y)=(2\,;3)$.
\end{enumerate}
\textbf{Points de vigilance :} citer la condition $\det(A)\neq0$ \emph{avant} d'écrire $A^{-1}$ ; la solution du système $AX=B$ est $X=A^{-1}B$ (multiplication à gauche, dans cet ordre).
\end{corrige}

\begin{corrige}[Exercice 6 — Matrice d'une application linéaire]
\begin{enumerate}
\item Les images des vecteurs de base se rangent en colonnes : $f(\vec{\imath})=2\vec{\imath}-\vec{\jmath}$ donne la première colonne $\begin{pmatrix} 2 \\ -1 \end{pmatrix}$ et $f(\vec{\jmath})=\vec{\imath}+3\vec{\jmath}$ la deuxième $\begin{pmatrix} 1 \\ 3 \end{pmatrix}$. D'où
\[ M=\begin{pmatrix} 2 & 1 \\ -1 & 3 \end{pmatrix}. \]

\item Les coordonnées de $f(\vec{u})$ s'obtiennent par produit matriciel :
\[ \begin{pmatrix} 2 & 1 \\ -1 & 3 \end{pmatrix}\begin{pmatrix} 4 \\ -2 \end{pmatrix}=\begin{pmatrix} 8-2 \\ -4-6 \end{pmatrix}=\begin{pmatrix} 6 \\ -10 \end{pmatrix}. \]
Donc $f(\vec{u})$ a pour coordonnées $(6\,;-10)$.

\item $\det(M)=2\times3-1\times(-1)=6+1=7\neq0$ : la matrice $M$ est inversible, donc $f$ est un \cle{automorphisme} du plan.

\item La matrice de $f^{-1}$ est l'inverse de $M$ :
\[ M^{-1}=\frac{1}{7}\begin{pmatrix} 3 & -1 \\ 1 & 2 \end{pmatrix}. \]
\end{enumerate}
\textbf{Points de vigilance :} ne pas inverser lignes et colonnes lors de la construction de $M$ ; justifier le caractère bijectif par le déterminant non nul.
\end{corrige}

\begin{corrige}[Exercice 7 — Somme et composée d'applications linéaires]
\begin{enumerate}
\item La matrice d'une somme est la somme des matrices :
\[ \mathrm{Mat}(f+g)=M+N=\begin{pmatrix} 1+2 & 2+0 \\ 0+1 & 1-1 \end{pmatrix}=\begin{pmatrix} 3 & 2 \\ 1 & 0 \end{pmatrix}. \]

\item La matrice de $g\circ f$ est $NM$ :
\[ NM=\begin{pmatrix} 2 & 0 \\ 1 & -1 \end{pmatrix}\begin{pmatrix} 1 & 2 \\ 0 & 1 \end{pmatrix}=\begin{pmatrix} 2+0 & 4+0 \\ 1+0 & 2-1 \end{pmatrix}=\begin{pmatrix} 2 & 4 \\ 1 & 1 \end{pmatrix}. \]
Celle de $f\circ g$ est $MN$ :
\[ MN=\begin{pmatrix} 1 & 2 \\ 0 & 1 \end{pmatrix}\begin{pmatrix} 2 & 0 \\ 1 & -1 \end{pmatrix}=\begin{pmatrix} 2+2 & 0-2 \\ 0+1 & 0-1 \end{pmatrix}=\begin{pmatrix} 4 & -2 \\ 1 & -1 \end{pmatrix}. \]
\textbf{Comparaison :} $NM\neq MN$, donc $g\circ f\neq f\circ g$ : la composition des applications linéaires n'est pas commutative en général.

\item On applique le produit matriciel avec la matrice de $g\circ f$ :
\[ \begin{pmatrix} 2 & 4 \\ 1 & 1 \end{pmatrix}\begin{pmatrix} 3 \\ -1 \end{pmatrix}=\begin{pmatrix} 6-4 \\ 3-1 \end{pmatrix}=\begin{pmatrix} 2 \\ 2 \end{pmatrix}. \]
Donc $g\circ f(\vec{u})$ a pour coordonnées $(2\,;2)$. Vérification directe : $f(\vec{u})$ a pour coordonnées $(3-2\,;-1)=(1\,;-1)$, puis $g$ envoie $(1\,;-1)$ sur $(2\times1\,;1-(-1))=(2\,;2)$.

\item $\det(M)=1\times1-2\times0=1\neq0$, donc $M$ est inversible et $f$ est un automorphisme. La matrice de $f^{-1}$ est
\[ M^{-1}=\frac{1}{1}\begin{pmatrix} 1 & -2 \\ 0 & 1 \end{pmatrix}=\begin{pmatrix} 1 & -2 \\ 0 & 1 \end{pmatrix}. \]
\end{enumerate}
\end{corrige}

\begin{corrige}[Exercice 8 — Une transformation géométrique]
\begin{enumerate}
\item $M(x\,;y)$ est fixe si et seulement si $x'=x$ et $y'=y$, c'est-à-dire
\[ \begin{cases} x=x \\ -y=y \end{cases}\iff y=0. \]
L'ensemble des points fixes de $s$ est donc \textbf{l'axe des abscisses} $(O\,;\vec{\imath})$.

\item Par produit matriciel :
\[ s(O)=(0\,;0)=O,\qquad s(A)=(1\,;0)=A,\qquad s(B)=(0\,;-1)=B'. \]

\item Le triangle $OAB$ est rectangle en $O$ avec $OA=1$ et $OB=1$, donc
\[ \mathcal{A}(OAB)=\frac{OA\times OB}{2}=\frac{1}{2}. \]
Son image $OA'B'$ avec $A'=A(1\,;0)$ et $B'(0\,;-1)$ est aussi rectangle en $O$ avec $OA'=1$ et $OB'=1$, donc $\mathcal{A}(OA'B')=\dfrac{1}{2}$. \textbf{Constat :} la transformation $s$ conserve les aires (on remarque que la valeur absolue du déterminant de sa matrice vaut $|{-1}|=1$).

\item $s$ fixe tous les points de l'axe des abscisses et envoie tout point sur son symétrique par rapport à cet axe : $s$ est la \cle{symétrie axiale} (réflexion) d'axe $(O\,;\vec{\imath})$.
\end{enumerate}
\textbf{Points de vigilance :} un point fixe vérifie $M'=M$, ce qui donne un système à résoudre ; ne pas conclure « $s$ est l'identité » parce que $O$ et $A$ sont fixes : $B$ ne l'est pas.
\end{corrige}

\begin{corrige}[Exercice 9 — Type Bac : matrice, inverse et système]
\textbf{Partie A.}
\begin{enumerate}
\item $\det(A)=3\times1-2\times1=3-2=1\neq0$, donc $A$ est inversible.

\item $A^{-1}=\dfrac{1}{1}\begin{pmatrix} 1 & -2 \\ -1 & 3 \end{pmatrix}=\begin{pmatrix} 1 & -2 \\ -1 & 3 \end{pmatrix}$.

\item On calcule :
\[ A^2=\begin{pmatrix} 3 & 2 \\ 1 & 1 \end{pmatrix}\begin{pmatrix} 3 & 2 \\ 1 & 1 \end{pmatrix}=\begin{pmatrix} 9+2 & 6+2 \\ 3+1 & 2+1 \end{pmatrix}=\begin{pmatrix} 11 & 8 \\ 4 & 5 \end{pmatrix}. \]
D'autre part :
\[ 4A-I_2=\begin{pmatrix} 12 & 8 \\ 4 & 4 \end{pmatrix}-\begin{pmatrix} 1 & 0 \\ 0 & 1 \end{pmatrix}=\begin{pmatrix} 11 & 8 \\ 4 & 5 \end{pmatrix}=A^2. \]
La relation $A^2=4A-I_2$ est bien vérifiée.

\item De $A^2=4A-I_2$, on tire $I_2=4A-A^2=A(4I_2-A)$. Ainsi $A\times(4I_2-A)=I_2$, ce qui prouve que $A$ est inversible et que
\[ A^{-1}=4I_2-A=\begin{pmatrix} 4 & 0 \\ 0 & 4 \end{pmatrix}-\begin{pmatrix} 3 & 2 \\ 1 & 1 \end{pmatrix}=\begin{pmatrix} 1 & -2 \\ -1 & 3 \end{pmatrix}. \]
On retrouve bien le résultat de la question 2.
\end{enumerate}
\textbf{Partie B.}
\begin{enumerate}
\item Le premier lot donne $3x+2y=74$ et le deuxième $x+y=27$ (en milliers de francs). D'où le système
\[ \begin{cases} 3x+2y=74 \\ x+y=27 \end{cases}\iff \underbrace{\begin{pmatrix} 3 & 2 \\ 1 & 1 \end{pmatrix}}_{A}\begin{pmatrix} x \\ y \end{pmatrix}=\begin{pmatrix} 74 \\ 27 \end{pmatrix}. \]

\item Comme $A$ est inversible :
\[ \begin{pmatrix} x \\ y \end{pmatrix}=A^{-1}\begin{pmatrix} 74 \\ 27 \end{pmatrix}=\begin{pmatrix} 1 & -2 \\ -1 & 3 \end{pmatrix}\begin{pmatrix} 74 \\ 27 \end{pmatrix}=\begin{pmatrix} 74-54 \\ -74+81 \end{pmatrix}=\begin{pmatrix} 20 \\ 7 \end{pmatrix}. \]
Vérification : $3\times20+2\times7=60+14=74$ et $20+7=27$. Le sac de riz coûte \SI{20000}{F} et le bidon d'huile \SI{7000}{F}.
\end{enumerate}
\textbf{Barème indicatif (sur 5 pts) :} déterminant + conclusion (0,5) ; $A^{-1}$ (0,5) ; $A^2$ et vérification (1) ; expression $A^{-1}=4I_2-A$ (1) ; système et écriture matricielle (1) ; résolution et interprétation en francs (1).

\textbf{Points de vigilance :} dans la question A.4, il faut \emph{factoriser} $4A-A^2=A(4I_2-A)$ et non diviser par $A$ (une matrice ne se divise pas) ; penser à convertir les milliers de francs en francs dans la réponse finale ; vérifier la solution en la substituant dans le système.
\end{corrige}

\begin{corrige}[Exercice 10 — Type Bac : applications linéaires et composition]
\begin{enumerate}
\item On lit les images des vecteurs de base. Pour $f$ : $f(\vec{\imath})=\vec{\imath}+\vec{\jmath}$ et $f(\vec{\jmath})=\vec{\imath}-\vec{\jmath}$, rangés en colonnes :
\[ M=\begin{pmatrix} 1 & 1 \\ 1 & -1 \end{pmatrix}. \]
Pour $g$ : $g(\vec{\imath})=2\vec{\imath}+\vec{\jmath}$ et $g(\vec{\jmath})=\vec{\jmath}$, d'où
\[ N=\begin{pmatrix} 2 & 0 \\ 1 & 1 \end{pmatrix}. \]

\item $\det(M)=1\times(-1)-1\times1=-2\neq0$ : $M$ est inversible, donc $f$ est un automorphisme du plan. La matrice de $f^{-1}$ est
\[ M^{-1}=\frac{1}{-2}\begin{pmatrix} -1 & -1 \\ -1 & 1 \end{pmatrix}=\begin{pmatrix} \tfrac{1}{2} & \tfrac{1}{2} \\[2pt] \tfrac{1}{2} & -\tfrac{1}{2} \end{pmatrix}. \]

\item La matrice de $g\circ f$ est $P=NM$ :
\[ P=\begin{pmatrix} 2 & 0 \\ 1 & 1 \end{pmatrix}\begin{pmatrix} 1 & 1 \\ 1 & -1 \end{pmatrix}=\begin{pmatrix} 2+0 & 2+0 \\ 1+1 & 1-1 \end{pmatrix}=\begin{pmatrix} 2 & 2 \\ 2 & 0 \end{pmatrix}. \]

\item \textbf{Première méthode} (par le produit matriciel) :
\[ P\begin{pmatrix} 1 \\ 3 \end{pmatrix}=\begin{pmatrix} 2 & 2 \\ 2 & 0 \end{pmatrix}\begin{pmatrix} 1 \\ 3 \end{pmatrix}=\begin{pmatrix} 2+6 \\ 2+0 \end{pmatrix}=\begin{pmatrix} 8 \\ 2 \end{pmatrix}. \]
\textbf{Deuxième méthode} (en composant les expressions) : $f(\vec{u})=(1+3)\vec{\imath}+(1-3)\vec{\jmath}=4\vec{\imath}-2\vec{\jmath}$, puis
\[ g(4\vec{\imath}-2\vec{\jmath})=(2\times4)\vec{\imath}+(4-2)\vec{\jmath}=8\vec{\imath}+2\vec{\jmath}. \]
Les deux méthodes donnent les mêmes coordonnées : $g\circ f(\vec{u})$ a pour coordonnées $(8\,;2)$.

\item On cherche $\vec{v}(x\,;y)$ tel que $f(\vec{v})=5\vec{\imath}+\vec{\jmath}$. Comme $f$ est bijective :
\[ \begin{pmatrix} x \\ y \end{pmatrix}=M^{-1}\begin{pmatrix} 5 \\ 1 \end{pmatrix}=\begin{pmatrix} \tfrac{1}{2} & \tfrac{1}{2} \\[2pt] \tfrac{1}{2} & -\tfrac{1}{2} \end{pmatrix}\begin{pmatrix} 5 \\ 1 \end{pmatrix}=\begin{pmatrix} \tfrac{5+1}{2} \\[2pt] \tfrac{5-1}{2} \end{pmatrix}=\begin{pmatrix} 3 \\ 2 \end{pmatrix}. \]
Vérification : $f(3\vec{\imath}+2\vec{\jmath})=(3+2)\vec{\imath}+(3-2)\vec{\jmath}=5\vec{\imath}+\vec{\jmath}$. Donc $\vec{v}$ a pour coordonnées $(3\,;2)$.

\item $\det(N)=2\times1-0\times1=2\neq0$ : la matrice $N$ est inversible, donc \textbf{oui}, $g$ est un automorphisme du plan.
\end{enumerate}
\textbf{Barème indicatif (sur 6 pts) :} matrices $M$ et $N$ (1) ; déterminant, conclusion et $M^{-1}$ (1,5) ; matrice $P$ (1) ; double calcul de l'image (1) ; vecteur $\vec{v}$ via $M^{-1}$ (1) ; étude de $g$ (0,5).

\textbf{Points de vigilance :} la matrice de $g\circ f$ est $NM$ et non $MN$ ; pour la question 5, utiliser $f^{-1}$ est plus rapide et plus sûr que résoudre un système, à condition d'avoir justifié que $f$ est bijective ; conclure sur la bijectivité uniquement après avoir calculé le déterminant.
\end{corrige}

\begin{corrige}[Exercice 11 — Problème type Bac : la coopérative agricole]
\begin{enumerate}
\item On applique la relation de récurrence avec $m_1=10$ et $c_1=4$ :
\[ \begin{pmatrix} m_2 \\ c_2 \end{pmatrix}=T\begin{pmatrix} 10 \\ 4 \end{pmatrix}=\begin{pmatrix} 2 & 1 \\ 1 & 2 \end{pmatrix}\begin{pmatrix} 10 \\ 4 \end{pmatrix}=\begin{pmatrix} 20+4 \\ 10+8 \end{pmatrix}=\begin{pmatrix} 24 \\ 18 \end{pmatrix}. \]
Au début du deuxième mois : \textbf{24 tonnes de maïs et 18 tonnes de manioc}.

\item On calcule :
\[ T^2=\begin{pmatrix} 2 & 1 \\ 1 & 2 \end{pmatrix}\begin{pmatrix} 2 & 1 \\ 1 & 2 \end{pmatrix}=\begin{pmatrix} 4+1 & 2+2 \\ 2+2 & 1+4 \end{pmatrix}=\begin{pmatrix} 5 & 4 \\ 4 & 5 \end{pmatrix}. \]
Appliquer deux fois la transformation revient à multiplier par $T^2$ :
\[ \begin{pmatrix} m_3 \\ c_3 \end{pmatrix}=T^2\begin{pmatrix} 10 \\ 4 \end{pmatrix}=\begin{pmatrix} 5 & 4 \\ 4 & 5 \end{pmatrix}\begin{pmatrix} 10 \\ 4 \end{pmatrix}=\begin{pmatrix} 50+16 \\ 40+20 \end{pmatrix}=\begin{pmatrix} 66 \\ 56 \end{pmatrix}. \]
Au début du troisième mois : \textbf{66 tonnes de maïs et 56 tonnes de manioc}. Vérification par récurrence directe : $T\begin{pmatrix} 24 \\ 18 \end{pmatrix}=\begin{pmatrix} 48+18 \\ 24+36 \end{pmatrix}=\begin{pmatrix} 66 \\ 56 \end{pmatrix}$.

\item $\det(T)=2\times2-1\times1=3\neq0$, donc $T$ est inversible et
\[ T^{-1}=\frac{1}{3}\begin{pmatrix} 2 & -1 \\ -1 & 2 \end{pmatrix}. \]

\item Si les quantités en fin de mois sont $\begin{pmatrix} 46 \\ 38 \end{pmatrix}=T\begin{pmatrix} m \\ c \end{pmatrix}$, alors les quantités en début de mois sont
\[ \begin{pmatrix} m \\ c \end{pmatrix}=T^{-1}\begin{pmatrix} 46 \\ 38 \end{pmatrix}=\frac{1}{3}\begin{pmatrix} 2 & -1 \\ -1 & 2 \end{pmatrix}\begin{pmatrix} 46 \\ 38 \end{pmatrix}=\frac{1}{3}\begin{pmatrix} 92-38 \\ -46+76 \end{pmatrix}=\frac{1}{3}\begin{pmatrix} 54 \\ 30 \end{pmatrix}=\begin{pmatrix} 18 \\ 10 \end{pmatrix}. \]
La coopérative disposait au début de ce mois de \textbf{18 tonnes de maïs et 10 tonnes de manioc}. Vérification : $T\begin{pmatrix} 18 \\ 10 \end{pmatrix}=\begin{pmatrix} 36+10 \\ 18+20 \end{pmatrix}=\begin{pmatrix} 46 \\ 38 \end{pmatrix}$.

\item L'agronome a \textbf{tort}. Appliquer deux fois la transformation correspond à la composée, dont la matrice est le \emph{produit} $T^2$, et non le double $2T$. Or
\[ T^2=\begin{pmatrix} 5 & 4 \\ 4 & 5 \end{pmatrix}\qquad\text{alors que}\qquad 2T=\begin{pmatrix} 4 & 2 \\ 2 & 4 \end{pmatrix}. \]
Comme $T^2\neq2T$, l'affirmation est fausse. (Multiplier par 2 reviendrait à doubler chaque coefficient, ce qui ne correspond à aucune composition.)
\end{enumerate}
\textbf{Barème indicatif (sur 5 pts) :} calcul du mois 2 (0,5) ; $T^2$ et mois 3 (1,5) ; déterminant et $T^{-1}$ (1) ; remontée par $T^{-1}$ (1) ; discussion et justification (1).

\textbf{Points de vigilance :} bien distinguer $T^2$ (composer deux fois) et $2T$ (doubler la matrice) ; pour « remonter le temps », c'est $T^{-1}$ qu'il faut appliquer, après avoir justifié l'inversibilité par $\det(T)\neq0$ ; toujours vérifier les quantités obtenues en réappliquant $T$.
\end{corrige}

\newpage

\coursec{Fiche bilan}

\begin{popfigure}
\begin{tikzpicture}[mindmap, grow cyclic, every node/.style={concept, text=white, font=\small\bfseries},
  root concept/.append style={concept color=popdark, font=\bfseries},
  level 1/.append style={level distance=3.6cm, sibling angle=60},
  level 2/.append style={level distance=2.4cm, sibling angle=45, font=\scriptsize}]
\node[root concept]{Matrices $2\times 2$ et applications linéaires}
  child[concept color=popblue]{ node {Définition et opérations}
    child { node[concept color=popblue!70] {Somme, produit par un scalaire} }
    child { node[concept color=popblue!70] {Produit $AB$ (non commutatif)} }
  }
  child[concept color=poporange]{ node {Déterminant}
    child { node[concept color=poporange!70] {$\det A = ad-bc$} }
    child { node[concept color=poporange!70] {Interprétation aire} }
  }
  child[concept color=popgreen]{ node {Matrice inverse}
    child { node[concept color=popgreen!70, align=center] {$A^{-1}=\frac{1}{\det A}\begin{pmatrix} d & -b \\ -c & a \end{pmatrix}$} }
    child { node[concept color=popgreen!70] {Existe ssi $\det A\neq 0$} }
  }
  child[concept color=poppurple]{ node {Systèmes linéaires}
    child { node[concept color=poppurple!70] {$AX=B \Rightarrow X=A^{-1}B$} }
    child { node[concept color=poppurple!70] {Cas $\det A=0$} }
  }
  child[concept color=poppink]{ node {Matrice d'une application linéaire}
    child { node[concept color=poppink!70] {Colonnes = images des vecteurs de base} }
    child { node[concept color=poppink!70] {Coordonnées image : $MX$} }
  }
  child[concept color=popgold]{ node {Composée et réciproque}
    child { node[concept color=popgold!70] {$M_{g\circ f}=M_g M_f$} }
    child { node[concept color=popgold!70] {$M_{f^{-1}}=M_f^{-1}$} }
  };
\end{tikzpicture}
\legende{Carte mentale de la leçon : les six compétences clés autour des matrices $2\times 2$.}
\end{popfigure}

\coursec{Matrices $2\times 2$ : définition et opérations}

\courssub{Définition et notations}
\begin{definition}[Matrice carrée d'ordre 2]
Une \cle{matrice carrée d'ordre 2} est un tableau carré à 2 lignes et 2 colonnes à coefficients dans $\mathbb{R}$ :
\[
A = \begin{pmatrix} a & b \\ c & d \end{pmatrix} = (a_{ij})_{\substack{1\le i\le 2\\1\le j\le 2}}
\]
où $a_{ij}$ désigne le coefficient de la ligne $i$ et de la colonne $j$. L'ensemble de ces matrices est noté $\mathcal{M}_2(\mathbb{R})$.
\end{definition}

\begin{exemplebox}[Exemples de matrices]
\begin{itemize}
  \item Matrice nulle : $O_2 = \begin{pmatrix} 0 & 0 \\ 0 & 0 \end{pmatrix}$.
  \item Matrice identité : $I_2 = \begin{pmatrix} 1 & 0 \\ 0 & 1 \end{pmatrix}$.
  \item Matrice diagonale : $D = \begin{pmatrix} 3 & 0 \\ 0 & -2 \end{pmatrix}$.
  \item Matrice quelconque : $A = \begin{pmatrix} 1 & 4 \\ -2 & 5 \end{pmatrix}$.
\end{itemize}
\end{exemplebox}

\courssub{Addition et multiplication par un scalaire}
\begin{propriete}[Opérations externes]
Soient $A=\begin{pmatrix} a & b \\ c & d \end{pmatrix}$, $B=\begin{pmatrix} a' & b' \\ c' & d' \end{pmatrix}$ et $k\in\mathbb{R}$.
\begin{itemize}
  \item \textbf{Somme} : $A+B = \begin{pmatrix} a+a' & b+b' \\ c+c' & d+d' \end{pmatrix}$.
  \item \textbf{Produit par un scalaire} : $kA = \begin{pmatrix} ka & kb \\ kc & kd \end{pmatrix}$.
\end{itemize}
$(\mathcal{M}_2(\mathbb{R}), +, \times)$ est un \cle{espace vectoriel} de dimension 4 sur $\mathbb{R}$.
\end{propriete}

\courssub{Produit de deux matrices}
\begin{definition}[Produit matriciel]
Le produit $AB$ est défini par :
\[
AB = \begin{pmatrix} a & b \\ c & d \end{pmatrix}\begin{pmatrix} a' & b' \\ c' & d' \end{pmatrix}
= \begin{pmatrix} aa'+bc' & ab'+bd' \\ ca'+dc' & cb'+dd' \end{pmatrix}
\]
Le coefficient de la ligne $i$, colonne $j$ s'obtient par produit scalaire de la ligne $i$ de $A$ et de la colonne $j$ de $B$.
\end{definition}

\begin{attention}[Pièges classiques sur le produit]
\begin{itemize}
  \item \textbf{Non-commutativité} : en général $AB \neq BA$. L'ordre est crucial.
  \item \textbf{Produit nul} : $AB = O_2$ n'implique pas $A=O_2$ ou $B=O_2$ (diviseurs de zéro).
  \item \textbf{Dimension} : le produit n'est défini que si le nombre de colonnes de $A$ égale le nombre de lignes de $B$ (ici toujours vrai pour des $2\times 2$).
\end{itemize}
\end{attention}

\begin{exoresolu}[Calcul de produits]
Soient $A=\begin{pmatrix} 1 & 2 \\ 3 & 4 \end{pmatrix}$ et $B=\begin{pmatrix} 0 & 1 \\ -1 & 2 \end{pmatrix}$.
Calculer $AB$ et $BA$.
\tcblower
$AB = \begin{pmatrix} 1\times0+2\times(-1) & 1\times1+2\times2 \\ 3\times0+4\times(-1) & 3\times1+4\times2 \end{pmatrix} = \begin{pmatrix} -2 & 5 \\ -4 & 11 \end{pmatrix}$.
$BA = \begin{pmatrix} 0\times1+1\times3 & 0\times2+1\times4 \\ -1\times1+2\times3 & -1\times2+2\times4 \end{pmatrix} = \begin{pmatrix} 3 & 4 \\ 5 & 6 \end{pmatrix}$.
On voit bien que $AB \neq BA$.
\end{exoresolu}

\coursec{Déterminant et inverse d'une matrice $2\times 2$}

\courssub{Déterminant}
\begin{definition}[Déterminant]
Le \cle{déterminant} de $A=\begin{pmatrix} a & b \\ c & d \end{pmatrix}$ est le réel :
\[
\det A = |A| = ad - bc.
\]
\end{definition}

\begin{propriete}[Propriétés du déterminant]
\begin{itemize}
  \item $\det(AB) = \det A \times \det B$.
  \item $\det(kA) = k^2 \det A$ (car $k$ multiplie les 2 lignes).
  \item $\det A = 0 \iff$ les lignes (ou colonnes) sont proportionnelles $\iff$ les vecteurs colonnes sont colinéaires.
  \item \textbf{Interprétation géométrique} : $|\det A|$ est l'aire du parallélogramme engendré par les vecteurs colonnes de $A$ dans une base orthonormée.
\end{itemize}
\end{propriete}

\courssub{Matrice inverse}
\begin{propriete}[Inversibilité]
Une matrice $A\in\mathcal{M}_2(\mathbb{R})$ est \cle{inversible} (ou \cle{régulière}) si et seulement si $\det A \neq 0$.
L'inverse $A^{-1}$ est l'unique matrice telle que $AA^{-1}=A^{-1}A=I_2$.
\end{propriete}

\begin{propriete}[Formule de l'inverse]
Si $\det A = ad-bc \neq 0$, alors :
\[
A^{-1} = \frac{1}{ad-bc}\begin{pmatrix} d & -b \\ -c & a \end{pmatrix}.
\]
\emph{Méthode mnémotechnique :} on échange $a$ et $d$, on change les signes de $b$ et $c$, on divise par $\det A$.
\end{propriete}

\begin{exemplebox}[Calcul d'inverse]
$A = \begin{pmatrix} 2 & 1 \\ 3 & 4 \end{pmatrix}$.
$\det A = 2\times4 - 1\times3 = 8-3 = 5 \neq 0$.
$A^{-1} = \frac{1}{5}\begin{pmatrix} 4 & -1 \\ -3 & 2 \end{pmatrix} = \begin{pmatrix} 0,8 & -0,2 \\ -0,6 & 0,4 \end{pmatrix}$.
Vérification : $AA^{-1} = \begin{pmatrix} 2 & 1 \\ 3 & 4 \end{pmatrix}\begin{pmatrix} 0,8 & -0,2 \\ -0,6 & 0,4 \end{pmatrix} = \begin{pmatrix} 1,6-0,6 & -0,4+0,4 \\ 2,4-2,4 & -0,6+1,6 \end{pmatrix} = I_2$.
\end{exemplebox}

\begin{attention}[Erreurs fréquentes]
\begin{itemize}
  \item Oublier de vérifier $\det A \neq 0$ avant d'écrire $A^{-1}$.
  \item Confondre l'inverse de la matrice et l'inverse des coefficients ($\frac{1}{a}$ au lieu de $\frac{d}{\det A}$).
  \item Erreur de signe sur $b$ et $c$ (doivent devenir $-b$ et $-c$).
  \item Écrire $A^{-1} = \frac{1}{\det A}\begin{pmatrix} a & -b \\ -c & d \end{pmatrix}$ (oubli d'échanger $a$ et $d$).
\end{itemize}
\end{attention}

\coursec{Matrices et systèmes linéaires $2\times 2$}

\courssub{Écriture matricielle}
Un système linéaire à deux inconnues :
\[
(S) : \begin{cases} ax + by = e \\ cx + dy = f \end{cases}
\]
s'écrit sous la forme matricielle $AX = B$ avec :
\[
A = \begin{pmatrix} a & b \\ c & d \end{pmatrix},\quad
X = \begin{pmatrix} x \\ y \end{pmatrix},\quad
B = \begin{pmatrix} e \\ f \end{pmatrix}.
\]

\courssub{Résolution par l'inverse}
\begin{propriete}[Résolution de $AX=B$]
Si $\det A \neq 0$, le système admet une \cle{solution unique} donnée par :
\[
X = A^{-1}B.
\]
Si $\det A = 0$, le système n'a pas de solution unique (soit aucune, soit une infinité) ; on ne peut pas utiliser la méthode de l'inverse.
\end{propriete}

\begin{methode}[Résoudre un système $2\times 2$ par matrice inverse]
\begin{enumerate}
  \item Écrire le système sous la forme $AX=B$ (identifier $A, X, B$).
  \item Calculer $\det A = ad-bc$.
  \item Si $\det A \neq 0$, calculer $A^{-1} = \frac{1}{\det A}\begin{pmatrix} d & -b \\ -c & a \end{pmatrix}$.
  \item Calculer $X = A^{-1}B$ (produit matrice $\times$ vecteur colonne).
  \item Vérifier la solution dans le système initial.
\end{enumerate}
\end{methode}

\begin{exoresolu}[Résolution d'un système]
Résoudre $\begin{cases} 2x + y = 5 \\ 3x + 4y = 11 \end{cases}$.
\tcblower
$A=\begin{pmatrix} 2 & 1 \\ 3 & 4 \end{pmatrix}$, $B=\begin{pmatrix} 5 \\ 11 \end{pmatrix}$.
$\det A = 8-3=5 \neq 0$.
$A^{-1} = \frac{1}{5}\begin{pmatrix} 4 & -1 \\ -3 & 2 \end{pmatrix}$.
$X = A^{-1}B = \frac{1}{5}\begin{pmatrix} 4 & -1 \\ -3 & 2 \end{pmatrix}\begin{pmatrix} 5 \\ 11 \end{pmatrix}
= \frac{1}{5}\begin{pmatrix} 20-11 \\ -15+22 \end{pmatrix} = \frac{1}{5}\begin{pmatrix} 9 \\ 7 \end{pmatrix} = \begin{pmatrix} 1,8 \\ 1,4 \end{pmatrix}$.
Solution : $x=1,8$ ; $y=1,4$. Vérification : $2(1,8)+1,4=5$ ; $3(1,8)+4(1,4)=5,4+5,6=11$. OK.
\end{exoresolu}

\begin{savaistu}[Contexte camerounais : Facturation télécom]
Un abonné MTN a deux forfaits : voix et data. En janvier, 2h voix + 1 Go data coûtent 1500 FCFA. En février, 3h voix + 4 Go data coûtent 3400 FCFA. On modélise par $\begin{cases} 2x+y=1500 \\ 3x+4y=3400 \end{cases}$ où $x$ = prix 1h voix, $y$ = prix 1 Go data. La matrice $A=\begin{pmatrix}2&1\\3&4\end{pmatrix}$ a $\det=5$. $A^{-1}=\frac{1}{5}\begin{pmatrix}4&-1\\-3&2\end{pmatrix}$. $X=\frac{1}{5}\begin{pmatrix}4&-1\\-3&2\end{pmatrix}\begin{pmatrix}1500\\3400\end{pmatrix}=\frac{1}{5}\begin{pmatrix}2600\\2300\end{pmatrix}=\begin{pmatrix}520\\460\end{pmatrix}$. 1h voix = 520 FCFA, 1 Go data = 460 FCFA.
\end{savaistu}

\coursec{Matrice d'une application linéaire du plan}

\courssub{Application linéaire : définition}
\begin{definition}[Application linéaire]
Une application $f : \mathbb{R}^2 \to \mathbb{R}^2$ est \cle{linéaire} si pour tous vecteurs $\vec{u},\vec{v}$ et tout réel $\lambda$ :
\[
f(\vec{u}+\vec{v}) = f(\vec{u})+f(\vec{v}) \quad\text{et}\quad f(\lambda\vec{u}) = \lambda f(\vec{u}).
\]
Conséquence immédiate : $f(\vec{0})=\vec{0}$ et $f(\vec{u}-\vec{v})=f(\vec{u})-f(\vec{v})$.
\end{definition}

\begin{exemplebox}[Exemples d'applications linéaires]
\begin{itemize}
  \item Homothétie de rapport $k$ : $h_k(\vec{u}) = k\vec{u}$.
  \item Symétrie axiale, symétrie centrale (rotation de $\pi$).
  \item Rotation d'angle $\alpha$.
  \item Projection orthogonale sur une droite.
  \item \textbf{Contre-exemple} : translation $\vec{u} \mapsto \vec{u}+\vec{a}$ ($\vec{a}\neq\vec{0}$) n'est pas linéaire car $f(\vec{0})=\vec{a}\neq\vec{0}$.
\end{itemize}
\end{exemplebox}

\courssub{Matrice d'une application linéaire dans une base}
Soit $\mathcal{B}=(\vec{\imath},\vec{\jmath})$ une base du plan (souvent la base canonique orthonormée).
Soit $f$ une application linéaire. Les images des vecteurs de base s'écrivent :
\[
f(\vec{\imath}) = a\vec{\imath} + c\vec{\jmath},\qquad f(\vec{\jmath}) = b\vec{\imath} + d\vec{\jmath}.
\]

\begin{definition}[Matrice de $f$ dans la base $\mathcal{B}$]
La \cle{matrice de $f$} dans la base $\mathcal{B}$ est la matrice $M_f = \begin{pmatrix} a & b \\ c & d \end{pmatrix}$ dont les \textbf{colonnes} sont les coordonnées de $f(\vec{\imath})$ et $f(\vec{\jmath})$ dans $\mathcal{B}$.
\[
M_f = \big( \text{coordonnées de } f(\vec{\imath}) \mid \text{coordonnées de } f(\vec{\jmath}) \big).
\]
\end{definition}

\begin{propriete}[Coordonnées de l'image]
Si $\vec{u}$ a pour coordonnées $X = \begin{pmatrix} x \\ y \end{pmatrix}$ dans $\mathcal{B}$, alors $f(\vec{u})$ a pour coordonnées $Y = M_f X$ dans $\mathcal{B}$.
\[
Y = \begin{pmatrix} a & b \\ c & d \end{pmatrix}\begin{pmatrix} x \\ y \end{pmatrix} = \begin{pmatrix} ax+by \\ cx+dy \end{pmatrix}.
\]
\end{propriete}

\begin{methode}[Déterminer la matrice d'une application linéaire]
\begin{enumerate}
  \item Choisir une base $\mathcal{B}=(\vec{\imath},\vec{\jmath})$ (souvent canonique).
  \item Calculer $f(\vec{\imath})$ et $f(\vec{\jmath})$.
  \item Exprimer ces images dans $\mathcal{B}$ : $f(\vec{\imath}) = \begin{pmatrix} a \\ c \end{pmatrix}$, $f(\vec{\jmath}) = \begin{pmatrix} b \\ d \end{pmatrix}$.
  \item Écrire la matrice $M_f = \begin{pmatrix} a & b \\ c & d \end{pmatrix}$ (colonnes = images).
\end{enumerate}
\end{methode}

\begin{exoresolu}[Matrice d'une rotation]
Soit $r$ la rotation de centre $O$ et d'angle $+\frac{\pi}{2}$ (90°) dans le plan muni de la base canonique orthonormée $(\vec{\imath},\vec{\jmath})$.
Déterminer la matrice de $r$.
\tcblower
$r(\vec{\imath}) = \vec{\jmath} = 0\vec{\imath} + 1\vec{\jmath} \Rightarrow$ coordonnées $\begin{pmatrix} 0 \\ 1 \end{pmatrix}$.
$r(\vec{\jmath}) = -\vec{\imath} = -1\vec{\imath} + 0\vec{\jmath} \Rightarrow$ coordonnées $\begin{pmatrix} -1 \\ 0 \end{pmatrix}$.
Matrice : $M_r = \begin{pmatrix} 0 & -1 \\ 1 & 0 \end{pmatrix}$.
Vérification : image de $\vec{u}=\begin{pmatrix} x \\ y \end{pmatrix}$ est $M_r X = \begin{pmatrix} -y \\ x \end{pmatrix}$, bien la rotation de 90°.
\end{exoresolu}

\begin{attention}[Piège : lignes vs colonnes]
Les coordonnées de $f(\vec{\imath})$ forment la \textbf{première colonne}, celles de $f(\vec{\jmath})$ la \textbf{deuxième colonne}. Ne jamais les mettre en lignes.
\end{attention}

\coursec{Opérations sur les applications linéaires et automorphismes}

\courssub{Somme et composée}
\begin{propriete}[Matrice de la somme]
Si $f$ et $g$ sont deux applications linéaires de matrices $M_f$ et $M_f$ dans la même base, alors $f+g$ est linéaire et sa matrice est $M_f + M_g$.
\end{propriete}

\begin{propriete}[Matrice de la composée]
Si $f$ et $g$ sont deux applications linéaires de matrices $M_f$ et $M_g$ dans la même base, alors $g\circ f$ est linéaire et sa matrice est le \textbf{produit} $M_g M_f$ (dans cet ordre : la matrice de la seconde appliquée est à gauche).
\[
M_{g\circ f} = M_g \times M_f.
\]
\end{propriete}

\begin{attention}[Ordre dans la composée]
$g\circ f$ signifie « $f$ puis $g$ ». La matrice est $M_g M_f$ (on applique $M_f$ au vecteur, puis $M_g$ au résultat : $M_g(M_f X) = (M_g M_f)X$). L'ordre des matrices est \textbf{inverse} de l'ordre d'écriture usuel des fonctions.
\end{attention}

\courssub{Automorphismes et application réciproque}
\begin{definition}[Automorphisme]
Une application linéaire $f$ est un \cle{automorphisme} si elle est bijective (injective et surjective).
\end{definition}

\begin{propriete}[Caractérisation par la matrice]
$f$ est un automorphisme $\iff$ sa matrice $M_f$ est inversible $\iff \det M_f \neq 0$.
Dans ce cas, l'application réciproque $f^{-1}$ est linéaire et sa matrice est l'inverse de $M_f$ :
\[
M_{f^{-1}} = (M_f)^{-1}.
\]
\end{propriete}

\begin{exoresolu}[Composée et réciproque]
Soient $f$ l'homothétie de rapport 2 ($M_f = \begin{pmatrix} 2 & 0 \\ 0 & 2 \end{pmatrix}$) et $g$ la symétrie par rapport à l'axe des abscisses ($M_g = \begin{pmatrix} 1 & 0 \\ 0 & -1 \end{pmatrix}$).
1. Matrice de $g\circ f$.
2. $f$ est-il un automorphisme ? Si oui, matrice de $f^{-1}$.
\tcblower
1. $M_{g\circ f} = M_g M_f = \begin{pmatrix} 1 & 0 \\ 0 & -1 \end{pmatrix}\begin{pmatrix} 2 & 0 \\ 0 & 2 \end{pmatrix} = \begin{pmatrix} 2 & 0 \\ 0 & -2 \end{pmatrix}$.
2. $\det M_f = 4 \neq 0$, donc $f$ est un automorphisme. $M_{f^{-1}} = M_f^{-1} = \frac{1}{4}\begin{pmatrix} 2 & 0 \\ 0 & 2 \end{pmatrix} = \begin{pmatrix} 0,5 & 0 \\ 0 & 0,5 \end{pmatrix}$ (homothétie de rapport 1/2).
\end{exoresolu}

\coursec{Applications géométriques et synthèse}

\courssub{Matrices des transformations usuelles (base canonique orthonormée)}
\begin{center}
\begin{tabularx}{\linewidth}{|l|X|c|}
\hline
\rowcolor{popblueL} \textbf{Transformation} & \textbf{Matrice} & \textbf{Déterminant} \\
\hline
Homothétie rapport $k$ & $\begin{pmatrix} k & 0 \\ 0 & k \end{pmatrix}$ & $k^2$ \\
\hline
Symétrie axiale (axe $Ox$) & $\begin{pmatrix} 1 & 0 \\ 0 & -1 \end{pmatrix}$ & $-1$ \\
\hline
Symétrie centrale (rotation $\pi$) & $\begin{pmatrix} -1 & 0 \\ 0 & -1 \end{pmatrix}$ & $1$ \\
\hline
Rotation angle $\alpha$ & $\begin{pmatrix} \cos\alpha & -\sin\alpha \\ \sin\alpha & \cos\alpha \end{pmatrix}$ & $1$ \\
\hline
Projection orthogonale sur $Ox$ & $\begin{pmatrix} 1 & 0 \\ 0 & 0 \end{pmatrix}$ & $0$ \\
\hline
Cisaillement horizontal ($x\mapsto x+ky$) & $\begin{pmatrix} 1 & k \\ 0 & 1 \end{pmatrix}$ & $1$ \\
\hline
\end{tabularx}
\end{center}

\begin{propriete}[Interprétation du déterminant pour une isométrie]
Pour une application linéaire $f$ de matrice $M$ dans une base orthonormée :
\begin{itemize}
  \item $|\det M| = 1$ : $f$ conserve les aires (isométrie ou rotation).
  \item $\det M = 1$ : $f$ conserve l'orientation (rotation).
  \item $\det M = -1$ : $f$ inverse l'orientation (symétrie axiale).
  \item $\det M = 0$ : $f$ écrase le plan sur une droite (projection) ou un point (nulle), aire nulle.
\end{itemize}
\end{propriete}

\begin{exemplebox}[Composition de transformations géométriques]
Une rotation $r_{\pi/2}$ suivie d'une symétrie axiale $s_x$ (axe $Ox$).
$M_r = \begin{pmatrix} 0 & -1 \\ 1 & 0 \end{pmatrix}$, $M_s = \begin{pmatrix} 1 & 0 \\ 0 & -1 \end{pmatrix}$.
Matrice de $s_x \circ r_{\pi/2}$ : $M_s M_r = \begin{pmatrix} 1 & 0 \\ 0 & -1 \end{pmatrix}\begin{pmatrix} 0 & -1 \\ 1 & 0 \end{pmatrix} = \begin{pmatrix} 0 & -1 \\ -1 & 0 \end{pmatrix}$.
C'est la matrice de la symétrie par rapport à la droite $y=-x$ (bisectrice des 2e et 4e quadrants). $\det = -1$, orientation inversée.
\end{exemplebox}

\begin{savaistu}[Application : Graphismes 2D et Jeux vidéo]
Dans les moteurs graphiques (Unity, Godot) ou le traitement d'images (OpenCV), chaque pixel/vertex est un vecteur $(x,y)$. Appliquer une transformation (rotation, zoom, translation via matrices homogènes $3\times3$) revient à multiplier des millions de fois des matrices $2\times2$ (ou $3\times3$) par des vecteurs. La rapidité du calcul matriciel (optimisé GPU) permet l'affichage temps réel à 60 images/seconde. Au Cameroun, les startups de \emph{gamedev} (ex: Kiro'o Games) utilisent ces mathématiques quotidiennement.
\end{savaistu}

\courssub{Exercices types Bac}

\begin{exercice}[Exercice 1 : Matrice inverse et système]
Soit $A = \begin{pmatrix} 3 & -1 \\ 2 & 4 \end{pmatrix}$.
\begin{enumerate}
  \item Calculer $\det A$. En déduire si $A$ est inversible.
  \item Si oui, calculer $A^{-1}$.
  \item Résoudre le système $\begin{cases} 3x - y = 7 \\ 2x + 4y = 10 \end{cases}$ par la méthode de l'inverse.
\end{enumerate}
\end{exercice}

\begin{corrige}[Exercice 1]
\begin{enumerate}
  \item $\det A = 3\times4 - (-1)\times2 = 12+2=14 \neq 0$. $A$ est inversible.
  \item $A^{-1} = \frac{1}{14}\begin{pmatrix} 4 & 1 \\ -2 & 3 \end{pmatrix} = \begin{pmatrix} 2/7 & 1/14 \\ -1/7 & 3/14 \end{pmatrix}$.
  \item $X = A^{-1}B = \frac{1}{14}\begin{pmatrix} 4 & 1 \\ -2 & 3 \end{pmatrix}\begin{pmatrix} 7 \\ 10 \end{pmatrix} = \frac{1}{14}\begin{pmatrix} 28+10 \\ -14+30 \end{pmatrix} = \frac{1}{14}\begin{pmatrix} 38 \\ 16 \end{pmatrix} = \begin{pmatrix} 19/7 \\ 8/7 \end{pmatrix}$.
  Solution : $x=\frac{19}{7}$, $y=\frac{8}{7}$.
\end{enumerate}
\end{corrige}

\begin{exercice}[Exercice 2 : Matrice d'une application linéaire]
Dans le plan muni d'un repère orthonormé $(O;\vec{\imath},\vec{\jmath})$, on considère l'application $f$ qui à tout point $M(x,y)$ associe $M'(x',y')$ tel que :
\[
\begin{cases} x' = 2x + y \\ y' = x - 3y \end{cases}
\]
\begin{enumerate}
  \item Montrer que $f$ est une application linéaire.
  \item Déterminer la matrice $M_f$ de $f$ dans la base $(\vec{\imath},\vec{\jmath})$.
  \item Calculer $\det M_f$. $f$ est-elle un automorphisme ?
  \item Déterminer les coordonnées de l'image du vecteur $\vec{u}=\begin{pmatrix} 1 \\ -2 \end{pmatrix}$.
  \item Si $f$ est un automorphisme, donner la matrice de $f^{-1}$.
\end{enumerate}
\end{exercice}

\begin{corrige}[Exercice 2]
\begin{enumerate}
  \item $f(\vec{u}+\vec{v}) = \begin{pmatrix} 2(x_1+x_2)+(y_1+y_2) \\ (x_1+x_2)-3(y_1+y_2) \end{pmatrix} = f(\vec{u})+f(\vec{v})$. $f(\lambda\vec{u}) = \lambda f(\vec{u})$. Donc $f$ est linéaire.
  \item $f(\vec{\imath}) = f(1,0) = (2,1)^T$, $f(\vec{\jmath}) = f(0,1) = (1,-3)^T$.
  $M_f = \begin{pmatrix} 2 & 1 \\ 1 & -3 \end{pmatrix}$.
  \item $\det M_f = 2\times(-3) - 1\times1 = -7 \neq 0$. $f$ est un automorphisme.
  \item $M_f \begin{pmatrix} 1 \\ -2 \end{pmatrix} = \begin{pmatrix} 2-2 \\ 1+6 \end{pmatrix} = \begin{pmatrix} 0 \\ 7 \end{pmatrix}$.
  \item $M_f^{-1} = -\frac{1}{7}\begin{pmatrix} -3 & -1 \\ -1 & 2 \end{pmatrix} = \begin{pmatrix} 3/7 & 1/7 \\ 1/7 & -2/7 \end{pmatrix}$.
\end{enumerate}
\end{corrige}

\begin{exercice}[Exercice 3 : Composition de transformations]
Soit $r$ la rotation d'angle $+\frac{\pi}{3}$ et $h$ l'homothétie de rapport $2$, toutes deux de centre $O$.
\begin{enumerate}
  \item Écrire les matrices $M_r$ et $M_h$.
  \item Déterminer la matrice de $h\circ r$ et celle de $r\circ h$. Sont-elles égales ?
  \item Quelle transformation géométrique représente $h\circ r$ ?
\end{enumerate}
\end{exercice}

\begin{corrige}[Exercice 3]
\begin{enumerate}
  \item $M_r = \begin{pmatrix} \cos\frac{\pi}{3} & -\sin\frac{\pi}{3} \\ \sin\frac{\pi}{3} & \cos\frac{\pi}{3} \end{pmatrix} = \begin{pmatrix} 1/2 & -\sqrt{3}/2 \\ \sqrt{3}/2 & 1/2 \end{pmatrix}$.
  $M_h = \begin{pmatrix} 2 & 0 \\ 0 & 2 \end{pmatrix} = 2I_2$.
  \item $M_{h\circ r} = M_h M_r = 2I_2 M_r = 2M_r = \begin{pmatrix} 1 & -\sqrt{3} \\ \sqrt{3} & 1 \end{pmatrix}$.
  $M_{r\circ h} = M_r M_h = M_r (2I_2) = 2M_r = \begin{pmatrix} 1 & -\sqrt{3} \\ \sqrt{3} & 1 \end{pmatrix}$.
  Elles sont égales car $M_h = 2I_2$ commute avec toute matrice.
  \item $h\circ r$ est une rotation d'angle $\frac{\pi}{3}$ composée avec une homothétie de rapport 2 : c'est une \textbf{similitude directe} de rapport 2 et d'angle $\frac{\pi}{3}$ (ou rotation-homothétie).
\end{enumerate}
\end{corrige}

\begin{aretenir}[Checklist finale avant le Bac — Matrices $2\times 2$]
\begin{itemize}
  \item[$\square$] Je définis une matrice $2\times 2$, je fais la somme, le produit par un scalaire, le produit $AB$ (ligne $\times$ colonne).
  \item[$\square$] Je sais que $AB \neq BA$ en général ; je donne un contre-exemple simple.
  \item[$\square$] Je calcule $\det A = ad-bc$ sans erreur de signe ; j'interprète $\det A=0$ (vecteurs colinéaires, aire nulle).
  \item[$\square$] Je vérifie $\det A \neq 0$ \textbf{avant} d'écrire $A^{-1}$ ; j'applique la formule $\frac{1}{\det A}\begin{pmatrix} d & -b \\ -c & a \end{pmatrix}$.
  \item[$\square$] Je traduis un système $2\times 2$ en $AX=B$ et je le résous par $X=A^{-1}B$ (si $\det A\neq 0$).
  \item[$\square$] Je reconnais une application linéaire ($f(\vec{0})=\vec{0}$, conservation combinaisons linéaires).
  \item[$\square$] Je construis la matrice de $f$ : \textbf{colonnes = images des vecteurs de base}.
  \item[$\square$] Je calcule l'image d'un vecteur par $Y = M_f X$.
  \item[$\square$] Matrice de $g\circ f$ : $M_g M_f$ (ordre inverse !). Matrice de $f+g$ : $M_f+M_g$.
  \item[$\square$] Automorphisme $\iff \det M_f \neq 0$ ; alors $M_{f^{-1}} = M_f^{-1}$.
  \item[$\square$] Je connais les matrices des transformations usuelles (rotation, symétrie, homothétie, projection) et leurs déterminants.
\end{itemize}
\end{aretenir}

\findecours

\end{document}
